% Using formatting, editing features of Word processors for producing documents — ICT Upper Sixth — Cours signé OnBuch+
% Official MINESEC syllabus. Compiler avec : tectonic ICT21-using-formatting-editing-features-of-word-processors-for-pro.tex
\newcommand{\DOCMATIERE}{ICT}
\newcommand{\DOCNIVEAU}{Upper Sixth}
\newcommand{\DOCMODULE}{Upper Sixth — ICT}
\newcommand{\DOCLECON}{21}
\newcommand{\DOCTITRE}{Using formatting, editing features of Word processors for producing documents}
% OnBuch+ — préambule « cours signés » ENRICHI (compilé avec tectonic / XeLaTeX)
% Système visuel « Pop » d'OnBuch+ : fond crème (jamais blanc), bordures encre
% épaisses, ombres « dures » décalées, titres Archivo Black en capitales.
% Version MATHÉMATIQUES : graphes (pgfplots/tikz), boîtes pédagogiques
% (définition, propriété, méthode, attention, le savais-tu, exercice résolu,
% exercice, TP, bilan), page de garde et signature OnBuch+ ancrée en bas.
%
% Macros attendues dans main.tex AVANT \input{preamble} :
%   \DOCMATIERE  \DOCNIVEAU  \DOCMODULE  \DOCLECON (numéro)  \DOCTITRE
\documentclass[11pt]{article}

\usepackage{fontspec}
\usepackage[english]{babel}
\usepackage[a4paper,margin=2cm,top=3.4cm,bottom=2.8cm,headheight=1.4cm,footskip=1.3cm]{geometry}
\usepackage{amsmath,amssymb,amsfonts}
\usepackage{mathtools} % \xmapsto, \coloneqq... souvent utilisés par les modèles
\usepackage{cancel} % simplifications barrées (bilans de forces)
% \abs{z} (module, valeur absolue) et \norm{u} (norme), utilisés par les modèles
\providecommand{\abs}{}\let\abs\relax\DeclarePairedDelimiter{\abs}{\lvert}{\rvert}
\providecommand{\norm}{}\let\norm\relax\DeclarePairedDelimiter{\norm}{\lVert}{\rVert}
\usepackage[table]{xcolor} % [table] : fournit aussi \rowcolors (alternance de lignes)
\usepackage{pagecolor}
\usepackage{tikz}
\usetikzlibrary{shapes.symbols,circuits.logic.US,circuits.logic.IEC,arrows.meta,positioning,calc,shapes.geometric,shapes.misc,shapes.arrows,decorations.pathmorphing,decorations.pathreplacing,decorations.markings,patterns,shadows,mindmap,trees,fit,backgrounds,matrix,intersections,angles,quotes}
\usepackage{pgfplots}
\usepackage[version=4]{mhchem} % équations nucléaires \ce{^{235}_{92}U}
\usepackage[european,siunitx]{circuitikz} % schémas électriques
\pgfplotsset{compat=1.18}
\usepgfplotslibrary{fillbetween,groupplots,polar,statistics,dateplot} % aires entre courbes (calcul intégral)
\usepackage{enumitem}
\usepackage{fancyhdr}
% [most] charge toutes les bibliothèques tcolorbox (listings, minted, raster,
% documentation...) et gonfle inutilement la mémoire TeX sur un document long
% (~300 boîtes) : on ne charge que ce qui est réellement utilisé.
\usepackage[skins,breakable]{tcolorbox}
\usepackage{graphicx}
\usepackage{colortbl}
\usepackage{booktabs}
\usepackage{tabularx}
\usepackage{diagbox} % case d'en-tête diagonale des tableaux à double entrée
% Colonnes extensibles centrée / gauche / droite (C, L, R), souvent utilisées par les modèles
\newcolumntype{C}{>{\centering\arraybackslash}X}
\newcolumntype{L}{>{\raggedright\arraybackslash}X}
\newcolumntype{R}{>{\raggedleft\arraybackslash}X}
\usepackage{array}
\usepackage{multicol}
\usepackage{siunitx}
% parse-numbers=false : accepte aussi \SI{2,5 \times 10^{-3}}{...}
\sisetup{locale=UK,output-decimal-marker={.},group-minimum-digits=4,parse-numbers=false}
\DeclareSIUnit\ohm{\ensuremath{\symup{\Omega}}} % U+2126 absent de la police
\DeclareSIUnit\division{div} % réglages d'oscilloscope (ms/div, V/div)
\DeclareSIUnit\year{yr}\DeclareSIUnit\annum{yr}\DeclareSIUnit\Ma{Ma}\DeclareSIUnit\tr{tr} % tours (vitesse de rotation, ex. \SI{1500}{\tr\per\minute})
\usepackage{qrcode}
% \degree : commande fréquemment utilisée par les modèles pour un angle ou une
% température en dehors de \SI{}{} (non fournie par les paquets chargés).
\providecommand{\degree}{^{\circ}}
% \qed (fin de démonstration), utilisé par les modèles sans amsthm
\providecommand{\qed}{\hfill$\square$}
% \barre : double barre des valeurs interdites dans un tableau de variations (tkz-tab)
\providecommand{\barre}{\ensuremath{\|}}
% environnement proof (amsthm non chargé) utilisé par les modèles
\ifdefined\proof\else\newenvironment{proof}[1][Proof]{\par\noindent\textit{#1.}\ }{\qed\par}\fi
\usepackage{lastpage}
\usepackage{hyperref}
\hypersetup{hidelinks}

% ============================================================
% Palette « Pop » OnBuch+ — mêmes hex que la classe P (plus_ui.dart)
% ============================================================
\definecolor{poporange}{HTML}{F59321}
\definecolor{popdark}{HTML}{E07A0C}
\definecolor{popcream}{HTML}{FFF6E8}
\definecolor{popink}{HTML}{1C1714}
\definecolor{poppurple}{HTML}{7A5AE0}
\definecolor{popgreen}{HTML}{2E9E6A}
\definecolor{poppink}{HTML}{E0457B}
\definecolor{popblue}{HTML}{2F7DE1}
\definecolor{popgold}{HTML}{C98A12}
\definecolor{popmuted}{HTML}{8A7F76}
\definecolor{popline}{HTML}{EADFD2}
\definecolor{popgray}{HTML}{8A7F76}
\colorlet{popgrey}{popgray}
% Alias tolérants (le modèle nomme parfois les couleurs différemment)
\colorlet{popwhite}{white}
\colorlet{popblack}{popink}
\colorlet{popred}{poppink}
\colorlet{popyellow}{popgold}
% Teintes claires pour fonds de tableaux / zones
\colorlet{popblueL}{popblue!10!white}
\colorlet{poporangeL}{poporange!14!white}
\colorlet{popgreenL}{popgreen!12!white}
\colorlet{poppurpleL}{poppurple!12!white}
\colorlet{poppinkL}{poppink!12!white}
\colorlet{popgoldL}{popgold!14!white}

\pagecolor{popcream}

% ============================================================
% Typographie
% ============================================================
\defaultfontfeatures{Ligatures=TeX}
\setmainfont{PlusJakartaSans-Medium.ttf}[Path=../../../fonts/,BoldFont=PlusJakartaSans-Bold.ttf]
% Maths en sans-serif (Fira Math) avec chiffres et lettres droites de Plus
% Jakarta, pour que les formules se fondent dans le texte.
\usepackage[math-style=ISO,bold-style=ISO]{unicode-math}
\setmathfont{FiraMath-Regular.otf}[Path=../../../fonts/]
\setmathfont{PlusJakartaSans-Medium.ttf}[Path=../../../fonts/,range={up/num,up/{Latin,latin},\mathup}]
% (icomma retiré : décimales avec point en anglais)
% Caractères mathématiques Unicode absents des polices de texte (Plus Jakarta,
% Archivo Black) : rendus par leur équivalent mathématique, en texte comme en
% formule — sinon ils s'affichent en carré vide (ex. « Arithmétique dans ℤ »).
\usepackage{newunicodechar}
\newunicodechar{ℝ}{\ensuremath{\mathbb{R}}}
\newunicodechar{ℤ}{\ensuremath{\mathbb{Z}}}
\newunicodechar{ℂ}{\ensuremath{\mathbb{C}}}
\newunicodechar{ℕ}{\ensuremath{\mathbb{N}}}
\newunicodechar{ℚ}{\ensuremath{\mathbb{Q}}}
\newunicodechar{𝔻}{\ensuremath{\mathbb{D}}}
\newunicodechar{α}{\ensuremath{\alpha}}
\newunicodechar{β}{\ensuremath{\beta}}
\newunicodechar{γ}{\ensuremath{\gamma}}
\newunicodechar{δ}{\ensuremath{\delta}}
\newunicodechar{ε}{\ensuremath{\varepsilon}}
\newunicodechar{θ}{\ensuremath{\theta}}
\newunicodechar{λ}{\ensuremath{\lambda}}
\newunicodechar{μ}{\ensuremath{\mu}}
\newunicodechar{σ}{\ensuremath{\sigma}}
\newunicodechar{τ}{\ensuremath{\tau}}
\newunicodechar{φ}{\ensuremath{\varphi}}
\newunicodechar{ψ}{\ensuremath{\psi}}
\newunicodechar{ω}{\ensuremath{\omega}}
\newunicodechar{Σ}{\ensuremath{\Sigma}}
% majuscules produites par \MakeUppercase dans les titres (α-aminés -> Α)
\newunicodechar{Α}{\ensuremath{\alpha}}
\newunicodechar{Β}{\ensuremath{\beta}}
\newunicodechar{Γ}{\ensuremath{\Gamma}}
\newunicodechar{Δ}{\ensuremath{\Delta}}
\newunicodechar{Ε}{\ensuremath{\varepsilon}}
\newunicodechar{Θ}{\ensuremath{\Theta}}
\newunicodechar{Λ}{\ensuremath{\Lambda}}
\newunicodechar{Μ}{\ensuremath{\mu}}
\newunicodechar{Τ}{\ensuremath{\tau}}
\newunicodechar{Φ}{\ensuremath{\Phi}}
\newunicodechar{Ψ}{\ensuremath{\Psi}}
\newunicodechar{Ω}{\ensuremath{\Omega}}
\newunicodechar{↦}{\ensuremath{\mapsto}}
\newunicodechar{∈}{\ensuremath{\in}}
\newunicodechar{∉}{\ensuremath{\notin}}
\newunicodechar{∧}{\ensuremath{\wedge}}
\newunicodechar{⇔}{\ensuremath{\Leftrightarrow}}
\newunicodechar{⟺}{\ensuremath{\Longleftrightarrow}}
\newunicodechar{⇒}{\ensuremath{\Rightarrow}}
\newunicodechar{⟹}{\ensuremath{\Longrightarrow}}
\newunicodechar{∘}{\ensuremath{\circ}}
\newunicodechar{∪}{\ensuremath{\cup}}
\newunicodechar{∩}{\ensuremath{\cap}}
\newunicodechar{≡}{\ensuremath{\equiv}}
\newunicodechar{⊂}{\ensuremath{\subset}}
\newunicodechar{⊥}{\ensuremath{\perp}}
\newunicodechar{∃}{\ensuremath{\exists}}
\newunicodechar{∀}{\ensuremath{\forall}}
\newunicodechar{∛}{\ensuremath{\sqrt[3]{\,}}}
\newunicodechar{∜}{\ensuremath{\sqrt[4]{\,}}}
\newunicodechar{⃗}{\ensuremath{{}^{\to}}}
\newunicodechar{ⁿ}{\ifmmode^{n}\else\textsuperscript{\ensuremath{n}}\fi}
\newunicodechar{ˣ}{\ifmmode^{x}\else\textsuperscript{\ensuremath{x}}\fi}
\newunicodechar{ᵇ}{\ifmmode^{b}\else\textsuperscript{\ensuremath{b}}\fi}
\newunicodechar{ʳ}{\ifmmode^{r}\else\textsuperscript{\ensuremath{r}}\fi}
\newunicodechar{ᵘ}{\ifmmode^{u}\else\textsuperscript{\ensuremath{u}}\fi}
\newunicodechar{ʸ}{\ifmmode^{y}\else\textsuperscript{\ensuremath{y}}\fi}
\newunicodechar{ᵃ}{\ifmmode^{a}\else\textsuperscript{\ensuremath{a}}\fi}
\newunicodechar{ᵉ}{\ifmmode^{e}\else\textsuperscript{\ensuremath{e}}\fi}
\newunicodechar{ᵏ}{\ifmmode^{k}\else\textsuperscript{\ensuremath{k}}\fi}
\newunicodechar{ᵖ}{\ifmmode^{p}\else\textsuperscript{\ensuremath{p}}\fi}
\newunicodechar{ᴺ}{\ifmmode^{N}\else\textsuperscript{\ensuremath{N}}\fi}
\newunicodechar{ᶜ}{\ifmmode^{c}\else\textsuperscript{\ensuremath{c}}\fi}
\newunicodechar{ʰ}{\ifmmode^{h}\else\textsuperscript{\ensuremath{h}}\fi}
\newunicodechar{ᵠ}{\ifmmode^{q}\else\textsuperscript{\ensuremath{q}}\fi}
\newunicodechar{ₙ}{\ifmmode_{n}\else\textsubscript{\ensuremath{n}}\fi}
\newunicodechar{ₐ}{\ifmmode_{a}\else\textsubscript{\ensuremath{a}}\fi}
\newunicodechar{ᵢ}{\ifmmode_{i}\else\textsubscript{\ensuremath{i}}\fi}
\newunicodechar{ᵣ}{\ifmmode_{r}\else\textsubscript{\ensuremath{r}}\fi}
\newunicodechar{ₖ}{\ifmmode_{k}\else\textsubscript{\ensuremath{k}}\fi}
\newunicodechar{ᵦ}{\ifmmode_{\beta}\else\textsubscript{\ensuremath{\beta}}\fi}
\newunicodechar{ₓ}{\ifmmode_{x}\else\textsubscript{\ensuremath{x}}\fi}
\newfontfamily\popdisplay{ArchivoBlack-Regular.ttf}[Path=../../../fonts/]
\newfontfamily\poplogo{SpaceGrotesk-Bold.ttf}[Path=../../../fonts/]
\newfontfamily\popsemi{PlusJakartaSans-SemiBold.ttf}[Path=../../../fonts/]
\newfontfamily\popxbold{PlusJakartaSans-ExtraBold.ttf}[Path=../../../fonts/]
\newfontfamily\popxboldls{PlusJakartaSans-ExtraBold.ttf}[Path=../../../fonts/,LetterSpace=3.0]

\setlist[enumerate]{leftmargin=1.7em,itemsep=5pt,topsep=4pt}
\setlist[itemize]{leftmargin=1.7em,itemsep=4pt,topsep=4pt,label={\color{poporange}\textbullet}}
\setlength{\parindent}{0pt}
\setlength{\parskip}{5pt}
\renewcommand{\arraystretch}{1.25}

% pgfplots : style Pop par défaut
\pgfplotsset{
  every axis/.append style={axis line style={popink,line width=1pt},tick style={popink},
    grid style={popline,line width=0.5pt},label style={font=\small},tick label style={font=\footnotesize}},
  popcurve/.style={poporange,line width=1.8pt,no markers,smooth},
}
% Même style disponible dans un \draw TikZ simple (hors pgfplots)
\tikzset{popcurve/.style={poporange,line width=1.8pt}}
\tikzset{popgrid/.style={popline,very thin}}
\colorlet{popcurve}{poporange} % le nom du style est parfois utilisé comme couleur

% ============================================================
% Pill / squiggle
% ============================================================
\newcommand{\poppill}[3]{%
  \tikz[baseline=-0.55ex]\node[draw=popink,line width=1.2pt,rounded corners=2.6pt,%
    fill=#2,inner xsep=6.5pt,inner ysep=3pt]{%
    \popxboldls\fontsize{7}{7}\selectfont\color{#3}\MakeUppercase{#1}};%
}
\newcommand{\popsquiggle}[1]{%
  \tikz[baseline=-0.4ex]{
    \draw[#1,line width=2.6pt,line cap=round]
      plot [smooth] coordinates {(0,0.09) (0.34,0.01) (0.68,0.21) (1.02,0.01) (1.36,0.10)};
  }%
}
% Mot-clé surligné (marqueur orange sous le texte)
\newcommand{\cle}[1]{{\popxbold\color{popdark}#1}}

% ============================================================
% En-tête / pied
% ============================================================
\pagestyle{fancy}
\fancyhf{}
\renewcommand{\headrulewidth}{0pt}
\renewcommand{\footrulewidth}{0pt}
\lhead{\raisebox{-0.15ex}{\poplogo\fontsize{13}{13}\selectfont\color{popink}OnBuch\color{poporange}+}}
\rhead{\poppill{\DOCMATIERE\ \textbullet\ \DOCNIVEAU\ \textbullet\ Chapter \DOCLECON}{white}{popink}}
\lfoot{\popxbold\fontsize{7.6}{7.6}\selectfont\color{popmuted}\MakeUppercase{Signed course OnBuch+}\ \textbullet\ {\popsemi\DOCTITRE}}
\rfoot{\tikz[baseline=-0.6ex]\node[rounded corners=8.5pt,fill=popink,minimum height=17pt,minimum width=17pt,inner xsep=5pt,inner sep=0]{\popxbold\fontsize{8}{8}\selectfont\color{popcream}\thepage\,/\,\pageref*{LastPage}};}

% ============================================================
% Titres
% ============================================================
\newcommand{\chaptitre}[2]{%
  \begin{tcolorbox}[enhanced,colback=poporange,colframe=popink,boxrule=2.6pt,arc=11pt,
      drop shadow=popink,left=15pt,right=15pt,top=12pt,bottom=11pt,before skip=3pt,after skip=18pt]
    \poppill{#2}{popink}{popcream}\par\vspace{9pt}
    {\raggedright\hyphenpenalty=10000\exhyphenpenalty=10000\popdisplay\fontsize{22}{24}\selectfont\color{popcream}\MakeUppercase{#1}\par}
  \end{tcolorbox}
}
% Section numérotée : pastille orange avec le numéro + titre Archivo
\newcounter{popsec}
\newcommand{\coursec}[1]{%
  \stepcounter{popsec}\setcounter{popsub}{0}%
  \par\vspace{16pt}\noindent
  \begin{minipage}{\linewidth}
  \tikz[baseline=-0.7ex]\node[draw=popink,line width=1.6pt,fill=poporange,rounded corners=4pt,minimum size=22pt,inner sep=0]{\popdisplay\fontsize{12}{12}\selectfont\color{popcream}\thepopsec};%
  \hspace{8pt}{\popdisplay\fontsize{15}{17}\selectfont\color{popink}\MakeUppercase{#1}}\par
  \vspace{4pt}\noindent{\color{poporange}\rule{36pt}{3.6pt}}
  \end{minipage}\par\nopagebreak\vspace{8pt}
}
\newcounter{popsub}
\newcommand{\courssub}[1]{%
  \stepcounter{popsub}%
  \par\vspace{9pt}\noindent{\popxbold\fontsize{11.5}{13}\selectfont\color{popink}{\color{poporange}\thepopsec.\thepopsub}\hspace{6pt}#1}\par\nopagebreak\vspace{4pt}
}

% ============================================================
% Boîtes pédagogiques — même recette : fond blanc, bordure épaisse colorée,
% ombre dure encre, badge plein. Titre optionnel [#1].
% ============================================================
\tcbset{popbase/.style={breakable,enhanced,colback=white,boxrule=2.3pt,arc=9pt,
  shadow={3pt}{-3pt}{0pt}{popink},left=14pt,right=14pt,top=11pt,bottom=11pt,before skip=11pt,after skip=15pt,
  fontupper=\popsemi\fontsize{10.6}{14.5}\selectfont\color{popink},
  fontlower=\popsemi\fontsize{10.6}{14.5}\selectfont\color{popink}}}
% Coche dessinée (le glyphe ✓ n'existe pas dans les polices du document)
\newcommand{\popcheck}[1]{\tikz[baseline=-0.5ex]\draw[#1,line width=1.6pt,line cap=round,line join=round](-0.13,0.0)--(-0.04,-0.09)--(0.13,0.1);}
\providecommand{\coche}{\popcheck{popgreen}}
\providecommand{\resultat}[1]{\par\smallskip\noindent\fcolorbox{popgreen}{popgreenL}{\parbox{\dimexpr\linewidth-2\fboxsep-2\fboxrule\relax}{\textbf{Result.} #1}}\par\smallskip}
\newcommand{\popboxhead}[3]{\poppill{#1}{#2}{white}\ifx\relax#3\relax\else\hspace{6pt}{\popxbold\fontsize{10.6}{12}\selectfont\color{popink}#3}\fi\par\vspace{7pt}}

\newtcolorbox{aretenir}[1][]{popbase,colframe=popgreen,before upper={\popboxhead{Key points}{popgreen}{#1}}}
\newtcolorbox{exemplebox}[1][]{popbase,colframe=poporange,before upper={\popboxhead{Exemple}{poporange}{#1}}}
\newtcolorbox{definition}[1][]{popbase,colframe=popblue,before upper={\popboxhead{Definition}{popblue}{#1}}}
\newtcolorbox{propriete}[1][]{popbase,colframe=poppurple,before upper={\popboxhead{Property}{poppurple}{#1}}}
\newtcolorbox{methode}[1][]{popbase,colframe=popgold,before upper={\popboxhead{Method}{popgold}{#1}}}
\newtcolorbox{attention}[1][]{popbase,colframe=poppink,before upper={\popboxhead{Warning}{poppink}{#1}}}
\newtcolorbox{experience}[1][]{popbase,colframe=popblue,colback=popblueL,before upper={\popboxhead{Experiment}{popblue}{#1}}}
% « Le savais-tu ? » : carte violette pleine (contexte, culture, Cameroun)
\newtcolorbox{savaistu}[1][]{popbase,colback=poppurple,colframe=popink,
  fontupper=\popsemi\fontsize{10.4}{14.2}\selectfont\color{white},
  before upper={\poppill{Did you know?}{white}{poppurple}\ifx\relax#1\relax\else\hspace{6pt}{\popxbold\color{white}#1}\fi\par\vspace{7pt}}}
% Exercice résolu : énoncé en haut, \tcblower puis la solution sur fond crème
\newcounter{popexo}\newcounter{popexr}
\newtcolorbox{exoresolu}[1][]{popbase,colframe=poporange,colbacklower=popcream,
  segmentation style={popink,line width=1.2pt,dashed},
  before upper={\stepcounter{popexr}\popboxhead{Worked example \thepopexr}{poporange}{#1}},
  before lower={\poppill{Solution}{popink}{popcream}\par\vspace{6pt}}}
% Exercice (non résolu en place) — la correction est dans la partie Corrigés
\newtcolorbox{exercice}[1][]{popbase,colframe=popink,
  before upper={\stepcounter{popexo}\popboxhead{Exercise \thepopexo}{popink}{#1}}}
% Corrigé
\newtcolorbox{corrige}[1][]{popbase,colframe=popgreen,colback=popgreenL,
  before upper={\popboxhead{Solution}{popgreen}{#1}}}
% Objectifs / prérequis / bilan
\newtcolorbox{objectifs}{popbase,colframe=popink,colback=poporangeL,
  before upper={\popboxhead{Chapter objectives}{popink}{}}}
\newtcolorbox{prerequis}{popbase,colframe=popink,colback=popblueL,
  before upper={\popboxhead{Prerequisites}{popblue}{}}}
\newtcolorbox{bilan}{popbase,colframe=popink,colback=white,boxrule=2.8pt,
  before upper={\popboxhead{Fiche bilan}{poporange}{L'essentiel à connaître pour l'examen}}}
% Figure encadrée avec légende
\newtcolorbox{popfigure}[1][]{enhanced,breakable=false,colback=white,colframe=popline,boxrule=1.4pt,arc=8pt,
  left=10pt,right=10pt,top=10pt,bottom=8pt,before skip=10pt,after skip=12pt,halign=center,
  lower separated=false,halign lower=center,
  fontlower=\popsemi\fontsize{9}{11.5}\selectfont\color{popmuted},
  #1}
\newcommand{\legende}[1]{\par\vspace{6pt}{\popsemi\fontsize{9}{11.5}\selectfont\color{popmuted}#1}}

% ============================================================
% Signature OnBuch+
% ============================================================
\newcommand{\poplogoinline}[1]{{\poplogo\fontsize{#1}{#1}\selectfont\color{popink}OnBuch\color{poporange}+}}
\newcommand{\signatureonbuch}{%
  \begin{tcolorbox}[enhanced,colback=popink,colframe=popink,boxrule=2.2pt,arc=10pt,
      drop shadow=poporange,left=14pt,right=14pt,top=10pt,bottom=10pt,before skip=0pt,after skip=0pt]
    \tikz[baseline=-0.7ex]\node[circle,fill=popgreen,draw=popcream,line width=1.3pt,minimum size=22pt,inner sep=0]{\popcheck{white}};%
    \hspace{9pt}\begin{minipage}[c]{0.62\linewidth}
      {\popxbold\fontsize{12}{13}\selectfont\color{popcream}Signed\ \textbullet\ {\poplogo OnBuch\color{poporange}+}}\par\vspace{2pt}
      {\raggedright\popsemi\fontsize{8}{10.5}\selectfont\color{popline}\MakeUppercase{OnBuch+ teaching team}\\\MakeUppercase{Follows the official MINESEC syllabus}\par}
    \end{minipage}\hfill
    {\poplogo\fontsize{22}{22}\selectfont\color{popcream}OnBuch\color{poporange}+}
  \end{tcolorbox}%
}

% ============================================================
% Page de garde
% #1 titre · #2 matière · #3 niveau · #4 module · #5 n° leçon · #6 durée ·
% #7 description / objectifs (texte ou itemize)
% ============================================================
\newcommand{\pagegarde}[7]{%
  \thispagestyle{empty}
  \newgeometry{a4paper,margin=1.7cm,top=1.6cm,bottom=1.4cm}%
  \begin{tikzpicture}[remember picture,overlay]
    \fill[poporange!18] ([xshift=-3.2cm,yshift=-3.6cm]current page.north east) circle (4.6cm);
    \fill[poppurple!14] ([xshift=2.4cm,yshift=7.6cm]current page.south west) circle (3.8cm);
  \end{tikzpicture}%
  {\poplogo\fontsize{30}{30}\selectfont\color{popink}OnBuch\color{poporange}+}\hfill
  \raisebox{0.6ex}{\poppill{Signed course \textbullet\ Premium}{popink}{popcream}}\par
  \vspace{1pt}\popsquiggle{poppurple}\par
  \vspace{8pt}
  \begin{tcolorbox}[enhanced,colback=poporange,colframe=popink,boxrule=2.8pt,arc=13pt,
      drop shadow=popink,left=18pt,right=18pt,top=12pt,bottom=13pt,after skip=10pt]
    \poppill{#2 \textbullet\ #3}{popink}{popcream}\hspace{5pt}\poppill{#4}{white}{popink}\par\vspace{8pt}
    {\popdisplay\fontsize{14}{14}\selectfont\color{popink}CHAPTER #5}\par\vspace{4pt}
    {\raggedright\hyphenpenalty=10000\exhyphenpenalty=10000\popdisplay\fontsize{25}{27}\selectfont\color{popcream}\MakeUppercase{#1}\par}
    \vspace{8pt}
    {\popxbold\fontsize{9.5}{9.5}\selectfont\color{popink}\MakeUppercase{Suggested time: #6}}
  \end{tcolorbox}
  \begin{tcolorbox}[enhanced,colback=white,colframe=popink,boxrule=2.2pt,arc=10pt,
      drop shadow=popink,left=16pt,right=16pt,top=10pt,bottom=10pt,after skip=8pt]
    \poppill{In this chapter}{poppurple}{white}\par\vspace{5pt}
    \popsemi\fontsize{9.3}{12.6}\selectfont\color{popink}#7
  \end{tcolorbox}
  \noindent\begin{minipage}[t]{0.487\linewidth}
    \begin{tcolorbox}[enhanced,colback=popgreen,colframe=popink,boxrule=2.2pt,arc=10pt,
        drop shadow=popink,left=11pt,right=11pt,top=10pt,bottom=10pt,height=3.1cm,valign=center]
      \tcbox[colback=white,colframe=popink,boxrule=1.3pt,arc=5pt,boxsep=0pt,nobeforeafter,
        left=3pt,right=3pt,top=3pt,bottom=3pt,valign=center]{\qrcode[height=1.9cm]{https://play.google.com/store/apps/details?id=com.luvvix.onbuch&hl=en}}%
      \hspace{8pt}\begin{minipage}[c]{0.5\linewidth}
        {\popdisplay\fontsize{11}{12.5}\selectfont\color{white}\MakeUppercase{Download OnBuch}}\par\vspace{4pt}
        {\raggedright\popsemi\fontsize{8.4}{11}\selectfont\color{white}Free \textbullet\ Google Play\\AI tutor, past papers, results\par}
      \end{minipage}
    \end{tcolorbox}
  \end{minipage}\hfill
  \begin{minipage}[t]{0.487\linewidth}
    \begin{tcolorbox}[enhanced,colback=poppurple,colframe=popink,boxrule=2.2pt,arc=10pt,
        drop shadow=popink,left=11pt,right=11pt,top=10pt,bottom=10pt,height=3.1cm,valign=center]
      \tikz[baseline=-0.7ex]\node[circle,fill=white,draw=popink,line width=1.3pt,minimum size=1.6cm,inner sep=0]{\popdisplay\fontsize{24}{24}\selectfont\color{poppurple}+};%
      \hspace{8pt}\begin{minipage}[c]{0.6\linewidth}
        {\popdisplay\fontsize{11}{12.5}\selectfont\color{white}\MakeUppercase{Go OnBuch+}}\par\vspace{4pt}
        {\raggedright\popsemi\fontsize{8.4}{11}\selectfont\color{white}All signed courses, unlimited AI tutor, PDF sheets — OnBuch+ tab in the app\par}
      \end{minipage}
    \end{tcolorbox}
  \end{minipage}
  \vspace{8pt}
  \popcontenu
  \vfill
  \signatureonbuch
  \restoregeometry
  \newpage
}

% Rangée « Ce cours contient » (page de garde)
\newcommand{\poptile}[3]{% #1 couleur #2 gros texte #3 légende
  \begin{tcolorbox}[enhanced,colback=#1,colframe=popink,boxrule=2pt,arc=9pt,shadow={2.5pt}{-2.5pt}{0pt}{popink},
      left=6pt,right=6pt,top=9pt,bottom=9pt,halign=center,valign=center,height=1.75cm,width=\linewidth,nobeforeafter]
    {\popdisplay\fontsize{12.5}{13}\selectfont\color{white}\MakeUppercase{#2}}\par\vspace{3pt}
    {\centering\popsemi\fontsize{7.8}{9.5}\selectfont\color{white}#3\par}
  \end{tcolorbox}}
\newcommand{\popcontenu}{%
  \par\noindent{\popxboldls\fontsize{7.5}{7.5}\selectfont\color{popmuted}THIS SIGNED COURSE CONTAINS}\par\vspace{7pt}
  \noindent\begin{minipage}[t]{0.235\linewidth}\poptile{popblue}{Course}{complete, illustrated, step by step}\end{minipage}\hfill
  \begin{minipage}[t]{0.235\linewidth}\poptile{poporange}{Methods}{and worked examples}\end{minipage}\hfill
  \begin{minipage}[t]{0.235\linewidth}\poptile{poppink}{Exercises}{exam style + solutions}\end{minipage}\hfill
  \begin{minipage}[t]{0.235\linewidth}\poptile{popgreen}{Summary sheet}{the essentials to remember}\end{minipage}\par}

% Sommaire
\newtcolorbox{sommaire}{enhanced,colback=white,colframe=popink,boxrule=2.2pt,arc=10pt,
  drop shadow=popink,left=18pt,right=18pt,top=13pt,bottom=13pt,before skip=6pt,after skip=20pt,
  before upper={\poppill{Contents}{poppurple}{white}\par\vspace{9pt}\popsemi\fontsize{10.6}{14.5}\selectfont\color{popink}}}

% Signature de fin de cours — poussée tout en bas de la dernière page
\newcommand{\findecours}{%
  \par\vfill\nopagebreak
  \begin{center}\popsquiggle{poporange}\end{center}\vspace{4pt}
  \signatureonbuch
}
\AtBeginDocument{\def\llbracket{\mathopen{[\mkern-3.5mu[}}\def\rrbracket{\mathclose{]\mkern-3.5mu]}}} % intervalles d'entiers (glyphe ⟦ absent de la police)
% Glyphes absents de Fira Math : redéfinis à partir de glyphes disponibles
\AtBeginDocument{%
  \def\setminus{\mathbin{\backslash}}\let\smallsetminus\setminus
  \def\vdots{\mathord{\vbox{\baselineskip3.2pt\lineskiplimit0pt\kern4pt\hbox{.}\hbox{.}\hbox{.}}}}%
  \def\ddots{\mathinner{\mkern1mu\raise6.5pt\hbox{.}\mkern2mu\raise3.6pt\hbox{.}\mkern2mu\raise0.7pt\hbox{.}\mkern1mu}}%
  \def\triangle{\mathord{\scalebox{0.9}{$\symup{\Delta}$}}}%
  \def\nabla{\mathord{\raisebox{\depth}{\scalebox{1}[-1]{$\symup{\Delta}$}}}}%
  \def\checkmark{\popcheck{popdark}}%
  \def\star{\mathbin{\ast}}\def\bigstar{\mathord{\ast}}%
  \def\diamond{\mathbin{\scalebox{0.6}{$\Diamond$}}}%
  \def\bigcup{\mathop{\mathchoice{\raisebox{-0.3em}{\scalebox{1.7}{$\cup$}}}{\scalebox{1.25}{$\cup$}}{\cup}{\cup}}\displaylimits}%
  \def\bigcap{\mathop{\mathchoice{\raisebox{-0.3em}{\scalebox{1.7}{$\cap$}}}{\scalebox{1.25}{$\cap$}}{\cap}{\cap}}\displaylimits}%
  \def\lesseqqgtr{\mathrel{\lessgtr}}\def\bigoplus{\mathop{\oplus}\displaylimits}%
}
\providecommand{\ssi}{if and only if} % abréviation fréquente
\AtBeginDocument{\providecommand{\wideparen}{\overparen}} % arc de cercle

% Fonctions usuelles que les modèles utilisent sans que amsmath les fournisse
\providecommand{\arsinh}{\operatorname{arsinh}}\providecommand{\arcosh}{\operatorname{arcosh}}
\providecommand{\artanh}{\operatorname{artanh}}\providecommand{\arsech}{\operatorname{arsech}}
\providecommand{\arcsch}{\operatorname{arcsch}}\providecommand{\arcoth}{\operatorname{arcoth}}
\providecommand{\arcsinh}{\operatorname{arsinh}}\providecommand{\arccosh}{\operatorname{arcosh}}
\providecommand{\arctanh}{\operatorname{artanh}}\providecommand{\sech}{\operatorname{sech}}
\providecommand{\csch}{\operatorname{csch}}\providecommand{\arcsec}{\operatorname{arcsec}}
\providecommand{\arccsc}{\operatorname{arccsc}}\providecommand{\arccot}{\operatorname{arccot}}
\providecommand{\sgn}{\operatorname{sgn}}\providecommand{\lcm}{\operatorname{lcm}}
\providecommand{\Var}{\operatorname{Var}}\providecommand{\Cov}{\operatorname{Cov}}
\providecommand{\permil}{\ensuremath{{}^{0}\!/_{00}}}\providecommand{\textperthousand}{\ensuremath{{}^{0}\!/_{00}}}

%% FIGFIX-BEGIN v3 — correctifs globaux des figures (docs/onbuch-generation/tools/figfix)
\usepackage{adjustbox}\usepackage{etoolbox}
% toute figure TikZ/pgfplots plus large que la ligne est réduite à la largeur disponible
\BeforeBeginEnvironment{tikzpicture}{\begin{adjustbox}{max width=\linewidth}}
\AfterEndEnvironment{tikzpicture}{\end{adjustbox}}
% légendes pgfplots lisibles par-dessus les courbes
\pgfplotsset{every axis/.append style={legend style={fill=white,fill opacity=0.92,text opacity=1,draw=black!15,inner sep=3pt}}}
% étiquettes d'arêtes (syntaxe edge["..."]) avec fond blanc
\tikzset{every edge quotes/.append style={fill=white,inner sep=1.2pt}}
% bulles de cartes mentales : texte à la ligne dans la bulle, bulle un peu plus grande
\tikzset{every concept/.append style={text width=2.0cm,align=center,minimum size=2.3cm,inner sep=2pt}}
%% FIGFIX-END
\begin{document}

\pagegarde{Using formatting, editing features of Word processors for producing documents}{ICT}{Upper Sixth}{Upper Sixth — ICT}{21}{7 periods}{This chapter develops advanced mastery of word processors for producing well-formatted documents, then builds complete competence in spreadsheets: formatting data, cell referencing, formulae, functions, charts, linked worksheets and what-if analysis. It prepares learners for the practical and theory papers of the GCE Advanced Level in ICT.
\vspace{4pt}
\begin{multicols}{2}\small
\begin{itemize}
\item By the end of this chapter I can produce a professional document using formatting and editing features of a word processor (styles, tables, mail merge, headers/footers).
\item By the end of this chapter I can enter, format, sort, filter and manage data in a spreadsheet.
\item By the end of this chapter I can distinguish relative, absolute and mixed cell references and use them correctly in formulae.
\item By the end of this chapter I can write basic formulae and use advanced functions (IF, VLOOKUP, SUMIF, COUNTIF, statistical functions).
\item By the end of this chapter I can create and customise charts that correctly represent worksheet data.
\item By the end of this chapter I can link data across multiple worksheets and workbooks.
\end{itemize}\end{multicols}}

\begin{sommaire}
\begin{multicols}{2}
\begin{enumerate}
\item Producing Documents with Word Processing Features
\item Formatting and Managing Spreadsheet Data
\item Cell Referencing and Basic Formulae
\item Advanced Functions, Charts and Linked Worksheets
\item What-If Analysis and Integration Activities
\item Integration activity
\item Methods and worked examples
\item Exercises
\item Solutions to the exercises
\item Summary sheet
\end{enumerate}
\end{multicols}
\end{sommaire}

\begin{objectifs}
\begin{itemize}
\item By the end of this chapter I can produce a professional document using formatting and editing features of a word processor (styles, tables, mail merge, headers/footers).
\item By the end of this chapter I can enter, format, sort, filter and manage data in a spreadsheet.
\item By the end of this chapter I can distinguish relative, absolute and mixed cell references and use them correctly in formulae.
\item By the end of this chapter I can write basic formulae and use advanced functions (IF, VLOOKUP, SUMIF, COUNTIF, statistical functions).
\item By the end of this chapter I can create and customise charts that correctly represent worksheet data.
\item By the end of this chapter I can link data across multiple worksheets and workbooks.
\item By the end of this chapter I can perform what-if analysis using Goal Seek, scenarios and data tables.
\item By the end of this chapter I can integrate word processing and spreadsheet outputs into a single professional report.
\end{itemize}
\end{objectifs}

\begin{prerequis}
\begin{itemize}
\item Basic computer operations: files, folders, saving and opening documents (Lower Sixth).
\item Introduction to word processing: typing, selecting, copying and pasting text.
\item Introduction to spreadsheets: cells, rows, columns, worksheets and simple data entry.
\item Basic mathematical operations and percentages.
\item Keyboard and mouse handling skills acquired in Lower Sixth ICT practical sessions.
\end{itemize}
\end{prerequis}

\newpage

\coursec{Producing Documents with Word Processing Features}

At the end of every term, the bursar of a secondary school in Bamenda spends three days retyping fee records, recalculating totals by hand, and redrawing the same tables in Word. One wrong figure in a receipt book can trigger a dispute with a parent at the next PTA meeting. Meanwhile, a cybercafé operator in Douala produces the same reports in thirty minutes using Word formatting and Excel formulas that update themselves. What does he know that the bursar does not? He knows that a \cle{word processor} is not a typewriter: it is a document-production machine in which formatting is \emph{structured}, repetitive work is \emph{automated}, and one master document can generate hundreds of personalised copies. This section turns you into that efficient operator. We begin with the tool every GCE candidate must master: the word processor (Microsoft Word being the reference software in the syllabus).

\courssub{The Word Processor Interface: Knowing Your Workshop}

A carpenter who cannot find his tools wastes the day. The same is true of a Word user who cannot read the screen. When Word opens, the window presents several zones, each with a precise role.

\begin{definition}[Key zones of the Word window]
The \cle{ribbon} is the strip of tabs (Home, Insert, Layout, References, Mailings, Review, View) at the top of the window; each tab groups related commands into \emph{groups}. The \cle{ruler} (horizontal and vertical) shows margins, indents and tab stops. The \cle{insertion point} is the blinking vertical bar showing where the next typed character will appear. The \cle{status bar}, at the bottom, displays the current page number, word count, language and zoom level.
\end{definition}

Word also offers several \cle{views}, chosen from the View tab or the status bar:

\begin{itemize}
\item \textbf{Print Layout view}: shows the document exactly as it will be printed (margins, headers, page breaks visible). This is the default working view.
\item \textbf{Draft view}: hides headers, footers and page boundaries so you can focus on the text itself; useful for fast typing and editing.
\item \textbf{Read Mode, Web Layout and Outline views}: for reading on screen, previewing web pages, and reorganising long documents by headings respectively.
\end{itemize}

\begin{popfigure}
\begin{tikzpicture}[font=\small]
\draw[thick, popink, rounded corners=2pt] (0,0) rectangle (12,8);
\draw[fill=popblueL, draw=popink] (0,7.3) rectangle (12,8);
\node[anchor=west] at (0.2,7.65) {\textbf{Title bar} — document name};
\draw[fill=popblue!20, draw=popink] (0,6.4) rectangle (12,7.3);
\node[anchor=west] at (0.2,6.85) {\textbf{Ribbon}: Home \; Insert \; Layout \; References \; Mailings \; Review \; View};
\draw[fill=white, draw=popink] (0.4,5.9) rectangle (11.6,6.4);
\node[anchor=west] at (0.5,6.15) {\textbf{Horizontal ruler} (margins, indents, tabs)};
\draw[fill=white, draw=popline] (1.2,0.9) rectangle (10.8,5.9);
\draw[fill=popline!40, draw=popink] (0.4,0.9) rectangle (1.2,5.9);
\node[rotate=90] at (0.8,3.4) {\textbf{Vertical ruler}};
\draw[very thick, poporange] (3.0,4.6) -- (3.0,5.1);
\node[anchor=west, text=popdark] at (3.2,4.85) {\textbf{Insertion point} (blinking cursor)};
\node[anchor=west, text=popmuted] at (3.2,4.2) {Text typed here appears at the cursor.};
\draw[fill=popblueL, draw=popink] (0,0) rectangle (12,0.55);
\node[anchor=west] at (0.2,0.28) {\textbf{Status bar}: Page 2 of 5 \; Words: 1\,240 \; English (UK) \; Zoom 100\%};
\end{tikzpicture}
\legende{Mock-up of a Word window showing the ribbon, rulers, insertion point and status bar.}
\end{popfigure}

\begin{attention}[Do not confuse the ruler with margins]
Dragging the small markers on the ruler changes \emph{indentation} for the selected paragraphs, not the page margins. Margins are set in \textbf{Layout $\rightarrow$ Margins}. Mixing the two up is a classic source of crooked documents.
\end{attention}

\courssub{Character Formatting}

\cle{Character formatting} applies to individual letters, words or any selected text. All of it lives in the \textbf{Font} group of the Home tab.

\begin{itemize}
\item \textbf{Font (typeface)}: the design of the characters, e.g.\ Times New Roman (formal documents), Calibri, Arial. Choose one readable font and stick to it.
\item \textbf{Size}: measured in \emph{points} (pt). Body text is usually 11 or 12 pt; headings are larger.
\item \textbf{Bold, Italic, Underline} (Ctrl+B, Ctrl+I, Ctrl+U): used for emphasis. Underline is largely reserved for hyperlinks in modern documents.
\item \textbf{Superscript and subscript}: characters raised above or lowered below the baseline, e.g.\ $x^{2}$ (superscript) or H$_{2}$O (subscript). Essential in science reports.
\item \textbf{Text effects}: strikethrough, small caps, highlight colour, font colour, text outline and shadow.
\end{itemize}

\begin{exemplebox}[Formatting a chemistry note]
To type ``The formula of water is H$_2$O and $E = mc^{2}$'', select the ``2'' in H2O and click the subscript button (or Ctrl+=); select the ``2'' in $mc^2$ and click superscript (Ctrl+Shift+=). The formatting applies only to the selected characters.
\end{exemplebox}

\courssub{Paragraph Formatting}

\cle{Paragraph formatting} applies to whole paragraphs (a paragraph ends wherever the \emph{Enter} key was pressed). It is found in the \textbf{Paragraph} group of the Home tab and in the Paragraph dialog box.

\begin{itemize}
\item \textbf{Alignment}: left (default), centred, right, or \emph{justified} (both edges straight — used in newspapers and official letters).
\item \textbf{Line spacing}: the vertical space between lines (1.0, 1.5, 2.0\ldots). GCE reports are often required in 1.5 spacing for readability.
\item \textbf{Indentation}: the distance of a paragraph from the margin. The \emph{first-line indent} starts the first line further right; a \emph{hanging indent} does the opposite (used in bibliographies).
\item \textbf{Spacing before/after}: space inserted between paragraphs without pressing Enter.
\item \textbf{Bullets and numbering}: automatic lists. Numbering renumbers itself when items are inserted or deleted — a huge advantage over typing ``1.'', ``2.'' by hand.
\item \textbf{Borders and shading}: lines around, and background colour behind, a paragraph — useful for highlighting a warning or a title.
\end{itemize}

\begin{attention}[The Enter-key disease]
Beginners press \emph{Enter} several times to create space between paragraphs or to jump to a new page. This is wrong: the moment you add or delete text above, the whole layout collapses and pages shift unpredictably. Use \textbf{paragraph spacing} for gaps and \textbf{Layout $\rightarrow$ Breaks $\rightarrow$ Page Break} (Ctrl+Enter) to force a new page. A page break is a solid instruction; empty paragraphs are a house of cards.
\end{attention}

\courssub{Page Setup: Margins, Orientation, Size, Columns and Breaks}

Before printing, the page itself must be configured from the \textbf{Layout} tab:

\begin{itemize}
\item \textbf{Margins}: the white space around the text (Normal = 2.54 cm on all sides). Official letters often use wider left margins for binding.
\item \textbf{Orientation}: \emph{Portrait} (taller than wide, default) or \emph{Landscape} (wider than tall, good for large tables and certificates).
\item \textbf{Paper size}: A4 is the standard in Cameroon; Letter is the American size.
\item \textbf{Columns}: text can flow in two or three newspaper-style columns (Layout $\rightarrow$ Columns).
\item \textbf{Breaks}: a \emph{page break} starts a new page; a \emph{section break} starts a new \emph{section}, which can have its own orientation, margins, columns and headers. For example, one landscape page holding a wide table inside a portrait report requires a section break before and after it.
\end{itemize}

\courssub{Headers, Footers, Page Numbers, Footnotes and Endnotes}

The \cle{header} is the repeated zone at the top of every page; the \cle{footer} is its counterpart at the bottom. Double-click the top or bottom of a page (or Insert $\rightarrow$ Header/Footer) to edit them. Typical content: school name, document title, date. \cle{Page numbers} (Insert $\rightarrow$ Page Number) update automatically — you never type them by hand.

A \cle{footnote} is a note printed at the \emph{bottom of the page} where its reference mark appears; an \cle{endnote} is collected at the \emph{end of the document}. Both are inserted via References $\rightarrow$ Insert Footnote/Endnote and are numbered automatically. They are used to cite sources or add side remarks without interrupting the text.

\courssub{Styles, Themes and the Automatic Table of Contents}

Here lies the single most powerful idea of this section: \emph{never format headings by hand}.

\begin{definition}[Style]
A \cle{style} is a named set of formatting characteristics (font, size, colour, spacing, alignment) that is applied consistently to text. Examples: \emph{Heading 1}, \emph{Heading 2}, \emph{Title}, \emph{Normal}. A \cle{theme} is a coordinated set of colours and fonts applied to the whole document at once.
\end{definition}

Why are styles so important? First, \textbf{consistency}: every chapter heading looks identical. Second, \textbf{instant change}: modify the style once and every heading in a 60-page report updates at once. Third — and this is decisive — Word uses heading styles to build an \emph{automatic table of contents}.

\begin{methode}[Creating an automatic table of contents]
\begin{enumerate}
\item Apply \textbf{Heading 1} to each chapter title and \textbf{Heading 2} to each sub-title (Home tab, Styles gallery).
\item Click where the table of contents must appear (usually after the title page).
\item Go to \textbf{References $\rightarrow$ Table of Contents} and choose an automatic style.
\item After editing the document, right-click the table and choose \textbf{Update Field $\rightarrow$ Update entire table} so titles and page numbers are refreshed.
\end{enumerate}
\end{methode}

\courssub{Tables in Word}

A \cle{table} organises information in rows and columns; the intersection of a row and a column is a \emph{cell}. Insert one via Insert $\rightarrow$ Table (choose the grid size, or \emph{Insert Table} for exact dimensions). Key operations:

\begin{itemize}
\item \textbf{Merging cells}: select several cells, right-click $\rightarrow$ \emph{Merge Cells} — used for a title spanning several columns.
\item \textbf{Splitting cells}: divides one cell into several rows or columns.
\item \textbf{Sorting}: Table Tools $\rightarrow$ Layout $\rightarrow$ \emph{Sort} arranges rows alphabetically or numerically by any column (e.g.\ class list by marks, descending).
\item \textbf{Simple formulas}: Table Tools $\rightarrow$ Layout $\rightarrow$ \emph{Formula} accepts expressions such as \texttt{=SUM(ABOVE)} or \texttt{=AVERAGE(LEFT)} to total a column of figures — handy for a fee table, though a spreadsheet is the proper tool for heavy calculations (next sections).
\end{itemize}

\begin{exoresolu}[Producing a fee summary table]
The bursar must present a table of fees collected: Form Five 4\,500\,000 FCFA, Lower Sixth 3\,800\,000 FCFA, Upper Sixth 4\,100\,000 FCFA, with a total row.
\tcblower
\textbf{Solution.}
\begin{enumerate}
\item Insert $\rightarrow$ Table, 2 columns $\times$ 5 rows.
\item Row 1: type the headers \emph{Class} and \emph{Amount (FCFA)}; make them bold and shaded (Table Design tab).
\item Rows 2--4: enter the three classes and amounts.
\item Row 5, cell 1: type \emph{Total}. Merge nothing needed here.
\item Click in the last cell, Table Tools $\rightarrow$ Layout $\rightarrow$ Formula, accept \texttt{=SUM(ABOVE)}: Word displays \textbf{12\,400\,000}.
\item To rank classes, select rows 2--4 and use Sort on column 2, descending.
\end{enumerate}
Note: if a figure changes, the Word total does \emph{not} update automatically — select it and press F9 to refresh the field.
\end{exoresolu}

\courssub{Editing Features: Find and Replace, Proofing, Track Changes}

Professional documents are \emph{edited}, not merely typed.

\begin{itemize}
\item \textbf{Find} (Ctrl+F) locates a word or phrase. \textbf{Replace} (Ctrl+H) substitutes one text for another throughout the document — e.g.\ replacing ``2023'' by ``2024'' in a reused circular, or fixing a misspelt name everywhere at once.
\item \textbf{Spelling and grammar check} (Review $\rightarrow$ Spelling \& Grammar, or F7): red wavy underlines signal spelling errors, blue/green underlines signal grammar issues. Right-click an underlined word for suggestions.
\item \textbf{Track Changes} (Review tab): records every insertion and deletion so another person can \emph{accept} or \emph{reject} them. \textbf{Comments} attach margin notes to selected text without altering it. A supervisor correcting a student's project report uses exactly these two tools.
\end{itemize}

\courssub{Mail Merge: One Letter, Five Hundred Parents}

Back to our bursar. He must send a fee-reminder letter to 500 parents. Each letter is identical \emph{except} the parent's name, the child's name, the class and the balance due. Retyping 500 letters is three days of work and a guarantee of errors. \cle{Mail merge} does it in minutes.

\begin{definition}[Mail merge]
\cle{Mail merge} is a word-processing feature that combines a \emph{standard document} with a \emph{data source} to produce many \emph{personalised copies} automatically.
\end{definition}

\begin{aretenir}[The three components of mail merge — classic GCE question]
\begin{enumerate}
\item The \cle{main document}: the standard letter containing the fixed text and the \emph{merge fields} (placeholders such as \guillemotleft ParentName \guillemotright).
\item The \cle{data source}: the table of variable information (names, classes, balances) — often an Excel worksheet, an Access table or a Word table.
\item The \cle{merged document}: the final personalised output (500 individual letters), which can be printed, saved or sent by e-mail.
\end{enumerate}
Learn to \emph{list and define} these three components: it is one of the most frequent theory questions on this topic.
\end{aretenir}

\begin{popfigure}
\begin{tikzpicture}[font=\small, node distance=0.6cm]
\node[draw=popink, fill=popblueL, rounded corners=3pt, text width=3.4cm, align=center] (main) at (0,2.2) {\textbf{Main document}\\ standard letter + merge fields};
\node[draw=popink, fill=popgreenL, rounded corners=3pt, text width=3.4cm, align=center] (data) at (0,-0.6) {\textbf{Data source}\\ names, classes, balances};
\node[draw=poporange, very thick, circle, minimum size=1.5cm, align=center] (merge) at (5.6,0.8) {MERGE};
\node[draw=popink, fill=poppink!25, rounded corners=3pt, text width=3.6cm, align=center] (out) at (10.4,0.8) {\textbf{Merged documents}\\ 500 personalised letters};
\draw[-{Stealth}, very thick, popblue] (main.east) -- (merge.150);
\draw[-{Stealth}, very thick, popgreen!60!black] (data.east) -- (merge.210);
\draw[-{Stealth}, very thick, poporange] (merge.east) -- (out.west);
\end{tikzpicture}
\legende{The mail merge process: the main document and the data source are combined to produce personalised merged documents.}
\end{popfigure}

\begin{methode}[Producing personalised fee-reminder letters]
\begin{enumerate}
\item Type the standard letter in Word (the \emph{main document}).
\item Prepare the \emph{data source}: an Excel sheet with columns ParentName, ChildName, Class, Balance.
\item In Word: \textbf{Mailings $\rightarrow$ Start Mail Merge $\rightarrow$ Letters}.
\item \textbf{Select Recipients $\rightarrow$ Use an Existing List}, and choose the Excel file.
\item Place the cursor where each variable belongs and click \textbf{Insert Merge Field} (e.g.\ \guillemotleft ParentName\guillemotright, \guillemotleft Balance\guillemotright).
\item Click \textbf{Preview Results} to check a few records.
\item \textbf{Finish \& Merge $\rightarrow$ Print Documents} (or \emph{Edit Individual Documents} to save all 500 letters in one file).
\end{enumerate}
\end{methode}

\begin{savaistu}[Mail merge in Cameroonian life]
Banks, insurance companies and mobile money operators use mail-merge logic daily: the SMS you receive saying ``Dear \emph{Ngong Paul}, your balance is\ldots'' is generated by merging a standard message with a database of millions of customers. The same principle powers the personalised GCE registration slips printed for each candidate.
\end{savaistu}

\courssub{Printing and Exporting to PDF}

Before printing, always use \textbf{File $\rightarrow$ Print} and inspect the \emph{print preview}: check the page range, number of copies, printer, and whether to print one or both sides (duplex). To share a document that must not be modified and must look identical on every device, \textbf{export it to PDF} (File $\rightarrow$ Save As / Export $\rightarrow$ PDF). PDF preserves fonts, layout and images exactly — it is the format expected for official circulars and for submitting typed projects.

\begin{exercice}[Fee reminder project]
The principal of GBHS Bamenda asks you to send fee reminders to the parents of 40 Upper Sixth students.
\begin{enumerate}
\item Name the three components of the mail merge you will perform, and state the role of each.
\item Give two reasons why an automatic table of contents is preferable to a typed one in the end-of-year report.
\item Explain why inserting a page break is better than pressing Enter repeatedly at the end of a chapter.
\end{enumerate}
\end{exercice}

\begin{corrige}[Exercise 1]
\begin{enumerate}
\item \textbf{Main document}: the standard reminder letter with merge fields; \textbf{data source}: the table (e.g.\ Excel sheet) holding the 40 parents' names, students' names, classes and balances; \textbf{merged document}: the 40 personalised letters produced by combining the two.
\item (a) It updates automatically (Update Field) when titles or pages change; (b) it is generated from heading styles, so page numbers and titles are always exact, with no typing errors.
\item A page break is a fixed instruction: the next chapter always starts on a fresh page whatever edits occur above. Repeated Enters create empty paragraphs that shift whenever text is added or deleted, destroying the layout.
\end{enumerate}
\end{corrige}

\begin{aretenir}[Section summary]
\begin{itemize}
\item The Word window: ribbon, rulers, insertion point, status bar; Print Layout vs Draft view.
\item Character formatting (font, size, bold/italic/underline, super/subscript) vs paragraph formatting (alignment, spacing, indentation, bullets, borders).
\item Page setup: margins, orientation, paper size, columns, page and section breaks.
\item Headers/footers, automatic page numbers, footnotes and endnotes.
\item Styles guarantee consistency and feed the automatic table of contents.
\item Word tables: merge/split cells, sort, simple formulas such as \texttt{=SUM(ABOVE)}.
\item Editing tools: Find/Replace, spelling and grammar, track changes and comments.
\item Mail merge = main document + data source $\rightarrow$ merged documents; finish with printing or PDF export.
\end{itemize}
\end{aretenir}

\coursec{Formatting and Managing Spreadsheet Data}

In the previous section we used a word processor to \emph{present} information: letters, reports and tables that are essentially fixed once printed. But a school report card, a fee-collection record or a stock list at a Douala market is never fixed: values change, new rows appear, and we constantly need to rank, filter and check the data. This is exactly why the \cle{spreadsheet} exists. A spreadsheet such as Microsoft Excel, LibreOffice Calc or Google Sheets stores data in a grid where every value can be formatted, sorted, filtered and checked automatically. In this section we learn to \emph{enter data correctly}, \emph{format it professionally} and \emph{manage it efficiently} --- the foundations on which formulae, charts and what-if analysis (next sections) will rest.

\courssub{Spreadsheet Structure Revisited}

\begin{definition}[Workbook and worksheet]
A \cle{workbook} is the spreadsheet \emph{file} itself (for example \texttt{Term2\_Results.xlsx}). A \cle{worksheet} is a single grid of rows and columns \emph{within} a workbook file. A workbook may contain several worksheets, shown as tabs at the bottom of the window (e.g.\ ``Form 5'', ``Lower Sixth'', ``Upper Sixth'').
\end{definition}

The intersection of a row and a column is a \cle{cell}. Columns are labelled with letters (A, B, C, \dots, Z, AA, AB, \dots) and rows with numbers (1, 2, 3, \dots). Each cell therefore has a unique \cle{cell address} (or \emph{reference}): the column letter followed by the row number, for example \texttt{B3}. A rectangular group of cells is a \cle{range}, written with a colon: \texttt{A1:D10} means all cells from \texttt{A1} to \texttt{D10}.

Two screen elements are essential:
\begin{itemize}
\item the \cle{Name Box} (left of the formula bar) shows the address of the active cell;
\item the \cle{Formula Bar} shows the \emph{true content} of the active cell --- which may differ from what the cell \emph{displays} after formatting.
\end{itemize}

\begin{popfigure}
\begin{center}
\begin{tikzpicture}[scale=0.92, every node/.style={font=\small}]
% column headers
\foreach \c/\x in {A/0.9, B/2.7, C/4.5, D/6.3} {
  \fill[popblueL] (\x-0.9,4.4) rectangle (\x+0.9,5.0);
  \node at (\x,4.7) {\textbf{\c}};
}
% row headers
\foreach \r/\y in {1/3.7, 2/2.9, 3/2.1, 4/1.3} {
  \fill[popblueL] (-1.4,\y-0.4) rectangle (-0.0,\y+0.4);
  \node at (-0.7,\y) {\textbf{\r}};
}
% grid
\draw[popline, step=1.8] (0,0.9) grid (7.2,4.4);
\draw[popline] (0,0.9) rectangle (7.2,4.4);
% highlight B3
\fill[poporange, opacity=0.35] (1.8,1.7) rectangle (3.6,2.5);
\draw[popdark, very thick] (1.8,1.7) rectangle (3.6,2.5);
\node at (2.7,2.1) {B3};
% name box and formula bar
\fill[popgreenL] (-1.4,5.6) rectangle (0.6,6.2);
\node at (-0.4,5.9) {\textbf{B3}};
\node[popmuted] at (0.9,5.9) {$f_x$};
\draw[popline] (1.3,5.6) rectangle (7.2,6.2);
\node[anchor=west] at (1.5,5.9) {Ngo Bell Martine};
\node[popmuted, anchor=west] at (-1.4,6.5) {Name Box};
\node[popmuted, anchor=west] at (1.5,6.5) {Formula Bar};
\end{tikzpicture}
\end{center}
\legende{A worksheet grid: the active cell \texttt{B3} is the intersection of column B and row 3. The Name Box confirms the address; the Formula Bar shows the real content.}
\end{popfigure}

\begin{attention}[Displayed value vs real content]
Formatting changes only the \emph{display}. If a cell contains $12500.5$ and is formatted as currency with no decimals, it \emph{shows} \texttt{12\,501 FCFA} but the stored value is still $12500.5$. Always check the Formula Bar when a result surprises you.
\end{attention}

\courssub{Data Types and Correct Entry}

A spreadsheet does not treat everything alike: it recognises several \cle{data types}, and the type determines what you can do with the value (calculate, sort chronologically, etc.).

\begin{center}
\begin{tabularx}{\linewidth}{|l|X|X|l|}
\hline
\rowcolor{popblueL} \textbf{Type} & \textbf{Examples} & \textbf{Correct entry} & \textbf{Default alignment} \\
\hline
Text & Names, towns, codes & Typed directly & Left \\
\hline
Number & 15, 12500.5, $-3$ & Digits, one decimal \emph{point} or comma per regional settings & Right \\
\hline
Date & 12/01/2026 & \texttt{12/01/2026} or \texttt{12-Jan-2026} & Right \\
\hline
Currency & 12\,500 FCFA & Type the number, then apply the currency \emph{format} & Right \\
\hline
Percentage & 85\,\% & Type \texttt{85\%} or $0.85$ with \% format & Right \\
\hline
\end{tabularx}
\end{center}

\begin{attention}[Numbers typed as text --- the classic trap]
If you type an apostrophe before a number (\texttt{'12500}), or add a letter or a space inside it (\texttt{12500F}), the spreadsheet stores it as \emph{text}: it stays \textbf{left-aligned} and Excel shows a small \textbf{green triangle} in the corner of the cell. Such ``numbers'' are \textbf{ignored by calculations} --- a SUM over them gives a wrong total. The cure: re-enter the value as a pure number, or use the warning menu ``Convert to Number''.
\end{attention}

\begin{aretenir}[Entry rules]
\begin{itemize}
\item One piece of information per cell (do not type ``Etoudi Paul 14/20'' in one cell).
\item Dates are stored internally as numbers, so they can be sorted and subtracted.
\item Use the \emph{format} for FCFA and \%, never type the unit inside the cell.
\end{itemize}
\end{aretenir}

\courssub{Cell Formatting}

Formatting makes a sheet readable and professional. Select the cells, then use the \textbf{Home} tab or right-click $\rightarrow$ \textbf{Format Cells}.

\textbf{Number formats.} The \emph{Number} category controls how values appear:
\begin{itemize}
\item \textbf{Currency}: \texttt{12\,500 FCFA} --- choose the symbol and the number of decimal places;
\item \textbf{Percentage}: the stored value $0.85$ displays as \texttt{85\,\%};
\item \textbf{Date}: the same stored date can display as \texttt{12/01/2026}, \texttt{12-Jan-26} or \texttt{Monday 12 January 2026}.
\end{itemize}

\textbf{Alignment.} Horizontal (left, centre, right) and vertical alignment, plus \emph{Wrap Text} (long text stays inside its cell on several lines) and \emph{Merge \& Centre} (one title across several columns --- use sparingly, merged cells complicate sorting).

\textbf{Borders and fill colours.} Borders delimit the table (outside thick border, thin inside lines); fill colours distinguish headers or categories. Choose sober, contrasting colours: a report card is not a carnival poster.

\textbf{Cell styles.} A \cle{cell style} is a named bundle of formats (font, fill, border, number format) applied in one click, e.g.\ the built-in ``Heading 1'' or ``Total'' styles. Styles guarantee \emph{consistency} across a whole workbook.

\begin{exemplebox}[Formatting a marks table]
For the Upper Sixth class council, format the column of averages as \emph{Number, 2 decimals}; the column of fees paid as \emph{Currency, FCFA, 0 decimals}; give the header row a dark blue fill with bold white centred text; and put a thick border around the whole table. The sheet is immediately readable by the principal.
\end{exemplebox}

\courssub{Adjusting Rows and Columns}

Data rarely fits the default grid. Right-click a column letter or row number (or use \textbf{Home $\rightarrow$ Format}) to:
\begin{itemize}
\item \textbf{Width / Height}: set an exact size, or double-click the boundary between two headers for \textbf{AutoFit} (the column takes the width of its longest content);
\item \textbf{Insert / Delete}: inserting a column shifts the others right; inserting a row shifts rows down. Existing formulae adjust automatically;
\item \textbf{Hide / Unhide}: hiding a column (e.g.\ raw test marks) keeps its data and formulae active but removes it from view and from printing. Select the neighbouring columns and choose \emph{Unhide} to bring it back.
\end{itemize}

\begin{attention}[\#\#\#\#\# is not an error]
When a numeric column is too narrow, the cell shows \texttt{\#\#\#\#}. The value is intact: simply widen the column. (Text, by contrast, is merely truncated visually.)
\end{attention}

\courssub{AutoFill and Custom Series}

Observe the small square at the bottom-right corner of the active cell: the \cle{fill handle}. Dragging it copies or \emph{extends} the content intelligently.

\begin{propriete}[AutoFill series recognition]
AutoFill automatically recognises and extends built-in series such as \emph{Monday, Tuesday, \dots} and \emph{January, February, \dots}, as well as numeric patterns: dragging two selected cells containing $2$ and $4$ produces $6, 8, 10, \dots$; dragging a date produces the following days. A single number dragged alone is simply \emph{copied} (hold Ctrl to increment instead).
\end{propriete}

\begin{exemplebox}[Preparing a term planner]
Type \texttt{Week 1} in A1 and drag the fill handle down: the sheet writes \texttt{Week 2, Week 3, \dots} Type \texttt{January} in B1 and drag right: \texttt{February, March, \dots} appear. You can also define \emph{custom lists} (File $\rightarrow$ Options $\rightarrow$ Advanced $\rightarrow$ Edit Custom Lists), e.g.\ the sequence of your school's classes: Form 5, Lower Sixth, Upper Sixth.
\end{exemplebox}

\courssub{Sorting Data}

\cle{Sorting} rearranges the rows of a table according to the values of one or several columns, without altering the data itself.

\begin{methode}[Sorting on multiple keys]
To rank students by class, then by merit inside each class:
\begin{enumerate}
\item Click any single cell inside the table (do \emph{not} select only one column --- that would scramble the rows!).
\item Open \textbf{Data $\rightarrow$ Sort}.
\item Tick \emph{My data has headers} so the header row stays on top.
\item First level: Sort by \textbf{Class}, order A to Z.
\item Click \textbf{Add Level}: Then by \textbf{Average}, order \textbf{Largest to Smallest}.
\item Optionally add a third level (e.g.\ \textbf{Name}, A to Z) to break ties. Click \textbf{OK}.
\end{enumerate}
\end{methode}

\begin{attention}[The selection trap]
If you select only the ``Average'' column before sorting, the spreadsheet may sort \emph{that column alone}, detaching the averages from the students' names --- a disaster for a report card. Always click one cell \emph{inside} the table and let the software detect the whole range, or select the entire table including all columns.
\end{attention}

\courssub{Filtering Data with AutoFilter}

Sorting reorders; \cle{filtering} \emph{hides} the rows you do not want to see. With \textbf{Data $\rightarrow$ Filter} (AutoFilter), a drop-down arrow appears on each header. Untick values to hide them, or open \emph{Number Filters / Text Filters} for \textbf{custom criteria}: \emph{Greater Than}, \emph{Between}, \emph{Contains}, \dots

\begin{popfigure}
\begin{center}
\begin{tikzpicture}[node distance=0.9cm, every node/.style={font=\small}]
\node[draw=popblue, fill=popblueL, rounded corners, text width=3.1cm, align=center, minimum height=1.2cm] (full) {\textbf{Full table}\\ 45 students, all classes};
\node[draw=poppurple, fill=poppurple!15, rounded corners, text width=3.1cm, align=center, minimum height=1.2cm, right=1.6cm of full] (crit) {\textbf{Criteria}\\ Class = ``USC''\\ Average $\geq 10$};
\node[draw=popgreen, fill=popgreenL, rounded corners, text width=3.1cm, align=center, minimum height=1.2cm, right=1.6cm of crit] (sub) {\textbf{Filtered subset}\\ 12 rows shown\\ others hidden};
\draw[-{Stealth}, very thick, popdark] (full) -- (crit);
\draw[-{Stealth}, very thick, popdark] (crit) -- (sub);
\end{tikzpicture}
\end{center}
\legende{The filtering process: criteria applied to the full table reveal only the matching subset; hidden rows are not deleted.}
\end{popfigure}

\begin{exemplebox}[Custom filter]
To see only students with an average \emph{between 10 and 14}: open the arrow on the \emph{Average} header $\rightarrow$ \emph{Number Filters} $\rightarrow$ \emph{Between\ldots} $\rightarrow$ enter $10$ and $14$ $\rightarrow$ OK. Combine it with \emph{Class = USC} on another column: filters on several columns cumulate (logical AND). Use \textbf{Data $\rightarrow$ Clear} to remove all filters.
\end{exemplebox}

\begin{attention}[Filtered data is hidden, not deleted]
Rows hidden by a filter still exist. But beware: copying a filtered range copies only the \emph{visible} cells, and printing a filtered sheet prints only the visible rows. Check the row numbers (blue, non-consecutive) to know a filter is active.
\end{attention}

\courssub{Data Validation}

The best error is the one that can never be typed. \cle{Data validation} (\textbf{Data $\rightarrow$ Data Validation}) restricts what a cell accepts:
\begin{itemize}
\item \textbf{List}: the cell offers a drop-down, e.g.\ \texttt{USC;LSC;Form 5} --- impossible to misspell a class name;
\item \textbf{Whole number / Decimal between}: e.g.\ a mark must lie between $0$ and $20$;
\item \textbf{Date between}: e.g.\ only dates of the current school year.
\end{itemize}
You can add an \emph{Input Message} (guidance shown when the cell is selected) and an \emph{Error Alert} (message shown when the rule is broken).

\begin{exemplebox}[Protecting a mark-entry sheet]
Select the mark column $\rightarrow$ Data Validation $\rightarrow$ Allow: \emph{Decimal}, between $0$ and $20$ $\rightarrow$ Error Alert: ``A mark must be between 0 and 20.'' A teacher who types $75$ by mistake is stopped immediately instead of corrupting the term averages.
\end{exemplebox}

\courssub{Conditional Formatting}

\begin{definition}[Conditional formatting]
\cle{Conditional formatting} is formatting (colour, font, icon) applied \emph{automatically} to a cell whenever its value satisfies a stated condition, and removed automatically when the condition ceases to hold.
\end{definition}

Via \textbf{Home $\rightarrow$ Conditional Formatting} you can use:
\begin{itemize}
\item \textbf{Highlight Cells Rules}: Greater Than, Less Than, Between, Equal To, Text that Contains;
\item \textbf{Top/Bottom Rules}: Top 10 items, Below Average;
\item \textbf{Data Bars, Colour Scales, Icon Sets}: visual heat-maps of a whole column.
\end{itemize}

\begin{exemplebox}[Marks below 10 in red]
Select the averages column $\rightarrow$ Conditional Formatting $\rightarrow$ Highlight Cells Rules $\rightarrow$ Less Than $\rightarrow$ type $10$ $\rightarrow$ choose ``Light Red Fill with Dark Red Text''. Every weak average turns red instantly; if a mark is corrected and the average rises above $10$, the red disappears by itself.
\end{exemplebox}

\begin{attention}[Conditional formatting is not a calculation]
Conditional formatting only \emph{colours} cells; it does not change values and cannot replace a formula. Also, piling up many contradictory rules on the same range makes the sheet unreadable --- keep two or three meaningful rules.
\end{attention}

\courssub{Managing Large Sheets: Freezing, Splitting, Printing}

\textbf{Freeze panes.} In a sheet of 300 students, the header row disappears when you scroll down. \textbf{View $\rightarrow$ Freeze Panes $\rightarrow$ Freeze Top Row} (or \emph{Freeze First Column}) keeps headers visible while the rest scrolls. To freeze both, select cell \texttt{B2} and choose \emph{Freeze Panes}: everything above and left of \texttt{B2} stays fixed.

\textbf{Split.} \textbf{View $\rightarrow$ Split} divides the window into independent panes with separate scroll bars --- useful to compare the top and the bottom of a long list simultaneously.

\textbf{Page setup and print area.} Before printing (\textbf{Page Layout} tab):
\begin{itemize}
\item choose \textbf{Orientation} (landscape for wide tables) and \textbf{paper size} (A4);
\item \textbf{Scale} the sheet to fit one page wide;
\item define the \textbf{Print Area} (select the range $\rightarrow$ \emph{Print Area $\rightarrow$ Set Print Area}) so that only the useful table prints, not your scratch calculations;
\item use \emph{Print Titles} to repeat the header row on every page.
\end{itemize}

\begin{savaistu}[Spreadsheets in Cameroonian schools]
Many colleges and high schools in Cameroon now prepare term report cards with Excel templates: marks are typed once, averages and ranks are computed automatically, and the honour roll is produced by filtering and sorting exactly as in this section. Mastering these skills is directly useful for any secretarial, accounting or teaching career.
\end{savaistu}

\begin{exoresolu}[Preparing the honour roll of Upper Sixth]
The sheet \texttt{Results} contains, in columns A to D: Name, Class, Average, Fees Paid (FCFA) for 45 students. (1) Format the sheet professionally. (2) Produce the list of USC students ranked by decreasing average, showing only those with an average of at least $10$, with failing averages of the whole sheet highlighted in red. (3) Prepare printing of the honour roll alone.
\tcblower
\textbf{Solution.}
\begin{enumerate}
\item \emph{Formatting.} Select A1:D1 $\rightarrow$ bold, white text, dark blue fill, centred. Select C2:C46 $\rightarrow$ Format Cells $\rightarrow$ Number, 2 decimals. Select D2:D46 $\rightarrow$ Currency, symbol FCFA, 0 decimals. Select A1:D46 $\rightarrow$ All Borders, then a thick outside border. AutoFit all column widths (double-click the header boundaries).
\item \emph{Red highlight.} Select C2:C46 $\rightarrow$ Conditional Formatting $\rightarrow$ Less Than $10$ $\rightarrow$ red fill.
\item \emph{Filtering.} Click inside the table $\rightarrow$ Data $\rightarrow$ Filter. On the \emph{Class} arrow, tick only \texttt{USC}. On the \emph{Average} arrow $\rightarrow$ Number Filters $\rightarrow$ Greater Than Or Equal To $10$.
\item \emph{Sorting.} Data $\rightarrow$ Sort $\rightarrow$ Sort by \emph{Average}, order \emph{Largest to Smallest} $\rightarrow$ OK. (Filtering and sorting combine safely.)
\item \emph{Printing.} Select the visible table $\rightarrow$ Page Layout $\rightarrow$ Print Area $\rightarrow$ Set Print Area; orientation Portrait; Scale: Fit to 1 page wide; Print Titles: repeat row 1. Print Preview to verify, then print.
\end{enumerate}
The result is a clean, ranked honour roll of USC students with averages $\geq 10$, ready for the notice board.
\end{exoresolu}

\begin{exercice}[School canteen stock]
The canteen of a college in Bafoussam keeps a sheet with columns: Item, Category, Quantity, Unit Price (FCFA), Expiry Date.
\begin{enumerate}
\item Explain how to guarantee that \emph{Category} contains only ``Drink'', ``Food'' or ``Soap''.
\item Describe how to display only items whose quantity is below 20, sorted by expiry date, earliest first.
\item The Quantity column shows green triangles and the total is wrong. Diagnose and fix.
\item State how to keep the header row visible while scrolling through 500 items.
\end{enumerate}
\end{exercice}

\begin{corrige}[Exercise 1]
\begin{enumerate}
\item Select the Category column $\rightarrow$ Data Validation $\rightarrow$ Allow: \emph{List} $\rightarrow$ Source: \texttt{Drink,Food,Soap}. A drop-down appears and any other entry is rejected with an error alert.
\item Data $\rightarrow$ Filter; on \emph{Quantity}: Number Filters $\rightarrow$ Less Than $20$; then Data $\rightarrow$ Sort by \emph{Expiry Date}, order \emph{Oldest to Newest} (dates sort chronologically because they are true date values).
\item The quantities were entered as text (left-aligned, green triangles), so SUM ignores them. Fix: select the column, use the warning icon $\rightarrow$ \emph{Convert to Number} (or retype the values), then recompute the total.
\item Click in row 2 (or any cell below the headers) $\rightarrow$ View $\rightarrow$ Freeze Panes $\rightarrow$ Freeze Top Row.
\end{enumerate}
\end{corrige}

\begin{aretenir}[Key points of the section]
\begin{itemize}
\item A workbook contains worksheets; a cell is identified by its address (column letter + row number); the Formula Bar shows the true content.
\item Enter data in the correct \emph{type}; apply currency and percentage through \emph{formats}, never by typing symbols.
\item AutoFill extends days, months and numeric series via the fill handle.
\item Sort on several keys with \textbf{Data $\rightarrow$ Sort $\rightarrow$ Add Level}; filter with AutoFilter and custom criteria --- hidden rows are not deleted.
\item Data validation prevents wrong entries; conditional formatting colours cells automatically according to rules.
\item Freeze panes for large sheets; set the print area, orientation and print titles before printing.
\end{itemize}
\end{aretenir}

With clean, well-formatted and well-managed data, we are now ready to make the spreadsheet \emph{calculate}: the next section introduces cell referencing and basic formulae.

\coursec{Cell Referencing and Basic Formulae}

In the previous section we learned how to enter, format and manage data in a worksheet: adjusting column widths, applying number formats, sorting and filtering. But a well-formatted table that only \emph{displays} numbers is no better than a sheet of paper. The real power of a spreadsheet such as Microsoft Excel or LibreOffice Calc is that it \emph{computes}. Once you tell the spreadsheet \emph{how} to compute, it does the work for you --- and redoes it instantly whenever the data change. The key that unlocks this power is the \cle{formula}, and the ingredient that makes formulae so flexible is the \cle{cell reference}. This section builds both ideas rigorously, from the structure of a formula to the subtle but examinable difference between relative, absolute and mixed references.

\courssub{The Structure of a Formula}

Imagine a shopkeeper in Mokolo market who sells 15 bags of groundnuts at \SI{2500}{FCFA} each. On paper she writes $15 \times 2500$. In a spreadsheet, the quantity 15 is typed in cell \texttt{B2} and the price 2500 in cell \texttt{C2}. Instead of typing the result 37500 by hand (which would become wrong the moment a value changes), she types in \texttt{D2}:

\[ \texttt{=B2*C2} \]

The spreadsheet reads the values stored in \texttt{B2} and \texttt{C2}, multiplies them, and displays the result. This is a formula.

\begin{definition}[Formula]
A \cle{formula} is an expression entered in a cell that \textbf{always begins with the equal sign} \texttt{=} and instructs the spreadsheet to perform a calculation. A formula is built from \cle{operands} (constants such as $5$, or cell references such as \texttt{B2}) combined by \cle{operators}.
\end{definition}

\begin{definition}[Cell reference]
A \cle{cell reference} is the \textbf{address of a cell} (column letter followed by row number, e.g.\ \texttt{B2}) used inside a formula so that the formula works with the \emph{content} of that cell rather than with a fixed number. A range of cells is written with a colon, e.g.\ \texttt{B2:B10} means all the cells from \texttt{B2} down to \texttt{B10}.
\end{definition}

The arithmetic operators available in spreadsheet formulae are:

\begin{center}
\begin{tabularx}{\linewidth}{|c|l|X|}
\hline
\rowcolor{popblueL} \textbf{Operator} & \textbf{Meaning} & \textbf{Example (result)} \\
\hline
\texttt{+} & Addition & \texttt{=2+3} gives $5$ \\
\hline
\texttt{-} & Subtraction & \texttt{=10-4} gives $6$ \\
\hline
\texttt{*} & Multiplication & \texttt{=6*7} gives $42$ \\
\hline
\texttt{/} & Division & \texttt{=20/4} gives $5$ \\
\hline
\texttt{\^\{\}} & Exponentiation (power) & \texttt{=2\^\{\}3} gives $8$ \\
\hline
\texttt{\%} & Percentage (divides by 100) & \texttt{=19.25\%} gives $0.1925$ \\
\hline
\end{tabularx}
\end{center}

\begin{attention}[The \texttt{*} is compulsory]
In mathematics we write $3x$ to mean $3 \times x$. A spreadsheet does \textbf{not} accept this: \texttt{=3A1} is invalid. You must always write the multiplication sign explicitly: \texttt{=3*A1}. Similarly, the power symbol is the caret \texttt{\^\{\}}, not a superscript: $x^2$ is written \texttt{=A1\^\{\}2}.
\end{attention}

\courssub{Order of Operations and Parentheses}

A formula such as \texttt{=2+3*4} could give $20$ (if we added first) or $14$ (if we multiplied first). Spreadsheets follow the same convention as mathematics, often remembered as \emph{BODMAS/PEMDAS}:

\begin{enumerate}
\item \textbf{Parentheses} first;
\item \textbf{Exponentiation} (\texttt{\^\{\}});
\item \textbf{Multiplication and division} (\texttt{*}, \texttt{/}), from left to right;
\item \textbf{Addition and subtraction} (\texttt{+}, \texttt{-}), from left to right.
\end{enumerate}

Thus \texttt{=2+3*4} gives $14$, whereas \texttt{=(2+3)*4} gives $20$. Parentheses override the default priority and make your intention unambiguous.

\begin{exemplebox}[Computing a weighted average]
A student scores 12/20 in ICT (coefficient 4) and 15/20 in Mathematics (coefficient 2), stored in \texttt{B2} and \texttt{C2}. The weighted total is:
\[ \texttt{=B2*4 + C2*2} \quad\text{gives}\quad 12\times 4 + 15\times 2 = 78. \]
To compute the average we must divide by the sum of coefficients, and parentheses are essential:
\[ \texttt{=(B2*4 + C2*2)/(4+2)} = \frac{78}{6} = 13. \]
Without parentheses, \texttt{=B2*4+C2*2/4+2} would give $48 + 7.5 + 2 = 57.5$ --- completely wrong.
\end{exemplebox}

\begin{propriete}[Automatic recalculation]
A formula \textbf{always recalculates automatically} whenever any cell it refers to is modified. If the price in \texttt{C2} changes from 2500 to 3000, the formula \texttt{=B2*C2} immediately displays the new result $45000$. This property is the fundamental reason spreadsheets are preferred to calculators for repetitive work. (Not examinable as a proof, but frequently tested as a concept.)
\end{propriete}

\courssub{Relative Referencing}

Here is the observation that changes everything. Suppose column \texttt{D} contains \texttt{=B2*C2} in row 2, and we \emph{copy} this formula down to rows 3, 4, 5 using the fill handle (the small square at the bottom-right corner of the selected cell). We do \emph{not} obtain \texttt{=B2*C2} everywhere. The spreadsheet adjusts the addresses automatically:

\begin{center}
\begin{tabular}{|c|c|c|}
\hline
\rowcolor{popblueL} \textbf{Cell} & \textbf{Formula after copying} & \textbf{Meaning} \\
\hline
\texttt{D2} & \texttt{=B2*C2} & quantity 2 $\times$ price 2 \\
\hline
\texttt{D3} & \texttt{=B3*C3} & quantity 3 $\times$ price 3 \\
\hline
\texttt{D4} & \texttt{=B4*C4} & quantity 4 $\times$ price 4 \\
\hline
\end{tabular}
\end{center}

\begin{definition}[Relative reference]
A \cle{relative reference} (written without any \texttt{\$}, e.g.\ \texttt{B2}) is a reference that is \textbf{adjusted automatically} when the formula is copied to another cell: the column letter changes when copying horizontally, and the row number changes when copying vertically. The spreadsheet remembers the \emph{position of the referenced cell relative to the formula}, not the address itself.
\end{definition}

Intuitively, \texttt{=B2*C2} typed in \texttt{D2} really means ``multiply the cell two columns to my left by the cell one column to my left''. When copied to \texttt{D3}, the same instruction becomes \texttt{=B3*C3}. This is exactly what we want when the same calculation must be repeated on many rows of data --- one formula, copied once, serves a thousand rows.

\courssub{Absolute Referencing with \$}

Relative referencing is wonderful --- until it is not. Suppose all prices must be increased by a VAT rate of $19.25\%$ stored once in cell \texttt{E1}. In \texttt{D2} we write \texttt{=C2*E1} and copy down. In \texttt{D3} the formula becomes \texttt{=C3*E2}: the price reference correctly moved, but the VAT reference \emph{also moved} to the empty cell \texttt{E2}, producing zero. We need a way to say: ``this reference must \textbf{never} move''.

\begin{definition}[Absolute reference]
An \cle{absolute reference} is a cell reference written with dollar signs, e.g.\ \texttt{\$E\$1}, that \textbf{never changes} when the formula is copied, whether vertically or horizontally. The \texttt{\$} before the column letter fixes the column; the \texttt{\$} before the row number fixes the row.
\end{definition}

The correct formula in \texttt{D2} is therefore \texttt{=C2*\$E\$1}. Copied down, it becomes \texttt{=C3*\$E\$1}, \texttt{=C4*\$E\$1}, \ldots{} The price adapts to each row; the VAT rate stays pinned to \texttt{E1}.

\begin{popfigure}
\begin{tikzpicture}[scale=0.92, every node/.style={font=\small}]
% grid
\foreach \x/\l in {0/A, 1.6/B, 3.2/C, 4.8/D, 6.4/E}{
  \node[fill=popblueL, minimum width=1.6cm, minimum height=0.55cm] at (\x,0.3) {\textbf{\l}};
}
\foreach \y/\r in {-0.35/1, -1.05/2, -1.75/3, -2.45/4}{
  \node[fill=popblueL, minimum width=0.55cm, minimum height=0.7cm] at (-0.85,\y) {\textbf{\r}};
}
\node[draw, minimum width=1.6cm, minimum height=0.7cm] at (6.4,-0.35) {19.25\%};
\node[draw, minimum width=1.6cm, minimum height=0.7cm] at (3.2,-1.05) {2500};
\node[draw, minimum width=1.6cm, minimum height=0.7cm] at (3.2,-1.75) {3000};
\node[draw, minimum width=1.6cm, minimum height=0.7cm] at (3.2,-2.45) {1800};
\node[draw, fill=popgreenL, minimum width=1.6cm, minimum height=0.7cm] (d2) at (4.8,-1.05) {\texttt{=C2*\$E\$1}};
\node[draw, fill=popgreenL, minimum width=1.6cm, minimum height=0.7cm] (d3) at (4.8,-1.75) {\texttt{=C3*\$E\$1}};
\node[draw, fill=popgreenL, minimum width=1.6cm, minimum height=0.7cm] (d4) at (4.8,-2.45) {\texttt{=C4*\$E\$1}};
% arrows
\draw[->, very thick, poporange] (d2.south) -- (d3.north) node[midway, right, font=\footnotesize\itshape, popdark] {copy};
\draw[->, very thick, poporange] (d3.south) -- (d4.north);
% annotations
\draw[->, thick, popblue] (3.2,-2.9) .. controls (3.6,-3.5) and (5.2,-3.5) .. (5.6,-2.9);
\node[font=\footnotesize, text=popblue, align=center] at (4.4,-3.75) {\texttt{C2} $\to$ \texttt{C3} $\to$ \texttt{C4} : relative, changes};
\draw[->, thick, poppurple] (6.4,-0.75) .. controls (7.9,-1.2) and (7.9,-2.2) .. (6.2,-2.5);
\node[font=\footnotesize, text=poppurple, align=center, text width=3.4cm] at (8.6,-1.7) {\texttt{\$E\$1} : absolute,\\ stays fixed};
\end{tikzpicture}
\legende{Copying \texttt{=C2*\$E\$1} down a column: the relative reference \texttt{C2} follows the row, while the absolute reference \texttt{\$E\$1} (the VAT rate) remains fixed.}
\end{popfigure}

\begin{attention}[The classic exam trap: forgetting the \$]
The most frequent mistake at the GCE is writing \texttt{=C2*E1} (relative) for a fixed rate and copying it down. Every row below the first then refers to an empty or wrong cell, giving zeros or absurd values. \textbf{Rule of thumb:} any value stored in a \emph{single} cell and used by \emph{many} formulae (a tax rate, an exchange rate, a coefficient, a pass mark) must be referenced \emph{absolutely} with \texttt{\$}.
\end{attention}

\begin{methode}[Converting a reference with the F4 key]
To change the type of a reference without typing the \texttt{\$} signs by hand:
\begin{enumerate}
\item Double-click the cell containing the formula (or click in the \textbf{formula bar}).
\item Click on the reference you want to change (or place the cursor inside it).
\item Press the \textbf{F4} key repeatedly to cycle through the four forms:
\[ \texttt{E1} \;\xrightarrow{\;F4\;}\; \texttt{\$E\$1} \;\xrightarrow{\;F4\;}\; \texttt{E\$1} \;\xrightarrow{\;F4\;}\; \texttt{\$E1} \;\xrightarrow{\;F4\;}\; \texttt{E1} \]
\item Press \textbf{Enter} to validate.
\end{enumerate}
(On some laptops, hold \textbf{Fn} while pressing F4.)
\end{methode}

\courssub{Mixed Referencing}

Sometimes we want to fix \emph{only} the column or \emph{only} the row. This gives the two \cle{mixed references}:

\begin{itemize}
\item \texttt{\$B1} : the \textbf{column B is fixed}, the row is free to change;
\item \texttt{B\$1} : the \textbf{row 1 is fixed}, the column is free to change.
\end{itemize}

The classic application is a \textbf{multiplication table}: numbers 1 to 10 down column \texttt{A} (rows 2 to 11) and across row 1 (columns \texttt{B} to \texttt{K}). The single formula typed in \texttt{B2} is:

\[ \texttt{=\$A2 * B\$1} \]

When copied \emph{down}, \texttt{\$A2} keeps pointing to column \texttt{A} (its row changes), and \texttt{B\$1} keeps pointing to row 1. When copied \emph{across}, \texttt{\$A2} stays on column \texttt{A} while \texttt{B\$1} moves along the header row. One formula fills the whole $10 \times 10$ grid. The same technique builds a price grid (unit price $\times$ quantity) or a fees table (base fee $\times$ number of terms).

\begin{popfigure}
\begin{tikzpicture}[node distance=0.4cm and 0.9cm, every node/.style={font=\small}]
\node[draw, fill=popblueL, rounded corners, text width=4.6cm, align=center] (q) {\textbf{Will this cell be copied to other cells?}};
\node[draw, fill=popgreenL, rounded corners, text width=4.2cm, align=center, below left=0.9cm and -1.2cm of q] (single) {Single calculation:\\ any reference type works};
\node[draw, fill=popblueL, rounded corners, text width=4.8cm, align=center, below right=0.9cm and -1.2cm of q] (q2) {\textbf{Should the reference stay on the same cell when copied?}};
\node[draw, fill=popgreenL, rounded corners, text width=3.6cm, align=center, below left=0.9cm and -0.6cm of q2] (rel) {\textbf{Relative} (\texttt{B2}):\\ it follows the data};
\node[draw, fill=popblueL, rounded corners, text width=4.4cm, align=center, below right=0.9cm and -0.6cm of q2] (q3) {\textbf{Fix column, row, or both?}};
\node[draw, fill=popgreenL, rounded corners, text width=3.4cm, align=center, below left=0.9cm and -0.5cm of q3] (abs) {\textbf{Absolute} (\texttt{\$B\$1}):\\ fixed rate, one cell};
\node[draw, fill=popgreenL, rounded corners, text width=3.8cm, align=center, below right=0.9cm and -0.5cm of q3] (mix) {\textbf{Mixed} (\texttt{\$B1} or \texttt{B\$1}):\\ grids and tables};
\draw[->, thick, popdark] (q) -- (single) node[midway, above left, font=\footnotesize] {No};
\draw[->, thick, popdark] (q) -- (q2) node[midway, above right, font=\footnotesize] {Yes};
\draw[->, thick, popdark] (q2) -- (rel) node[midway, above left, font=\footnotesize] {No};
\draw[->, thick, popdark] (q2) -- (q3) node[midway, above right, font=\footnotesize] {Yes};
\draw[->, thick, popdark] (q3) -- (abs) node[midway, above left, font=\footnotesize] {Both};
\draw[->, thick, popdark] (q3) -- (mix) node[midway, above right, font=\footnotesize] {One only};
\end{tikzpicture}
\legende{Decision guide for choosing between relative, absolute and mixed referencing.}
\end{popfigure}

\courssub{Basic Functions: SUM, AVERAGE, MAX, MIN, COUNT, COUNTA}

Adding twenty cells with \texttt{=B2+B3+B4+...} would be tedious. Spreadsheets provide built-in \cle{functions}: named procedures applied to a \emph{range} of cells. The values given to a function between its parentheses are called its \cle{arguments}.

\begin{center}
\begin{tabularx}{\linewidth}{|l|X|l|}
\hline
\rowcolor{popblueL} \textbf{Function} & \textbf{What it computes} & \textbf{Example} \\
\hline
\texttt{SUM(range)} & Total of all numbers in the range & \texttt{=SUM(B2:B20)} \\
\hline
\texttt{AVERAGE(range)} & Arithmetic mean of the numbers & \texttt{=AVERAGE(B2:B20)} \\
\hline
\texttt{MAX(range)} & Largest value in the range & \texttt{=MAX(B2:B20)} \\
\hline
\texttt{MIN(range)} & Smallest value in the range & \texttt{=MIN(B2:B20)} \\
\hline
\texttt{COUNT(range)} & Number of cells containing \emph{numbers} & \texttt{=COUNT(B2:B20)} \\
\hline
\texttt{COUNTA(range)} & Number of \emph{non-empty} cells (text included) & \texttt{=COUNTA(A2:A20)} \\
\hline
\end{tabularx}
\end{center}

\begin{attention}[COUNT versus COUNTA]
\texttt{COUNT} ignores text and empty cells; \texttt{COUNTA} counts everything that is not empty, including names. To count how many students \emph{have a mark}, use \texttt{COUNT} on the marks column; to count how many students are \emph{registered}, use \texttt{COUNTA} on the names column. Confusing the two is a classic source of wrong class sizes.
\end{attention}

\textbf{AutoSum.} The \emph{AutoSum} button ($\Sigma$, on the Home tab) inserts \texttt{=SUM(...)} automatically, guessing the range of numbers directly above or to the left of the active cell. Always \emph{check the proposed range} before pressing Enter --- AutoSum stops at the first empty cell, so a gap in the data silently truncates the total. The drop-down arrow next to $\Sigma$ also gives quick access to AVERAGE, MAX, MIN and COUNT.

\textbf{Status bar quick calculations.} When you select a range of cells containing numbers, the \emph{status bar} at the bottom of the window instantly displays the \textbf{Average}, \textbf{Count} and \textbf{Sum} of the selection, without typing any formula. This is ideal for a quick check, but note that these values are not stored anywhere in the sheet.

\textbf{A note on separators.} Depending on the regional settings of the computer, the arguments of a function are separated either by a comma \texttt{,} (English settings) or by a semicolon \texttt{;} (French settings): \texttt{=SUM(B2,C2)} or \texttt{=SUM(B2;C2)}. In the examination, follow the convention shown on the machine in front of you; both are accepted.

\courssub{Common Error Values}

When a formula cannot be evaluated, the spreadsheet displays an error value beginning with \texttt{\#}. Reading these messages correctly is a diagnostic skill tested at the GCE.

\begin{center}
\begin{tabularx}{\linewidth}{|l|X|X|}
\hline
\rowcolor{popblueL} \textbf{Error} & \textbf{Cause} & \textbf{Remedy} \\
\hline
\texttt{\#DIV/0!} & Division by zero or by an empty cell, e.g.\ \texttt{=A2/B2} with \texttt{B2} empty & Check the denominator; test for zero before dividing \\
\hline
\texttt{\#VALUE!} & Wrong type of operand, e.g.\ multiplying a number by text & Ensure all referenced cells contain numbers \\
\hline
\texttt{\#NAME?} & Unrecognised name, usually a \emph{misspelt function} (\texttt{=SOMME} in an English version, \texttt{=AVRAGE}) & Correct the spelling of the function \\
\hline
\texttt{\#REF!} & Invalid reference: a cell used by the formula was \emph{deleted} (row or column removed) & Rebuild the formula; avoid deleting referenced cells \\
\hline
\end{tabularx}
\end{center}

\courssub{Practical Build: A Class Mark Sheet}

We now assemble everything into a realistic GCE-style task: the mark sheet of an Upper Sixth class, with marks out of 20 in three subjects.

\begin{exoresolu}[Building the mark sheet step by step]
\textbf{Statement.} Columns: \texttt{A} = Name, \texttt{B} = ICT, \texttt{C} = Maths, \texttt{D} = Physics (rows 2 to 6 for five students). Cell \texttt{G1} contains the number of subjects, $3$. In column \texttt{E} compute each student's \textbf{total}, in column \texttt{F} the \textbf{average}, and below the table the \textbf{class average}, \textbf{highest} and \textbf{lowest} averages.
\tcblower
\textbf{Solution.}
\begin{enumerate}
\item \textbf{Totals (relative references).} In \texttt{E2} type \texttt{=SUM(B2:D2)} and copy down to \texttt{E6}. Each row adjusts automatically: \texttt{E3} contains \texttt{=SUM(B3:D3)}, etc.
\item \textbf{Averages (absolute reference).} In \texttt{F2} type \texttt{=E2/\$G\$1} and copy down. The total \texttt{E2} is relative (it follows the row); the divisor \texttt{\$G\$1} is absolute because the number of subjects lives in one single cell. Equivalently, \texttt{=AVERAGE(B2:D2)} works here, but the absolute-reference version is the one examiners look for.
\item \textbf{Class statistics.} In \texttt{F8}: \texttt{=AVERAGE(F2:F6)}; in \texttt{F9}: \texttt{=MAX(F2:F6)}; in \texttt{F10}: \texttt{=MIN(F2:F6)}; in \texttt{F11}: \texttt{=COUNT(F2:F6)} to verify that five averages exist.
\item \textbf{Rank (preview).} In \texttt{G2}: \texttt{=RANK(F2,\$F\$2:\$F\$6)} and copy down. Note again the mixed need: the student's own average \texttt{F2} is relative, but the \emph{list} against which we rank must be absolute, otherwise it would slide down during the copy. (Ranking functions are studied in detail in the next section.)
\item \textbf{Test the recalculation.} Change one ICT mark: the total, the average, the class statistics and the ranks all update instantly --- the automatic recalculation property in action.
\end{enumerate}
\end{exoresolu}

\begin{exercice}[Fees at a bilingual college]
A college stores in \texttt{B1} the termly fee per student (\SI{25000}{FCFA}) and in column \texttt{A} (rows 2 to 8) the number of children per family. (a) Write the formula in \texttt{B2} computing each family's termly bill, designed to be copied down. (b) Write the formula in \texttt{B10} giving the total collected. (c) The fee rises to \SI{27000}{FCFA}: what single action updates every result?
\end{exercice}

\begin{corrige}[Exercise 1]
(a) \texttt{=A2*\$B\$1} --- the family size is relative, the fee cell absolute. (b) \texttt{=SUM(B2:B8)}. (c) Simply type the new value $27000$ in \texttt{B1}: every formula recalculates automatically, which is precisely why the fee was isolated in its own cell and referenced absolutely.
\end{corrige}

\begin{exercice}[Choosing the reference type]
Cell \texttt{D1} holds an exchange rate ($1$ dollar $= 600$ FCFA). Column \texttt{B} (rows 2 to 9) holds amounts in dollars. (a) Write the formula in \texttt{C2}, to be copied down, converting each amount to FCFA. (b) Predict, without a computer, what formula appears in \texttt{C5} after the copy. (c) What error would appear in \texttt{C3} if the rate had been referenced relatively as \texttt{D1} and the cell \texttt{D2} were empty?
\end{exercice}

\begin{corrige}[Exercise 2]
(a) \texttt{=B2*\$D\$1}. (b) \texttt{=B5*\$D\$1}: the relative part moved from row 2 to row 5, the absolute part stayed fixed. (c) With \texttt{=B2*D1} copied down, \texttt{C3} would contain \texttt{=B3*D2}; since \texttt{D2} is empty it is treated as $0$, so \texttt{C3} would display $0$ --- a silent, dangerous wrong result rather than an error message.
\end{corrige}

\begin{aretenir}[Key points]
\begin{itemize}
\item Every formula begins with \texttt{=}; operators are \texttt{+ - * / \^\{\} \%}; priority follows BODMAS, parentheses first.
\item A \textbf{relative} reference (\texttt{B2}) moves when copied; an \textbf{absolute} reference (\texttt{\$B\$1}) never moves; a \textbf{mixed} reference (\texttt{\$B1}, \texttt{B\$1}) fixes only the column or only the row. Use \textbf{F4} to cycle through the forms.
\item Fixed values used by many formulae (rates, coefficients) go in their own cell and are referenced with \texttt{\$}.
\item Master \texttt{SUM}, \texttt{AVERAGE}, \texttt{MAX}, \texttt{MIN}, \texttt{COUNT}, \texttt{COUNTA}; AutoSum and the status bar speed up the work.
\item Read errors as diagnoses: \texttt{\#DIV/0!} (division by zero), \texttt{\#VALUE!} (wrong type), \texttt{\#NAME?} (misspelt function), \texttt{\#REF!} (deleted cell).
\item Formulae recalculate automatically whenever a referenced cell changes.
\end{itemize}
\end{aretenir}

With referencing and basic formulae mastered, we are ready for the next step: powerful logical and statistical functions, charts that make data speak, and workbooks in which several worksheets cooperate.

\coursec{Advanced Functions, Charts and Linked Worksheets}

In the previous section you learnt how to reference cells (relative, absolute and mixed references) and how to write basic formulae such as \texttt{=B2+C2} or \texttt{=SUM(B2:B10)}. Those formulae always do the same thing: they compute. But real-life school and business data often require the spreadsheet to \emph{make decisions}, to \emph{search} for information in a table, to \emph{summarise} only part of the data, and to \emph{present} results visually. That is exactly what this section is about: advanced functions, charts, and the linking of several worksheets into one coherent workbook.

\courssub{Logical Functions: IF, Nested IF, AND, OR}

\textbf{Intuition.} Imagine the examinations office of a school in Bamenda. A mark of 50 or more means ``Admitted'', otherwise ``Failed''. The spreadsheet must \emph{test a condition} and react accordingly. This is the job of the \cle{IF} function.

\begin{methode}[Writing an IF function]
The general syntax is:
\[ \texttt{=IF(condition, value\_if\_true, value\_if\_false)} \]
\begin{enumerate}
\item Write the \cle{condition} using comparison operators: \texttt{=}, \texttt{>}, \texttt{<}, \texttt{>=}, \texttt{<=}, \texttt{<>} (not equal).
\item Write the result to display when the condition is \textbf{TRUE}.
\item Write the result to display when the condition is \textbf{FALSE}.
\item Text results must be enclosed in double quotes, e.g.\ \texttt{"Admitted"}.
\end{enumerate}
\end{methode}

\begin{exemplebox}[Simple IF: admission decision]
Cell B2 contains a candidate's average. In C2 we write:
\[ \texttt{=IF(B2>=50, "Admitted", "Failed")} \]
If B2 $= 63$, C2 displays \emph{Admitted}; if B2 $= 41$, C2 displays \emph{Failed}.
\end{exemplebox}

\textbf{Nested IFs.} When there are more than two possible outcomes, we place one IF inside another. For a grading scale (A: 80 and above, B: 60--79, C: 50--59, F: below 50):
\[ \texttt{=IF(B2>=80,"A", IF(B2>=60,"B", IF(B2>=50,"C","F")))} \]
Read it from left to right: each FALSE branch opens a new test. Note that the number of closing brackets equals the number of IFs opened.

\textbf{AND and OR.} Sometimes the condition itself is compound. \texttt{AND(t1, t2, ...)} is TRUE only if \emph{all} tests are true; \texttt{OR(t1, t2, ...)} is TRUE if \emph{at least one} test is true. They are normally used inside IF:
\[ \texttt{=IF(AND(B2>=50, C2>=50), "Admitted", "Repeat")} \]
admits a student only when both the written mark (B2) and the practical mark (C2) reach 50, whereas
\[ \texttt{=IF(OR(B2<30, C2<30), "Weak", "OK")} \]
flags a student as weak as soon as one of the two marks falls below 30.

\begin{attention}[Common mistakes with logical functions]
\begin{itemize}
\item Forgetting the quotes around text: \texttt{=IF(B2>=50, Admitted, Failed)} produces an error.
\item Writing \texttt{80<=B2} instead of \texttt{B2>=80}: the cell reference must come first.
\item In nested IFs, testing in the wrong order (e.g.\ testing \texttt{>=50} before \texttt{>=80}) makes the higher grades unreachable, because the first true condition stops the evaluation.
\end{itemize}
\end{attention}

\courssub{Conditional Aggregation: SUMIF, COUNTIF, AVERAGEIF}

\textbf{Intuition.} The bursar wants the total fees paid \emph{by Form Five students only}, not by the whole school. Filtering by hand is slow and error-prone. Conditional aggregation functions do it in one formula.

\begin{definition}[Conditional aggregation functions]
\begin{itemize}
\item \texttt{=SUMIF(range, criteria, sum\_range)} adds the cells of \texttt{sum\_range} whose corresponding cells in \texttt{range} satisfy the \texttt{criteria}.
\item \texttt{=COUNTIF(range, criteria)} counts how many cells in \texttt{range} satisfy the criteria.
\item \texttt{=AVERAGEIF(range, criteria, average\_range)} averages only the matching values.
\end{itemize}
The criteria can be a value (\texttt{"Form Five"}), a comparison (\texttt{">=50"}) or a cell reference.
\end{definition}

\begin{exoresolu}[Fees paid by Form Five only]
Column A (A2:A30) holds the class of each student, column C (C2:C30) the fees paid. Write the formula giving (a) the total fees paid by Form Five students, (b) the number of Form Five students.
\tcblower
\begin{itemize}
\item[(a)] \texttt{=SUMIF(A2:A30, "Form Five", C2:C30)} — the test is applied to column A, but the amounts are taken from column C.
\item[(b)] \texttt{=COUNTIF(A2:A30, "Form Five")}.
\end{itemize}
If the class name were stored in cell E1, we would write \texttt{=SUMIF(A2:A30, E1, C2:C30)} so that changing E1 changes the class analysed.
\end{exoresolu}

\courssub{Lookup Functions: VLOOKUP and HLOOKUP}

\textbf{Intuition.} A shop in Douala keeps a price list: product code in column A, unit price in column B. On an invoice, the cashier types only the code, and the spreadsheet must \emph{fetch} the price automatically. This is a lookup.

\begin{definition}[VLOOKUP]
\cle{VLOOKUP} is a function that \textbf{searches a value in the first column of a range and returns a value from another column of the same row}. Its syntax is:
\[ \texttt{=VLOOKUP(lookup\_value, table\_range, col\_index, [range\_lookup])} \]
where \texttt{col\_index} is the number of the column (counted from the left of the table) to return, and \texttt{range\_lookup} is \texttt{FALSE} for an exact match or \texttt{TRUE} for an approximate match. \cle{HLOOKUP} works the same way but searches the first \emph{row} of a horizontal table.
\end{definition}

\begin{exemplebox}[Retrieving a unit price]
The price list occupies \texttt{\$A\$2:\$B\$20} (code, price). In the invoice, the code is typed in D2. The price is obtained by:
\[ \texttt{=VLOOKUP(D2, \$A\$2:\$B\$20, 2, FALSE)} \]
If D2 contains ``P07'' and the row of the list whose code is ``P07'' shows a price of 1500, the formula returns 1500.
\end{exemplebox}

\begin{attention}[The two classic VLOOKUP traps]
\begin{itemize}
\item \textbf{Forgetting to fix the table range with \$.} If you write \texttt{A2:B20} and copy the formula down, the range slides to \texttt{A3:B21}, \texttt{A4:B22}, \dots{} and lookups fail. Always use \texttt{\$A\$2:\$B\$20}.
\item \textbf{Omitting FALSE.} Without the fourth argument, VLOOKUP assumes an approximate match and may return the price of the \emph{wrong} product. For codes and names, always end with \texttt{FALSE}.
\end{itemize}
\end{attention}

\courssub{Text and Rounding Functions}

Spreadsheets also manipulate text. The most useful functions at this level are:

\begin{center}
\begin{tabularx}{\linewidth}{|l|X|l|}
\hline
\rowcolor{popblueL} \textbf{Function} & \textbf{Effect} & \textbf{Example result} \\
\hline
\texttt{=CONCATENATE(A2," ",B2)} or \texttt{=A2\&" "\&B2} & Joins several text strings into one & ``Amina Bello'' \\
\hline
\texttt{=LEFT(A2,3)} & First 3 characters from the left & ``UST'' from ``UST2024'' \\
\hline
\texttt{=RIGHT(A2,4)} & Last 4 characters from the right & ``2024'' from ``UST2024'' \\
\hline
\texttt{=UPPER(A2)} & Converts text to capital letters & ``AMINA'' \\
\hline
\texttt{=ROUND(B2,2)} & Rounds a number to 2 decimal places & 12.46 from 12.457 \\
\hline
\end{tabularx}
\end{center}

The \cle{\&} operator is a quick alternative to CONCATENATE. A typical use in a school is building a student ID: \texttt{=UPPER(LEFT(B2,3))\&"-"\&C2} produces e.g.\ ``BEL-045''.

\courssub{Charts: Types and Appropriate Use}

A table of numbers is precise but hard to read; a \cle{chart} reveals patterns at a glance. Choosing the \emph{right} type is an examinable skill.

\begin{center}
\begin{tabularx}{\linewidth}{|l|X|X|}
\hline
\rowcolor{popblueL} \textbf{Chart type} & \textbf{Best for} & \textbf{Example} \\
\hline
Column & Comparing values across categories & Average mark per subject \\
\hline
Bar & Same as column, with long category names & Enrolment per school \\
\hline
Line & Showing a trend over time & Fees collected per month \\
\hline
Pie & Showing parts of a whole (percentages) & Budget distribution \\
\hline
\end{tabularx}
\end{center}

\begin{aretenir}[GCE exam tip — choosing the chart]
Know which chart type suits which data: \textbf{pie} for parts of a whole, \textbf{line} for trends over time, \textbf{column} (or bar) for comparisons between categories. Never use a pie chart to show a trend, and never use a line chart for unrelated categories.
\end{aretenir}

\begin{popfigure}
\begin{center}
\begin{tikzpicture}
\begin{axis}[
  ybar, width=11cm, height=6.5cm, bar width=14pt,
  ymin=0, ymax=20,
  ylabel={Average mark (/20)},
  symbolic x coords={Maths, Physics, ICT, English, French},
  xtick=data,
  nodes near coords, nodes near coords align={vertical},
  ymajorgrids=true, grid style={popline},
]
\addplot[fill=popblue, draw=popdark] coordinates {(Maths,12.4) (Physics,10.8) (ICT,14.6) (English,11.9) (French,9.7)};
\end{axis}
\end{tikzpicture}
\end{center}
\legende{Column chart: average marks per subject for an Upper Sixth class. ICT is immediately visible as the strongest subject.}
\end{popfigure}

\begin{popfigure}
\begin{center}
\begin{tikzpicture}
% Pie chart: 55\% salaries, 30\% fees, 10\% maintenance, 5\% other
\fill[popblue]  (0,0) -- (90:2.4) arc (90:-108:2.4) -- cycle;
\fill[poporange](0,0) -- (-108:2.4) arc (-108:-216:2.4) -- cycle;
\fill[popgreen] (0,0) -- (-216:2.4) arc (-216:-252:2.4) -- cycle;
\fill[popgold]  (0,0) -- (-252:2.4) arc (-252:-270:2.4) -- cycle;
\draw[popdark, thick] (0,0) circle (2.4);
\node[font=\small\bfseries, white] at (-9:1.5) {55\%};
\node[font=\small\bfseries, white] at (198:1.5) {30\%};
\node[font=\small\bfseries, white] at (126:1.5) {10\%};
\node[font=\small\bfseries] at (99:1.6) {5\%};
% Legend
\fill[popblue]   (3.1,1.5) rectangle (3.5,1.8); \node[right, font=\small] at (3.6,1.65) {Salaries};
\fill[poporange] (3.1,0.8) rectangle (3.5,1.1); \node[right, font=\small] at (3.6,0.95) {Fees and supplies};
\fill[popgreen]  (3.1,0.1) rectangle (3.5,0.4); \node[right, font=\small] at (3.6,0.25) {Maintenance};
\fill[popgold]   (3.1,-0.6) rectangle (3.5,-0.3); \node[right, font=\small] at (3.6,-0.45) {Other};
\end{tikzpicture}
\end{center}
\legende{Pie chart: distribution of a school budget. Each slice is a part of the whole (100\%).}
\end{popfigure}

\textbf{Creating and customising a chart.} Select the data (including headings), then \emph{Insert $\rightarrow$ Chart} and choose the type. A good chart is then customised with:
\begin{itemize}
\item a \cle{chart title} stating what is shown (``Average marks per subject'');
\item \cle{axis titles} with units (``Mark /20'');
\item a \cle{legend} when several data series are plotted;
\item \cle{data labels} displaying the exact value on each column or slice;
\item a \cle{chart style} (colours, gridlines) that keeps the chart readable.
\end{itemize}

\begin{propriete}[Charts are linked to their data]
A chart created from a range is \emph{linked} to that range: when the source values change, the chart \textbf{updates automatically}. There is no need to redraw it. (This property itself is not proved; it is verified practically and often asked as a short-answer question.)
\end{propriete}

\begin{experience}[Verifying the automatic update]
Create a column chart from the marks of five students. Then change one mark in the table (e.g.\ from 10 to 18). Observe that the corresponding column grows immediately. This confirms that the chart is a \emph{view} of the data, not a fixed picture.
\end{experience}

\courssub{Linking Worksheets and Consolidating Data}

\textbf{Intuition.} A school keeps one sheet per term: \texttt{Term1}, \texttt{Term2}, \texttt{Term3}. The annual summary must combine them \emph{without retyping}. The solution is the \cle{cross-sheet reference}.

\begin{definition}[Cross-sheet and cross-workbook references]
To use a cell of another worksheet in a formula, prefix the address with the sheet name and an exclamation mark:
\[ \texttt{=Sheet2!B5} \qquad \texttt{=Term1!C10} \]
If the sheet name contains spaces, enclose it in single quotes: \texttt{='Term 1'!C10}. To reference another \emph{workbook}, add the file name in square brackets: \texttt{=[Fees.xlsx]Sheet1!B5}.
\end{definition}

\textbf{Consolidation.} The annual total of a student whose mark is in cell C10 of each term sheet is:
\[ \texttt{=Term1!C10 + Term2!C10 + Term3!C10} \quad\text{or}\quad \texttt{=SUM(Term1:Term3!C10)} \]
The second form, a \cle{3-D reference}, adds the same cell across a range of consecutive sheets — very efficient for consolidation.

\begin{popfigure}
\begin{center}
\begin{tikzpicture}[
  sheet/.style={draw=popdark, fill=popblueL, rounded corners, minimum width=2.9cm, minimum height=1.9cm, align=center, font=\small},
  arr/.style={-{Stealth[length=3mm]}, thick, poporange}]
\node[sheet, align=center] (t1) at (0,0){\textbf{Term1}\\ marks C2:C31};
\node[sheet, align=center] (t2) at (0,-2.6){\textbf{Term2}\\ marks C2:C31};
\node[sheet, align=center] (sm) at (6.2,-1.3){\textbf{Summary}\\ \texttt{=Term1!C10+Term2!C10}};
\draw[arr] (t1.east) -- (sm.170);
\draw[arr] (t2.east) -- (sm.190);
\end{tikzpicture}
\end{center}
\legende{Cross-sheet references: the Summary sheet pulls values from Term1 and Term2. Changing a mark in Term1 updates the Summary automatically.}
\end{popfigure}

\textbf{Maintaining links.} Links to other workbooks store the \emph{path} of the source file. When the destination file opens, the spreadsheet offers to \emph{update links} so that recent changes in the source are reflected.

\begin{attention}[Broken links]
If the source workbook is \textbf{moved, renamed or deleted}, the link breaks: the formula keeps the last known value or returns an error such as \texttt{\#REF!}. To avoid this: keep linked files in the same folder, move them together, and use \emph{Edit Links} (or the equivalent command) to repair paths when necessary.
\end{attention}

\begin{exoresolu}[Consolidating term totals]
Sheet \texttt{Term1} has total fees collected in B20, sheet \texttt{Term2} in B20, sheet \texttt{Term3} in B20. In the \texttt{Annual} sheet, write (a) the grand total, (b) the average per term.
\tcblower
\begin{itemize}
\item[(a)] \texttt{=Term1!B20 + Term2!B20 + Term3!B20}, or more compactly \texttt{=SUM(Term1:Term3!B20)}.
\item[(b)] \texttt{=AVERAGE(Term1:Term3!B20)}.
\end{itemize}
If a fee figure is corrected in Term2, both results update automatically because the Annual sheet is linked to the term sheets.
\end{exoresolu}

\begin{savaistu}[Spreadsheets in Cameroonian administration]
Schools, councils and businesses across Cameroon use linked workbooks every day: a bursary consolidates fee payments per class into one annual workbook, and mobile-money agents (MTN MoMo, Orange Money) export daily transaction sheets that are summarised monthly with SUMIF and charts. Mastering cross-sheet links is therefore a directly employable skill.
\end{savaistu}

\begin{exercice}[Practice]
A workbook contains a sheet \texttt{Products} (A: code, B: unit price) and a sheet \texttt{Sales} (A: code sold, B: quantity).
\begin{enumerate}
\item In Sales!C2, write the VLOOKUP formula that fetches the unit price of the code in A2.
\item In Sales!D2, write the formula for the line total, then in D20 the grand total.
\item In a sheet \texttt{Report}, write the formula counting how many sales lines have a quantity greater than 10.
\item State which chart type you would use to show the share of each product in the grand total, and justify.
\end{enumerate}
\end{exercice}

\begin{corrige}[Exercise 1]
\begin{enumerate}
\item \texttt{=VLOOKUP(A2, Products!\$A\$2:\$B\$50, 2, FALSE)} — fixed range and exact match.
\item \texttt{=B2*C2} copied down; grand total \texttt{=SUM(D2:D19)}.
\item \texttt{=COUNTIF(Sales!B2:B19, ">10")}.
\item A \textbf{pie chart}, because each product's sales are parts of a whole (the grand total).
\end{enumerate}
\end{corrige}

\coursec{What-If Analysis and Integration Activities}

In the previous section we learnt to build powerful formulas, visualise results with charts and link several worksheets together. Our spreadsheet models are now complete: give them input values and they instantly compute results. But real life often asks the \emph{reverse} question. A student does not ask ``what is my average?'' but rather ``\cle{what mark do I need} in Sequence 5 to reach an average of 10/20?'' A bursar does not ask ``what is the profit?'' but ``\cle{what happens to the profit} if the price of exercise books rises to 350 FCFA?'' Answering such questions is the purpose of \cle{what-if analysis}, the last great family of spreadsheet tools. We shall then close the chapter with \cle{integration activities} that combine Word and Excel exactly as required in the GCE practical examination.

\courssub{The Concept of What-If Analysis}

\begin{definition}[What-if analysis]
\cle{What-if analysis} is the process of changing the \textbf{input values} of a spreadsheet model in order to observe the effect on the \textbf{results} computed by formulas, without destroying the model itself.
\end{definition}

The intuition is simple: a spreadsheet is a chain \emph{inputs} $\rightarrow$ \emph{formulas} $\rightarrow$ \emph{results}. What-if analysis lets us play with the first link and watch the last link move. Excel offers three main tools, each suited to a different kind of question:

\begin{center}
\rowcolors{1}{popblueL}{white}
\begin{tabularx}{\linewidth}{|l|X|X|}
\hline
\rowcolor{popblueL}\textbf{Tool} & \textbf{Question it answers} & \textbf{Example} \\
\hline
Goal Seek & What single input gives a chosen result? & Mark needed for average 10/20 \\
\hline
Scenario Manager & How do fixed sets of inputs compare? & Best / expected / worst budget \\
\hline
Data Table & How does the result vary over a range of inputs? & Profit for prices 300, 350, 400\ldots \\
\hline
\end{tabularx}
\end{center}

\courssub{Goal Seek: Working Backwards from a Target}

Consider a concrete situation. A student has sequence marks 8, 11, 9 and 12 out of 20, and the average is computed by a formula, say \texttt{=AVERAGE(B2:G2)} placed in cell \texttt{H2}, where \texttt{G2} is the (still empty) Sequence 5 mark. The student wants to know: \emph{what mark in Sequence 5 gives an average of exactly 10?} One could try 9, then 10, then 11\ldots\ by trial and error. \cle{Goal Seek} automates this search.

\begin{definition}[Goal Seek]
\cle{Goal Seek} is a what-if tool that works \textbf{backwards} from a desired result: you give it a \textbf{target value} for a formula cell, and it adjusts \textbf{one input cell} until the formula reaches that target.
\end{definition}

\begin{methode}[Using Goal Seek]
\begin{enumerate}
\item Click the formula cell whose result you want to control (e.g.\ the average \texttt{H2}).
\item Go to \textbf{Data} $\rightarrow$ \textbf{What-If Analysis} $\rightarrow$ \textbf{Goal Seek}.
\item Fill the three boxes: \textbf{Set cell} = the formula cell (\texttt{H2}); \textbf{To value} = the target (\texttt{10}); \textbf{By changing cell} = the input cell (\texttt{G2}, the Sequence 5 mark).
\item Click \textbf{OK}. Excel displays the solution; click \textbf{OK} again to keep it or \textbf{Cancel} to restore the old value.
\end{enumerate}
\end{methode}

\begin{popfigure}
\begin{tikzpicture}[node distance=0.9cm,
box/.style={draw=popblue, thick, rounded corners, fill=popblueL, text width=3.4cm, align=center, minimum height=1cm, font=\small},
arr/.style={-{Stealth[length=3mm]}, thick, popdark}]
\node[box, align=center] (target){\textbf{Target value}\\ average = 10/20};
\node[box, right=of target, align=center] (tool){\textbf{Goal Seek}\\ adjusts the input cell (Seq.\ 5 mark)};
\node[box, right=of tool, align=center] (sol){\textbf{Solution found}\\ Seq.\ 5 = 10};
\draw[arr] (target) -- (tool);
\draw[arr] (tool) -- (sol);
\draw[arr, poporange, dashed] (sol.south) to[bend left=25] node[below, font=\scriptsize, text=poporange]{if not exact, iterate again} (tool.south);
\end{tikzpicture}
\legende{Goal Seek logic: from the desired result back to the required input.}
\end{popfigure}

\begin{exemplebox}[Mark needed for the average]
With marks 8, 11, 9, 12 in \texttt{B2:E2}, an empty \texttt{G2} and \texttt{=AVERAGE(B2:G2)} in \texttt{H2}, Goal Seek with \emph{Set cell} \texttt{H2}, \emph{To value} \texttt{10}, \emph{By changing cell} \texttt{G2} returns \texttt{G2 = 10}. Indeed $(8+11+9+12+10)/5 = 50/5 = 10$.
\end{exemplebox}

\begin{attention}[Goal Seek traps]
The \emph{changing cell} must contain a \textbf{plain value} (or be empty), never a formula --- Goal Seek cannot overwrite a formula. Also, if no exact solution exists (e.g.\ asking for an average of 25/20), Goal Seek reports that it ``may not have found a solution''; always check the result with your own eyes.
\end{attention}

\courssub{Scenarios: Comparing Named Sets of Values}

A school bursar preparing the budget for the new academic year does not know in advance how many students will enrol or what the price of chalk will be. Rather than creating three separate files, he saves three \cle{scenarios} in one sheet.

\begin{definition}[Scenario]
A \cle{scenario} is a \textbf{named set of input values} stored inside a worksheet, which can be \textbf{switched on} at any moment to compare the resulting outcomes (for example \emph{best case}, \emph{expected case}, \emph{worst case}).
\end{definition}

\begin{methode}[Creating and viewing scenarios]
\begin{enumerate}
\item Build the model (e.g.\ \texttt{=B2*B3-B4} for profit from fees).
\item Go to \textbf{Data} $\rightarrow$ \textbf{What-If Analysis} $\rightarrow$ \textbf{Scenario Manager} $\rightarrow$ \textbf{Add}.
\item Name the scenario (\texttt{Worst case}), select the \textbf{changing cells}, type their values, then \textbf{OK}.
\item Repeat for \texttt{Expected case} and \texttt{Best case}.
\item Select a scenario and click \textbf{Show} to apply it; use \textbf{Summary} to generate a comparison report on a new sheet.
\end{enumerate}
\end{methode}

\begin{exemplebox}[School canteen budget]
Profit $=$ (meals sold $\times$ price) $-$ costs. Scenarios: \emph{Worst} (200 meals, 250 FCFA, costs 60\,000), \emph{Expected} (300 meals, 300 FCFA, costs 70\,000), \emph{Best} (400 meals, 300 FCFA, costs 75\,000). Switching scenarios instantly shows profits of $-10\,000$, $20\,000$ and $45\,000$ FCFA respectively --- the bursar sees at a glance that the canteen must sell at least about 240 meals to break even.
\end{exemplebox}

\courssub{Data Tables: Varying One or Two Inputs Systematically}

Goal Seek varies one input to hit \emph{one} target; scenarios compare a \emph{few} fixed cases. But sometimes we want a whole \emph{map} of results: profit for \emph{every} price from 250 to 500 FCFA, or for every combination of price \emph{and} quantity. This is the job of \cle{data tables}.

\begin{definition}[Data table]
A \cle{data table} is a range of cells in which Excel automatically re-evaluates a formula for a whole list of input values. A \textbf{one-variable} data table shows how results change with one input; a \textbf{two-variable} data table shows results for every pair of values of two inputs, arranged as a grid.
\end{definition}

\begin{methode}[Building a one-variable data table]
\begin{enumerate}
\item Build the model, e.g.\ profit \texttt{=B1*(B2-B3)} where \texttt{B1} is quantity, \texttt{B2} price, \texttt{B3} unit cost.
\item In a column, list the trial prices (250, 300, 350, 400); in the cell above and one column to the right, reference the result formula (e.g.\ \texttt{=B4} if the profit is in \texttt{B4}).
\item Select the whole block, then \textbf{Data} $\rightarrow$ \textbf{What-If Analysis} $\rightarrow$ \textbf{Data Table}; set the \textbf{Column input cell} to \texttt{B2} and click \textbf{OK}.
\end{enumerate}
For a \textbf{two-variable} table, place prices down the left column, quantities across the top row, the formula reference in the top-left corner cell, then fill both a \emph{Row input cell} and a \emph{Column input cell}.
\end{methode}

\begin{exemplebox}[Profit grid for the co-operative]
A two-variable data table with prices 300--500 FCFA (rows) and quantities 100--400 (columns) produces a grid of profits. The committee immediately sees which (price, quantity) pairs keep the co-operative profitable, without typing a single extra formula.
\end{exemplebox}

\begin{attention}[Choosing the right tool]
Exam questions often test this choice. Remember: \textbf{one input, one exact target} $\rightarrow$ Goal Seek; \textbf{a few named cases to compare} $\rightarrow$ Scenarios; \textbf{a whole range of values to explore} $\rightarrow$ Data table. Using Goal Seek to fill a table of 20 values is a classic waste of exam time.
\end{attention}

\courssub{Integration Activity 1: A Complete School Report}

The GCE practical paper rarely tests Word or Excel in isolation; it asks you to \emph{combine} them. The first integration activity is a school report: an Excel mark sheet (formulas, ranking, chart) presented inside a Word document.

\begin{methode}[Producing the combined report]
\begin{enumerate}
\item In Excel, build the mark sheet: totals with \texttt{SUM}, averages with \texttt{AVERAGE}, ranks with \texttt{RANK}, decisions with \texttt{IF}, and a bar chart of subject averages.
\item In Word, type the report text (heading with styles, introduction, comments).
\item Select the Excel chart, copy it (\texttt{Ctrl+C}), then in Word use \textbf{Home} $\rightarrow$ \textbf{Paste} $\rightarrow$ \textbf{Paste Special} $\rightarrow$ \textbf{Paste Link} $\rightarrow$ \emph{Microsoft Excel Chart Object}.
\item Save both files in the same folder. When the Excel marks are corrected, the chart in Word \textbf{updates automatically}.
\end{enumerate}
\end{methode}

\begin{popfigure}
\begin{tikzpicture}[node distance=1.6cm,
file/.style={draw=popdark, thick, rounded corners, text width=3.8cm, align=center, minimum height=2.2cm, font=\small},
arr/.style={-{Stealth[length=3mm]}, very thick, poporange}]
\node[file, fill=popgreenL, align=center] (excel){\textbf{Excel: marks.xlsx}\\[3pt]
\textcolor{popblue}{\rule{0.25cm}{0.7cm}}\ \textcolor{popblue}{\rule{0.25cm}{1.05cm}}\ \textcolor{popblue}{\rule{0.25cm}{0.5cm}}\ \textcolor{popblue}{\rule{0.25cm}{0.9cm}}\\[3pt]
chart of averages};
\node[file, fill=popblueL, right=of excel, align=center] (word){\textbf{Word: report.docx}\\[3pt]
\rule{2.6cm}{0.12cm}\\[2pt]
\rule{2.6cm}{0.12cm}\\[2pt]
\textcolor{popblue}{\rule{0.25cm}{0.7cm}}\ \textcolor{popblue}{\rule{0.25cm}{1.05cm}}\ \textcolor{popblue}{\rule{0.25cm}{0.5cm}}\ \textcolor{popblue}{\rule{0.25cm}{0.9cm}}\\[3pt]
linked chart};
\draw[arr] (excel) -- node[above, font=\small\itshape, text=poporange]{updates automatically} (word);
\end{tikzpicture}
\legende{A chart linked from Excel into a Word report: editing the workbook refreshes the document.}
\end{popfigure}

\courssub{Object Linking and Embedding (OLE)}

\emph{Paste Special} offers three behaviours that the GCE loves to ask about. Together they form the technology called \cle{OLE} (\textbf{O}bject \textbf{L}inking and \textbf{E}mbedding).

\begin{center}
\rowcolors{1}{popblueL}{white}
\begin{tabularx}{\linewidth}{|l|X|X|}
\hline
\rowcolor{popblueL}\textbf{Mode} & \textbf{Behaviour} & \textbf{When to use it} \\
\hline
Simple paste (picture/table) & A static copy; no connection with Excel & Final document, data will never change \\
\hline
Linking & Word stores a \emph{reference} to the Excel file; the chart updates when the workbook changes & Both files stay together and data evolves \\
\hline
Embedding & A full \emph{copy} of the Excel object is stored inside the Word file; double-clicking edits it with Excel tools & The Word file must travel alone (e-mail, USB key) \\
\hline
\end{tabularx}
\end{center}

\begin{attention}[The broken-link trap]
A very common mistake: a candidate links a chart into Word, then \textbf{deletes or moves the Excel file} (or copies only the \texttt{.docx} to a USB key). The link \textbf{breaks} and the chart can no longer update --- it may even disappear. Rule of thumb: if the Excel file will not travel with the document, \textbf{embed} the chart instead of linking it.
\end{attention}

\courssub{Integration Activity 2: Mail-Merging Result Slips}

Sending result slips to 200 parents does not mean typing 200 letters. \cle{Mail merge} combines one Word template with the Excel mark sheet used as a \textbf{data source}.

\begin{methode}[Mail merge of result slips]
\begin{enumerate}
\item In Excel, prepare a clean table: first row = field names (\texttt{Name}, \texttt{Class}, \texttt{Average}, \texttt{Rank}), one row per student. Save and close the file.
\item In Word: \textbf{Mailings} $\rightarrow$ \textbf{Start Mail Merge} $\rightarrow$ \textbf{Letters}.
\item \textbf{Select Recipients} $\rightarrow$ \textbf{Use an Existing List} $\rightarrow$ choose the Excel file and the correct worksheet.
\item Place the cursor in the letter and use \textbf{Insert Merge Field} to add \texttt{<<Name>>}, \texttt{<<Average>>}, \texttt{<<Rank>>}.
\item \textbf{Preview Results} to check a few records, then \textbf{Finish \& Merge} $\rightarrow$ \textbf{Print Documents} or \textbf{Edit Individual Documents}.
\end{enumerate}
\end{methode}

\begin{exoresolu}[Choosing and applying a what-if tool]
\textbf{Statement.} A candidate's model computes the term average in \texttt{H2} from five sequence marks. (a) Which tool finds the Sequence 5 mark needed for an average of exactly 10/20? (b) Which tool compares the budget outcomes for enrolments of 400, 500 and 600 students? (c) Which tool shows the profit for every ticket price from 100 to 1000 FCFA in steps of 100? Justify each answer.
\tcblower
\textbf{Solution.} (a) \textbf{Goal Seek}: a single input (one mark) must be adjusted to reach one exact target (10/20) --- set cell \texttt{H2}, to value \texttt{10}, by changing the Sequence 5 cell. (b) \textbf{Scenario Manager}: three \emph{named, fixed} sets of inputs (400/500/600 students) are stored and switched to compare outcomes. (c) A \textbf{one-variable data table}: the result must be evaluated over a whole \emph{range} of values of a single input (the price), which is exactly what a data table automates.
\end{exoresolu}

\courssub{Revision Synthesis and GCE Practice}

\begin{aretenir}[Chapter synthesis]
\begin{itemize}
\item Word processing: styles, headers/footers, tables of contents, mail merge produce professional documents.
\item Spreadsheets: formatting, sorting, filtering, absolute referencing (\texttt{\$B\$2}), functions (\texttt{SUM}, \texttt{AVERAGE}, \texttt{IF}, \texttt{VLOOKUP}, \texttt{RANK}), charts and linked sheets turn raw data into information.
\item What-if analysis: \textbf{Goal Seek} (one input, one target), \textbf{Scenarios} (named cases), \textbf{Data tables} (ranges of one or two inputs).
\item OLE: \textbf{paste} (static), \textbf{link} (updates, needs the source file), \textbf{embed} (portable, larger file).
\end{itemize}
\end{aretenir}

\begin{exercice}[GCE-style practice]
A school stores CA marks in \texttt{Sheet1} and exam marks in \texttt{Sheet2}. (a) Write the formula computing the term total on \texttt{Sheet1} that pulls the exam mark from \texttt{Sheet2}. (b) State the what-if tool and the exact dialog entries to find the exam mark giving a term average of 10/20. (c) The principal wants result slips for all parents: describe the mail-merge steps. (d) The report chart must remain correct even if the Excel file is deleted: should you link or embed? Justify.
\end{exercice}

\begin{corrige}[Exercise 1]
(a) \texttt{=B2+Sheet2!B2} (with appropriate cells). (b) Goal Seek: Set cell = average cell, To value = \texttt{10}, By changing cell = the exam mark cell on \texttt{Sheet2}. (c) Mailings $\rightarrow$ Start Mail Merge $\rightarrow$ Letters; Select Recipients $\rightarrow$ Use Existing List $\rightarrow$ the Excel sheet; Insert Merge Fields; Preview; Finish \& Merge. (d) \textbf{Embed}: linking would break when the source file is deleted, whereas embedding stores a full copy of the object inside the Word document.
\end{corrige}

\begin{savaistu}[GCE exam tip]
Practical papers very often require the \textbf{full chain}: a formula sheet in Excel, a chart, a linked or embedded object in a Word report, and a merged letter --- all within a limited time. Practise this complete workflow at home until you can do it in under 45 minutes; on exam day, save your files early and often, and never delete a source file before checking every linked object.
\end{savaistu}

\newpage

\coursec{Integration activity}

\courssub{Aims of the activity}

This integration activity closes the chapter by mobilising, in one realistic project, \emph{all} the competences studied in Lessons 58 to 63: word-processing features for producing documents, formatting and managing spreadsheet data, cell referencing and formulae, advanced functions and charts, linking worksheets, and what-if analysis. By the end of the activity, each pair of learners should be able to:

\begin{itemize}
\item structure a multi-sheet workbook (one sheet per sequence plus a summary sheet) and link sheets with \cle{cross-sheet references};
\item use \cle{relative}, \cle{absolute} and \cle{mixed references} correctly when copying formulae;
\item compute totals, averages and ranks, and apply logical and statistical functions (\texttt{IF}, \texttt{SUMIF}, \texttt{COUNTIF}, \texttt{AVERAGE}, \texttt{RANK}, \texttt{MAX}, \texttt{MIN});
\item build and format a \cle{bar chart} and a \cle{pie chart} from spreadsheet data;
\item use \cle{Goal Seek} to answer a what-if question;
\item produce a formal Word document with \cle{heading styles}, an \cle{automatic table of contents} and a linked chart;
\item perform a \cle{mail merge} using an Excel sheet as data source;
\item demonstrate that the package is \emph{automated}: changing one raw mark updates everything.
\end{itemize}

\begin{aretenir}[Success criteria]
The package is successful if (i) no result is typed twice — everything is computed by formulae or links; (ii) the charts and the Word report update when a mark changes; (iii) the mail merge produces 20 personalised result slips without retyping any name or mark.
\end{aretenir}

\courssub{Situation}

\begin{experience}[End-of-term reports at GHS Nkwen]
Government High School Nkwen (Bamenda, North-West Region) still prepares end-of-term reports by hand. The Principal, Mrs.\ A.\ Mbuh, has asked the Upper Sixth ICT class to digitalise the process for one pilot class, \textbf{Form 5A}, before generalising it to the whole school.

The class teacher gives you the raw marks (out of 20) of \textbf{20 students} in \textbf{6 subjects} for each of the two sequences of the term. The subjects and their coefficients are those of the Anglophone sub-system for Form 5: English Language (coef.\ 3), Mathematics (coef.\ 4), Physics (coef.\ 3), Chemistry (coef.\ 2), Biology (coef.\ 2), Computer Science (coef.\ 2). The \textbf{total coefficient is 16}; the weighted term average is computed out of 20. The school rule for the term decision is:

\begin{center}
\begin{tabularx}{\linewidth}{|l|X|}
\hline
\rowcolor{popblueL} \textbf{Weighted average} & \textbf{Decision} \\
\hline
Average $\geq 10/20$ & Admitted (passes to the next sequence in good standing) \\
\hline
Average $< 10/20$ & Repeat work / warned (parent summoned) \\
\hline
\end{tabularx}
\end{center}

Extract of the data for Sequence 1 (the full list of 20 students is provided by the teacher as a CSV file exported from the school register):

\begin{center}
\begin{tabular}{|l|c|c|c|c|c|c|}
\hline
\rowcolor{popblueL} \textbf{Name} & \textbf{ENG} & \textbf{MAT} & \textbf{PHY} & \textbf{CHE} & \textbf{BIO} & \textbf{CSC} \\
\hline
Abongwa Claris & 12 & 14 & 11 & 13 & 15 & 16 \\
\hline
Bih Terence & 9 & 7 & 8 & 10 & 11 & 14 \\
\hline
Che Mildred & 15 & 16 & 14 & 12 & 13 & 17 \\
\hline
Fru Collins & 6 & 5 & 7 & 8 & 9 & 12 \\
\hline
\multicolumn{7}{|c|}{\emph{\ldots 16 other students \ldots}} \\
\hline
\end{tabular}
\end{center}

Sequence 2 marks are given in a second CSV file with the same structure. One case worries the class teacher: \textbf{Bih Terence} has a Sequence 2 weighted average of $9.6/20$. The Principal wants to know: \emph{``What mark does Terence need in Sequence 3 Mathematics for his cumulative average to reach exactly $10/20$?''}

The Principal expects from each pair a complete package:
\begin{enumerate}
\item an \textbf{Excel workbook} \texttt{Form5A\_Term1.xlsx} with sheets \texttt{Seq1}, \texttt{Seq2}, \texttt{Summary} and \texttt{Statistics};
\item a \textbf{one-page Word report} to the Principal, with a table of contents and the chart of subject averages;
\item \textbf{20 individual result slips} for parents, produced by mail merge;
\item a \textbf{5-minute presentation}, ending with a live demonstration: the teacher changes one mark, and the whole package (averages, ranks, decision, charts) recalculates automatically.
\end{enumerate}
\end{experience}

\begin{popfigure}
\begin{tikzpicture}[
box/.style={draw=popblue, thick, rounded corners, fill=popblueL, text width=3.1cm, align=center, minimum height=1cm, font=\small},
arr/.style={->, thick, popdark}]
\node[box, align=center] (s1) at (0,2.2){\texttt{Seq1}\\ raw marks + formulae};
\node[box, align=center] (s2) at (0,0){\texttt{Seq2}\\ raw marks + formulae};
\node[box, align=center] (sum) at (4.6,1.1){\texttt{Summary}\\ links, rank, decision};
\node[box, align=center] (stat) at (9.2,2.2){\texttt{Statistics}\\ COUNTIF, charts};
\node[box] (word) at (9.2,0) {Word report + mail merge};
\draw[arr] (s1) -- (sum);
\draw[arr] (s2) -- (sum);
\draw[arr] (sum) -- (stat);
\draw[arr] (sum) -- (word);
\draw[arr] (stat) -- (word);
\end{tikzpicture}
\legende{Data flow of the package: raw marks are entered once; every other result is a formula or a link.}
\end{popfigure}

\courssub{Tasks}

\begin{exercice}[Task 1 — Building the workbook]
Import the two CSV files into sheets \texttt{Seq1} and \texttt{Seq2}. Add columns \textbf{Weighted total} and \textbf{Weighted average}. Write the formula for the weighted total of the first student using the coefficients stored in a header row with \emph{mixed references} (e.g.\ \texttt{=B2*B\$1+C2*C\$1+...}), then copy it down for the 20 students. Explain why the dollar sign is placed before the row number and not before the column letter.
\end{exercice}

\begin{exercice}[Task 2 — The linked summary sheet]
On the \texttt{Summary} sheet, retrieve each student's two sequence averages using \cle{cross-sheet references} (e.g.\ \texttt{=Seq1!I2}), compute the \textbf{term average} (mean of the two sequences), the \textbf{rank} with \texttt{RANK}, and the \textbf{decision} with \texttt{IF}. Why must the range of averages inside \texttt{RANK} be an absolute reference?
\end{exercice}

\begin{exercice}[Task 3 — Class statistics]
On the \texttt{Statistics} sheet, compute: the number of students (\texttt{COUNTA}), the class average (\texttt{AVERAGE}), the highest and lowest averages (\texttt{MAX}, \texttt{MIN}), the number of admitted students (\texttt{COUNTIF}) and the total of averages above 10 (\texttt{SUMIF}). Compute the \textbf{success rate} as a percentage.
\end{exercice}

\begin{exercice}[Task 4 — Charts]
Create (a) a \textbf{bar chart} of the six subject averages for the term, and (b) a \textbf{pie chart} showing the distribution Admitted / Repeat work. Give each chart a title, axis titles (bar chart) and data labels in percentage (pie chart).
\end{exercice}

\begin{exercice}[Task 5 — Goal Seek]
Using \cle{Goal Seek} (Data $\rightarrow$ What-If Analysis), determine the mark Terence needs in Sequence 3 Mathematics for his cumulative weighted average to equal exactly $10/20$. Set up the calculation so that the cumulative average is a formula depending on the unknown mark.
\end{exercice}

\begin{exercice}[Task 6 — Word report and mail merge]
(a) Write the one-page report to the Principal using \textbf{Heading 1 / Heading 2} styles, insert an \textbf{automatic table of contents}, and paste the bar chart as a \emph{linked} object. (b) Prepare the result slip (student name, subject averages, term average, rank, decision) and run a \textbf{mail merge} with the \texttt{Summary} sheet as data source to produce the 20 slips.
\end{exercice}

\begin{exercice}[Task 7 — Live demonstration]
Present the package in 5 minutes. Change Fru Collins' Sequence 2 Mathematics mark from 5 to 12 and show that the summary, statistics, charts and decision update without any retyping.
\end{exercice}

\courssub{Model answer}

\textbf{Task 1.} After import, row 1 holds the coefficients in \texttt{B1:G1} (3, 4, 3, 2, 2, 2) and \texttt{H1} contains the total coefficient 16. In \texttt{H2}:

\[ \texttt{=B2*B\$1+C2*C\$1+D2*D\$1+E2*E\$1+F2*F\$1+G2*G\$1} \]

and in \texttt{I2}: \texttt{=H2/\$H\$1}. The reference \texttt{B\$1} is \emph{mixed}: the column letter stays free so that the formula reads each subject's mark and coefficient along the row, while the row is locked with \texttt{\$1} so that copying \emph{down} to rows 3--21 still points to the coefficient row. For Abongwa Claris:
\begin{align*}
\text{Total} &= 12\times3+14\times4+11\times3+13\times2+15\times2+16\times2 \\
&= 36+56+33+26+30+32 = 213,
\end{align*}
giving a weighted average of $213/16 = 13.31/20$ (rounded to two decimal places).

\textbf{Task 2.} On \texttt{Summary}, \texttt{B2} = \texttt{=Seq1!I2}, \texttt{C2} = \texttt{=Seq2!I2}, \texttt{D2} = \texttt{=AVERAGE(B2:C2)}. Rank in \texttt{E2}: \texttt{=RANK(D2,\$D\$2:\$D\$21)} — the range must be \emph{absolute} (\texttt{\$D\$2:\$D\$21}) so that it does not shift when the formula is copied down; otherwise student 3 would be ranked against \texttt{D3:D22}, which is wrong. Decision in \texttt{F2}: \texttt{=IF(D2>=10,"Admitted","Repeat work")}.

\textbf{Task 3.} Typical formulae on \texttt{Statistics}:

\begin{center}
\begin{tabularx}{\linewidth}{|l|X|}
\hline
\rowcolor{popblueL} \textbf{Statistic} & \textbf{Formula} \\
\hline
Number of students & \texttt{=COUNTA(Summary!A2:A21)} \\
\hline
Class average & \texttt{=AVERAGE(Summary!D2:D21)} \\
\hline
Highest / lowest average & \texttt{=MAX(Summary!D2:D21)} / \texttt{=MIN(Summary!D2:D21)} \\
\hline
Number admitted & \texttt{=COUNTIF(Summary!F2:F21,"Admitted")} \\
\hline
Sum of averages $\geq 10$ & \texttt{=SUMIF(Summary!D2:D21,">=10")} \\
\hline
Success rate & \texttt{=COUNTIF(Summary!F2:F21,"Admitted")/20} formatted as \% \\
\hline
\end{tabularx}
\end{center}

\textbf{Task 4.} The subject averages are computed with \texttt{=AVERAGE(Seq1!B2:B21)} etc.\ on the \texttt{Statistics} sheet, then selected and inserted as a \textbf{clustered bar chart} titled ``Form 5A — Subject averages, Term 1'', with axis titles ``Subject'' and ``Average (/20)''. The pie chart uses the two \texttt{COUNTIF} results (Admitted / Repeat work) with percentage data labels.

\textbf{Task 5.} First build the model. In a free area, enter Terence's estimated Sequence 3 marks for the five other subjects, and a \emph{trial mark} (e.g.\ 10) for Mathematics in a single cell, say \texttt{K2}. Compute his Sequence 3 weighted total with the same coefficient formula (Mathematics contributes $4\times\texttt{K2}$), then his Sequence 3 average, then the cumulative average \texttt{=(Seq1avg+Seq2avg+Seq3avg)/3}. Open \textbf{Data $\rightarrow$ What-If Analysis $\rightarrow$ Goal Seek}: \emph{Set cell} = cumulative average cell, \emph{To value} = 10, \emph{By changing cell} = \texttt{K2}. Goal Seek iterates and returns the required Mathematics mark.

\emph{Numerical illustration.} Suppose Terence's Sequence 1 average is $9.0/20$ and his estimated Sequence 3 marks (without Mathematics) give a partial weighted total of 130. To reach a cumulative average of 10 he needs a Sequence 3 average of $3\times10-9.0-9.6 = 11.4$, i.e.\ a weighted total of $11.4\times16 = 182.4$, hence a Mathematics mark of $(182.4-130)/4 = 13.1/20$. Your own result will depend on his actual Sequence 1 average and estimated marks — that is precisely why Goal Seek is used rather than mental arithmetic. If Goal Seek returns a mark above 20, the honest conclusion to report is that the target of $10/20$ is unreachable through Mathematics alone.

\textbf{Task 6.} The Word report uses \textbf{Heading 1} for ``Form 5A — End-of-Term Report'' and \textbf{Heading 2} for ``Results overview'', ``Statistics'' and ``Recommendations''; the table of contents is inserted via \textbf{References $\rightarrow$ Table of Contents} and refreshed with F9. The chart is pasted with \textbf{Paste Special $\rightarrow$ Paste link} so it updates with the workbook. For the slips: \textbf{Mailings $\rightarrow$ Start Mail Merge $\rightarrow$ Letters}, \textbf{Select Recipients $\rightarrow$ Use an Existing List} pointing to \texttt{Form5A\_Term1.xlsx}, sheet \texttt{Summary}; merge fields \texttt{Name}, \texttt{Average}, \texttt{Rank}, \texttt{Decision} are inserted, then \textbf{Finish \& Merge} generates the 20 personalised slips.

\textbf{Task 7.} Changing Fru Collins' mark in \texttt{Seq2} propagates instantly: his average, rank and decision change, the \texttt{COUNTIF} statistics and success rate update, and both charts are redrawn — proving the package is fully automated.

\begin{attention}[Common mistakes to avoid]
\begin{itemize}
\item Typing computed values by hand instead of formulae — the live demonstration then fails immediately.
\item Forgetting the \texttt{\$} in the \texttt{RANK} range or in the coefficient row: copied formulae then give wrong results.
\item Averaging percentages or ranks instead of raw weighted totals.
\item Dividing by the number of subjects (6) instead of the total coefficient (16) when computing a weighted average.
\item Pasting the chart as a simple picture instead of a \emph{linked} object: the Word report no longer updates.
\end{itemize}
\end{attention}

\begin{savaistu}[Why this matters in Cameroon]
Schools that digitalise their reporting save days of manual calculation each term and reduce transcription errors. The same skills — linked sheets, mail merge, Goal Seek — are used daily in Cameroonian banks, mobile-money agent accounting (MTN MoMo, Orange Money) and the e-government services promoted by ANTIC. Mastering them makes you immediately useful in any administration.
\end{savaistu}

\newpage

\coursec{Methods and worked examples}

\coursec{Using Formatting, Editing Features of Word Processors for Producing Documents}

\begin{savaistu}[Why this chapter matters for the GCE]
The GCE Advanced Level ICT practical paper (Paper 3) and the theory paper (Paper 1) both test your ability to \emph{produce professional documents} and \emph{analyse data in spreadsheets}. This chapter chains the seven lessons of the official Second Term progression (Lessons 58--64) into a single workflow: from a blank page to a structured report, from raw data to a chart, from a formula to a what-if scenario, and finally to an integrated document that could be submitted to a principal, a bank manager, or an ANTIC inspector. Master every method below; they are the exact skills the examiner expects.
\end{savaistu}

\courssub{Producing Documents with Word Processing Features}

\begin{experience}[Observing a poorly formatted document]
Open the file \texttt{BadReport.docx} (provided in your resources). Notice: the title is centred by pressing \texttt{Space} repeatedly; headings are just bold, large text; the table of contents was typed by hand; page numbers are missing; the table columns do not align. Try to insert a new section in the middle --- everything shifts. This is what \emph{not} to do. The professional approach uses \cle{styles}, \cle{page breaks}, and \cle{automatic fields}.
\end{experience}

\begin{methode}[Producing a professional document in a word processor]
To produce a well-structured document (report, letter, CV, minutes) with Microsoft Word or LibreOffice Writer:
\begin{enumerate}
\item \textbf{Plan the structure} on paper first: title page, table of contents, headings (level 1, level 2), body paragraphs, tables, images, references, appendices.
\item \textbf{Set up the page}: \emph{Layout $\rightarrow$ Size: A4}; \emph{Margins: Normal (2.54 cm)}; \emph{Orientation: Portrait} (or Landscape for wide tables).
\item \textbf{Apply styles, not manual formatting}: select each heading and apply \emph{Heading 1}, \emph{Heading 2}, etc.; apply \emph{Normal} to body text. This enables automatic \cle{table of contents}, \cle{navigation pane}, and consistent reformatting.
\item \textbf{Format text properly}: font (Calibri/Arial 11 or 12), alignment (justified for body, centred for title), line spacing (1.15 or 1.5), bullets/numbering via the \emph{Home} tab.
\item \textbf{Insert objects correctly}:
\begin{itemize}
\item Tables: \emph{Insert $\rightarrow$ Table} $\rightarrow$ choose rows/columns $\rightarrow$ type data $\rightarrow$ \emph{Table Design} for header row shading.
\item Pictures: \emph{Insert $\rightarrow$ Pictures} $\rightarrow$ wrap text \emph{In Line with Text} or \emph{Tight}.
\item Page numbers: \emph{Insert $\rightarrow$ Page Number $\rightarrow$ Bottom of Page $\rightarrow$ Plain Number 2}.
\item Headers/Footers: \emph{Insert $\rightarrow$ Header/Footer} $\rightarrow$ type name, date, document code.
\item Footnotes/Endnotes: \emph{References $\rightarrow$ Insert Footnote}.
\item Page breaks: \texttt{Ctrl+Enter} (never multiple \texttt{Enter} keys).
\end{itemize}
\item \textbf{Use editing features}: \emph{Find \& Replace} (\texttt{Ctrl+H}), \emph{Spelling \& Grammar} (\emph{Review $\rightarrow$ Spelling}), \emph{Track Changes} for collaboration.
\item \textbf{Generate automatic table of contents}: place cursor where TOC should appear $\rightarrow$ \emph{References $\rightarrow$ Table of Contents $\rightarrow$ Automatic Table 1}.
\item \textbf{Save and export}: \texttt{Ctrl+S} regularly; final delivery: \emph{File $\rightarrow$ Export $\rightarrow$ Create PDF/XPS Document}.
\end{enumerate}
\textbf{Reflexes:} Never press \texttt{Enter} repeatedly to create vertical space --- use \emph{Spacing Before/After} in Paragraph settings. Always use styles for headings. Save every 5 minutes (\texttt{Ctrl+S}).
\end{methode}

\begin{aretenir}[Word processor essentials for GCE]
\begin{itemize}
\item \textbf{Styles} = structure + automation (TOC, cross-references).
\item \textbf{Page break} (\texttt{Ctrl+Enter}) = clean new page, immune to editing above.
\item \textbf{Section breaks} = different headers/footers, orientation, or columns in the same document.
\item \textbf{Paste Special $\rightarrow$ Paste Link} = live link to spreadsheet data (updates when source changes).
\item \textbf{PDF export} = fixed layout, non-editable, standard for submission.
\end{itemize}
\end{aretenir}

\begin{attention}[Classic GCE traps in word processing]
\begin{itemize}
\item \textbf{Manual formatting} (bold + large font) instead of \emph{Heading 1} style $\rightarrow$ TOC stays empty.
\item \textbf{Multiple Enters} for spacing $\rightarrow$ pagination breaks when text is edited.
\item \textbf{Typing page numbers} by hand $\rightarrow$ they don't update.
\item \textbf{Forgetting to update fields} after editing: select all (\texttt{Ctrl+A}) $\rightarrow$ \texttt{F9} to refresh TOC, page numbers, cross-references.
\item \textbf{Not saving as .docx/.odt} before PDF $\rightarrow$ you lose the editable source.
\end{itemize}
\end{attention}

\begin{exoresolu}[Formatting a school report --- GCE style]
A form teacher asks you to produce the end-of-term report of Government High School Mfoundi. The document must have: a title page, an automatic table of contents, two sections with headings, a table of results, page numbers in the footer, and must be exported as PDF. Describe the steps you would follow in a word processor.
\tcblower
\textbf{Solution.}
\begin{enumerate}
\item Open the word processor and set the page: \emph{Layout $\rightarrow$ Size: A4; Margins: Normal (2.54 cm); Orientation: Portrait}.
\item \textbf{Title page}: type the school name (\texttt{Government High School Mfoundi}), the report title (\texttt{End-of-Term Report}), and the term (\texttt{Second Term 2024/2025}); centre the text, increase font size (e.g.\ 24 pt for school name, 18 pt for title), then insert a \emph{page break} (\texttt{Ctrl+Enter}) so the next content starts on a new page.
\item On page 2, leave space for the table of contents (to be generated later). Insert a page break after the TOC placeholder.
\item Type the body of the report. Select each main section title (e.g.\ \texttt{1. Academic Performance}, \texttt{2. Discipline}) and apply the style \emph{Heading 1}; apply \emph{Heading 2} to sub-sections (e.g.\ \texttt{1.1 Class Averages}). Apply the \emph{Normal} style to paragraphs, with \emph{justified} alignment and 1.15 line spacing.
\item Insert the results table: \emph{Insert $\rightarrow$ Table}, choose the number of rows and columns (e.g.\ 6 rows $\times$ 5 columns for 5 classes + header), type the data, then format the header row (bold, shading, centred).
\item Insert page numbers: \emph{Insert $\rightarrow$ Page Number $\rightarrow$ Bottom of Page $\rightarrow$ Plain Number 2}.
\item Generate the table of contents: place the cursor on page 2 (the TOC placeholder), then \emph{References $\rightarrow$ Table of Contents $\rightarrow$ Automatic Table 1}. The headings styled \emph{Heading 1/2} appear automatically with their page numbers.
\item Run the spelling checker (\emph{Review $\rightarrow$ Spelling \& Grammar}), save the file as \texttt{GHS\_Mfoundi\_Report\_Term2.docx}, then export: \emph{File $\rightarrow$ Export $\rightarrow$ Create PDF/XPS Document} $\rightarrow$ \texttt{GHS\_Mfoundi\_Report\_Term2.pdf}.
\end{enumerate}
\textbf{Key point:} the automatic table of contents works \emph{only} because heading styles were used --- manual bold text would not be detected. Always update the TOC (\texttt{F9}) after any edit.
\end{exoresolu}

\courssub{Formatting and Managing Spreadsheet Data}

\begin{experience}[Observing a messy marksheet]
Open \texttt{MessyMarks.xlsx}. Column A has ``Name'', B has ``Class'', C has ``Mark/20'', but some cells contain ``15/20'', others ``15'', one has ``fifteen''. Column D ``Coefficient'' has ``2'' and ``two''. The total row uses \texttt{=C2+C3+...} instead of \texttt{SUM}. Try to sort by mark --- the header row gets mixed in. Try to filter ``Upper Sixth A'' --- nothing happens because the data isn't recognised as a table. This is why \emph{formatting and managing} come before formulae.
\end{experience}

\begin{methode}[Formatting and managing data in a spreadsheet]
To present and manage data in Excel or LibreOffice Calc:
\begin{enumerate}
\item \textbf{Enter data as a structured table}: one field per column, one record per row; first row = column headers (unique, no blanks). Example: \texttt{Name | Class | Mark (/20) | Coefficient | Weighted Mark}.
\item \textbf{Format cells for meaning, not just looks}:
\begin{itemize}
\item Number format: \emph{Home $\rightarrow$ Number Format $\rightarrow$ Currency (FCFA), Percentage, Date, Number (0 decimals for marks)}.
\item \textbf{Never type units inside numeric cells} (``25000 FCFA'' becomes text). Put the unit in the column header and format the cell as Currency (symbol FCFA).
\item Fonts, borders, fill colours for header row; alignment (centre for numbers, left for text); \emph{Wrap Text} for long headers; \emph{Merge \& Center} only for report titles above the table.
\end{itemize}
\item \textbf{Adjust layout}: double-click column boundary to auto-fit; \emph{View $\rightarrow$ Freeze Panes $\rightarrow$ Freeze Top Row} so headers stay visible when scrolling.
\item \textbf{Convert to Excel Table} (\texttt{Ctrl+T}): enables structured references, auto-expansion, filter buttons on every header.
\item \textbf{Sort data}: click any cell in the table $\rightarrow$ \emph{Data $\rightarrow$ Sort} $\rightarrow$ choose key column (e.g.\ \texttt{Mark}) and order (Largest to Smallest). \textbf{Always include the header row} in the selection.
\item \textbf{Filter data}: \emph{Data $\rightarrow$ Filter} (or \texttt{Ctrl+Shift+L}) $\rightarrow$ drop-down arrows appear $\rightarrow$ tick/untick values (e.g.\ only \texttt{Upper Sixth A}). Hidden rows are not deleted.
\item \textbf{Validate entries} (optional but examined): \emph{Data $\rightarrow$ Data Validation} $\rightarrow$ Allow: List $\rightarrow$ Source: \texttt{Pass,Fail} $\rightarrow$ creates a drop-down in the cell, preventing typos.
\end{enumerate}
\textbf{Reflexes:} One header row, no blank rows/columns inside the data, numeric cells contain \emph{only numbers}, use \texttt{Ctrl+T} for a managed table.
\end{methode}

\begin{aretenir}[Spreadsheet data management]
\begin{itemize}
\item \textbf{Raw data} $\rightarrow$ \textbf{Format as Table} (\texttt{Ctrl+T}) $\rightarrow$ \textbf{Sort/Filter} $\rightarrow$ \textbf{Analyse}.
\item \textbf{Freeze Panes} keeps headers visible; essential for lists $>20$ rows.
\item \textbf{Data Validation} reduces input errors (examinable: ``restrict entry to a list'').
\item \textbf{Sort} rearranges rows; \textbf{Filter} hides rows. Both keep records intact.
\end{itemize}
\end{aretenir}

\begin{attention}[Spreadsheet data traps]
\begin{itemize}
\item \textbf{Text in numeric column} (``15/20'', ``Absent'') $\rightarrow$ \texttt{SUM}, \texttt{AVERAGE} ignore it $\rightarrow$ wrong totals.
\item \textbf{Sorting only one column} $\rightarrow$ data corruption (names no longer match marks). Always select the whole table.
\item \textbf{Merged cells} inside the data area $\rightarrow$ break sorting, filtering, formulas. Merge only for titles \emph{above} the table.
\item \textbf{Blank rows} inside data $\rightarrow$ filter/sort stops at the blank.
\end{itemize}
\end{attention}

\begin{exoresolu}[Managing a marks list --- GCE style]
The marks of 5 candidates in ICT (out of 20) are entered in a worksheet: A1:E1 contain the headers \texttt{Name}, \texttt{Class}, \texttt{Mark (/20)}, \texttt{Coefficient}, \texttt{Weighted}; rows 2 to 6 contain the records. (a) Explain how to display only candidates of class ``Upper Sixth A''. (b) Explain how to sort the candidates by decreasing mark. (c) The teacher typed ``15/20'' in a Mark cell and the SUM function returned a wrong total. Explain why and correct it.
\tcblower
\textbf{Solution.}
\begin{enumerate}
\item[(a)] Click any cell in the table, then \emph{Data $\rightarrow$ Filter} (AutoFilter) or press \texttt{Ctrl+Shift+L}. Drop-down arrows appear on each header. Click the arrow of the \texttt{Class} column, untick ``(Select All)'', tick only \texttt{Upper Sixth A}, and click \texttt{OK}. Only the matching rows remain visible; the others are hidden, not deleted. The status bar shows ``3 of 5 records found''.
\item[(b)] Select the whole table (headers included, e.g.\ A1:E6), then \emph{Data $\rightarrow$ Sort}. In the dialog: ``Sort by: Mark (/20)'', ``Sort On: Cell Values'', ``Order: Largest to Smallest''. Ensure ``My data has headers'' is ticked. Click \texttt{OK}. The rows are rearranged while each record stays intact.
\item[(c)] ``15/20'' is interpreted either as text or as a date (15/20 is not a valid date, so it stays \emph{text}). Text is ignored by \texttt{SUM}, hence the wrong total. \textbf{Correction:} type only the number \texttt{15} in the cell, and indicate ``out of 20'' in the column header, e.g.\ \texttt{Mark (/20)}. Format the column as \emph{Number} with 0 decimals. The sum is then correct.
\end{enumerate}
\textbf{Key point:} sorting acts on entire rows; filtering only hides rows; numeric cells must contain numbers only.
\end{exoresolu}

\courssub{Cell Referencing and Basic Formulae}

\begin{experience}[The \$ sign mystery]
In a sheet, B1 = 10. In D1 you type \texttt{=B1*2} and drag right to E1, F1. D1=20, E1=0 (because C1 is empty), F1=0. Now type \texttt{=\$B\$1*2} in D1 and drag right: D1=20, E1=20, F1=20. The \$ locked the reference. This is the single most examined concept in GCE spreadsheet questions.
\end{experience}

\begin{methode}[Using cell references and basic formulae]
\begin{enumerate}
\item Every formula begins with the sign \cle{=}.
\item \textbf{Relative reference} (e.g.\ \texttt{B2}): changes when the formula is copied (row and column both shift).
\item \textbf{Absolute reference} (e.g.\ \texttt{\$B\$2}): fixed when copied; the \$ locks the column \textbf{and} the row.
\item \textbf{Mixed references}: \texttt{\$B2} (column locked, row relative), \texttt{B\$2} (row locked, column relative). Press \texttt{F4} (Windows) or \texttt{Cmd+T} (Mac) to cycle: \texttt{B2} $\rightarrow$ \texttt{\$B\$2} $\rightarrow$ \texttt{B\$2} $\rightarrow$ \texttt{\$B2} $\rightarrow$ \texttt{B2}.
\item \textbf{Basic functions} (all take a range):
\begin{itemize}
\item \texttt{=SUM(range)} --- total
\item \texttt{=AVERAGE(range)} --- arithmetic mean
\item \texttt{=MAX(range)} --- highest value
\item \texttt{=MIN(range)} --- lowest value
\item \texttt{=COUNT(range)} --- count of \emph{numeric} cells
\item \texttt{=COUNTA(range)} --- count of \emph{non-empty} cells (text or numbers)
\end{itemize}
\item To copy a formula, drag the \cle{fill handle} (small square at bottom-right of active cell) or double-click it to auto-fill down a column.
\end{enumerate}
\textbf{Reflex:} any value that must stay fixed in a copied formula (a coefficient, a rate, a single total, a tax percentage) must be referenced with \$.
\end{methode}

\begin{aretenir}[Reference types --- decision guide]
\begin{center}
\begin{tabularx}{\linewidth}{|l|X|X|}\hline
\textbf{Reference} & \textbf{Behaviour when copied} & \textbf{Typical use} \\ \hline
\texttt{B2} & Both row and column change & Normal data cells (marks, quantities) \\ \hline
\texttt{\$B\$2} & Neither changes & Constants: tax rate, exchange rate, total coefficient \\ \hline
\texttt{B\$2} & Row fixed, column changes & Row headers used across columns \\ \hline
\texttt{\$B2} & Column fixed, row changes & Column headers used down rows \\ \hline
\end{tabularx}
\end{center}
\end{aretenir}

\begin{attention}[Formula traps in GCE]
\begin{itemize}
\item \textbf{Forgetting =} $\rightarrow$ formula displays as text.
\item \textbf{Using comma instead of semicolon} (or vice versa) depending on regional settings. Cameroon GCE uses \textbf{semicolon} (;) as list separator in functions: \texttt{=IF(A1>10;"Pass";"Fail")}.
\item \textbf{Relative reference where absolute needed} $\rightarrow$ copied formula points to wrong cell (classic: dividing by a moving total).
\item \textbf{Circular reference} (formula refers to its own cell) $\rightarrow$ error or 0.
\item \texttt{COUNT} vs \texttt{COUNTA}: \texttt{COUNT} counts \emph{numbers only}; \texttt{COUNTA} counts \emph{everything non-blank}.
\end{itemize}
\end{attention}

\begin{exoresolu}[Weighted average with absolute referencing]
In a worksheet, B2:B6 contain the marks of 5 students, C2:C6 the coefficients, and cell E1 contains the total of coefficients (computed by \texttt{=SUM(C2:C6)}). (a) Write the formula in D2 computing the weighted mark (mark $\times$ coefficient). (b) Write the formula in D7 computing the weighted average of the class. (c) Explain the role of the \$ signs if you used \texttt{\$E\$1}.
\tcblower
\textbf{Solution.}
\begin{enumerate}
\item[(a)] In D2: \texttt{=B2*C2}. Both references are \emph{relative}: when the formula is dragged down to D6, it becomes \texttt{=B3*C3}, \texttt{=B4*C4}, etc., which is exactly what we want (each student's own mark $\times$ own coefficient).
\item[(b)] The weighted average is
\[ \bar{m}=\frac{\sum m_i c_i}{\sum c_i} \]
so in D7: \texttt{=SUM(D2:D6)/SUM(C2:C6)}, or equivalently \texttt{=SUM(D2:D6)/\$E\$1} if the total of coefficients is stored in E1.
\item[(c)] If the formula \texttt{=SUM(D2:D6)/E1} were copied elsewhere (e.g.\ to the right), the reference E1 would move (F1, G1, \ldots) and give errors or wrong results; writing \texttt{\$E\$1} locks \textbf{both} the column and the row so the divisor always points to the same cell E1, wherever the formula is copied.
\end{enumerate}
\textbf{Key point:} relative references ``follow'' the copy; absolute references with \$ stay fixed. Choosing correctly is a classic GCE question.
\end{exoresolu}

\courssub{Advanced Functions, Charts and Linked Worksheets}

\begin{experience}[From ``IF'' to a dashboard]
A school bursar in Buea wants a sheet that automatically shows ``Scholarship'' if average $\ge$ 14, ``Probation'' if average $<$ 10, else ``Normal''. Then a chart of averages per class. Then the same sheet linked to a master workbook. This lesson chains \texttt{IF}, \texttt{VLOOKUP}, charts, and cross-sheet references --- exactly the ``advanced'' skills of Lessons 61--62.
\end{experience}

\begin{methode}[Logical, lookup functions and charts]
\begin{enumerate}
\item \textbf{IF function} (single condition):
\texttt{=IF(logical\_test; value\_if\_true; value\_if\_false)}.
Example: \texttt{=IF(B2>=14;"Scholarship";IF(B2<10;"Probation";"Normal"))} (nested IF).
\item \textbf{Conditional counting/summing}:
\begin{itemize}
\item \texttt{=COUNTIF(range; criterion)} --- e.g.\ \texttt{=COUNTIF(C2:C50;">=14")}
\item \texttt{=SUMIF(range; criterion; sum\_range)} --- e.g.\ sum fees for class ``A'': \texttt{=SUMIF(B2:B50;"A";D2:D50)}
\end{itemize}
\item \textbf{Lookup functions}:
\begin{itemize}
\item \texttt{=VLOOKUP(lookup\_value; table\_array; col\_index\_num; FALSE)} --- searches \textbf{first column} of table, returns value from \texttt{col\_index\_num}. \texttt{FALSE} = exact match (required for codes, names).
\item \texttt{=HLOOKUP(...)} --- horizontal lookup (rare in GCE).
\end{itemize}
\item \textbf{Charts} --- select data \textbf{with headers} $\rightarrow$ \emph{Insert $\rightarrow$ Chart}:
\begin{itemize}
\item \textbf{Column/Bar chart}: compare categories (e.g.\ marks per student, sales per agent).
\item \textbf{Line chart}: show evolution over \emph{ordered} time (days, months, terms).
\item \textbf{Pie chart}: show \textbf{proportions} of a whole (e.g.\ budget allocation, gender distribution). Only one data series.
\item \textbf{Scatter (XY) chart}: relationship between two numeric variables (e.g.\ study hours vs mark).
\end{itemize}
\textbf{Always add}: Chart Title, Axis Titles, Legend (if multiple series), Data Labels (if few points).
\end{enumerate}
\textbf{Reflex:} text criteria in formulas are written between double quotes (\texttt{"Pass"}); numeric criteria without quotes (\texttt{>10}); cell references without quotes (\texttt{B2}).
\end{methode}

\begin{aretenir}[Chart choice --- examiner's checklist]
\begin{center}
\begin{tabularx}{\linewidth}{|l|X|X|}\hline
\textbf{Question asks for...} & \textbf{Choose} & \textbf{Why} \\ \hline
``Compare values across categories'' & Column / Bar & Length encodes magnitude \\ \hline
``Show trend / evolution over time'' & Line & Continuity emphasises change \\ \hline
``Show percentage share / proportion'' & Pie / Doughnut & Angle encodes part-of-whole \\ \hline
``Show relationship between two measures'' & Scatter (XY) & Both axes numeric \\ \hline
\end{tabularx}
\end{center}
\end{aretenir}

\begin{attention}[Advanced function traps]
\begin{itemize}
\item \texttt{VLOOKUP} without \texttt{FALSE} (or \texttt{0}) $\rightarrow$ approximate match $\rightarrow$ wrong results if first column not sorted.
\item \texttt{VLOOKUP} column index counted from \textbf{start of table\_array}, not from column A.
\item Nested IFs: close every parenthesis; maximum 64 levels (but 3--4 is typical in GCE).
\item \texttt{COUNTIF} criterion syntax: \texttt{">=10"} (quotes), \texttt{B2} (no quotes), \texttt{"Pass"} (quotes).
\item Charting the \textbf{wrong range} (e.g.\ including the total row) $\rightarrow$ distorted chart.
\end{itemize}
\end{attention}

\begin{exoresolu}[Decision function and chart for a canteen]
A school canteen in Bamenda records daily sales of meals in B2:B8 (Monday to Sunday). (a) Write the formula in C2 displaying ``Good day'' if sales exceed 100 meals, ``Low day'' otherwise. (b) Write the formula counting the number of good days. (c) Which chart type best shows the evolution of sales during the week? Justify.
\tcblower
\textbf{Solution.}
\begin{enumerate}
\item[(a)] In C2: \texttt{=IF(B2>100;"Good day";"Low day")}, then copy down to C8. The test \texttt{B2>100} is evaluated for each row; because B2 is relative, it adapts to each day.
\item[(b)] \texttt{=COUNTIF(C2:C8;"Good day")} counts the cells of the range containing exactly the text ``Good day''. One could also count directly on the numbers: \texttt{=COUNTIF(B2:B8;">100")}.
\item[(c)] A \textbf{line chart} is the most appropriate: the days form an ordered time sequence (Mon $\rightarrow$ Sun) and a line emphasises the \emph{evolution} (rises and falls) of sales during the week. A column chart would also be acceptable for comparing the days, but a pie chart would be wrong here, because a pie chart shows shares of a total, not an evolution.
\end{enumerate}
\textbf{Key point:} choose the chart according to the message: comparison $\rightarrow$ columns, proportion $\rightarrow$ pie, evolution $\rightarrow$ line.
\end{exoresolu}

\courssub{Linking Worksheets}

\begin{methode}[Linking worksheets and workbooks]
\begin{enumerate}
\item \textbf{Same workbook, different sheet}: \texttt{=SheetName!CellReference}. If the sheet name contains spaces or special characters, enclose in single quotes: \texttt{='Term 1'!B5}.
\item \textbf{3D reference (consolidation across adjacent sheets with identical layout)}: \texttt{=SUM(Sheet1:Sheet3!B2)} adds cell B2 from Sheet1, Sheet2, and Sheet3. Works for \texttt{SUM}, \texttt{AVERAGE}, \texttt{MAX}, \texttt{MIN}, \texttt{COUNT}.
\item \textbf{Different workbook (external link)}: \texttt{=[WorkbookName.xlsx]SheetName!CellReference}. Example: \texttt{=[Budget2025.xlsx]Sheet1!B5}. The source workbook must be open or the path included.
\item \textbf{Updating links}: when the source changes, the linked cell recalculates. On opening a workbook with external links, Excel asks ``Update links?'' $\rightarrow$ click \texttt{Update}.
\item \textbf{Break links} (if needed): \emph{Data $\rightarrow$ Edit Links $\rightarrow$ Break Link} $\rightarrow$ converts formulas to static values.
\end{enumerate}
\textbf{Reflex:} linking avoids duplicating data --- one single source, updated once, reflected everywhere. This is the principle of \cle{data integrity}.
\end{methode}

\begin{aretenir}[Linking syntax summary]
\begin{center}
\begin{tabular}{|l|l|}\hline
\textbf{Situation} & \textbf{Syntax example} \\ \hline
Same sheet & \texttt{=B2} \\ \hline
Different sheet, no spaces & \texttt{=Term1!B2} \\ \hline
Different sheet, with spaces & \texttt{='Term 1'!B2} \\ \hline
3D sum across Term1, Term2, Term3 & \texttt{=SUM(Term1:Term3!B2)} \\ \hline
External workbook (open) & \texttt{=[Budget.xlsx]Sheet1!B2} \\ \hline
External workbook (closed, full path) & \texttt{='C:\textbackslash{}Data\textbackslash{}[Budget.xlsx]Sheet1'!B2} \\ \hline
\end{tabular}
\end{center}
\end{aretenir}

\begin{exoresolu}[Consolidating term results --- GCE style]
A school keeps three worksheets in one workbook: \texttt{Term1}, \texttt{Term2}, \texttt{Term3}, each with the same layout; the total marks of each student are in column F (F2 for the first student). A summary sheet \texttt{Annual} must show, for the first student, the annual total and the annual average. Write the required formulas.
\tcblower
\textbf{Solution.}
\begin{enumerate}
\item Annual total in \texttt{Annual!B2}: \texttt{=Term1!F2+Term2!F2+Term3!F2}, or more compactly with a 3D sum: \texttt{=SUM(Term1:Term3!F2)}.
\item Annual average in \texttt{Annual!C2}: \texttt{=B2/3} or directly \texttt{=AVERAGE(Term1:Term3!F2)}.
\item Copy the formulas down for all students: the row number adapts automatically (relative references), while the sheet names remain unchanged.
\end{enumerate}
\textbf{Key point:} the syntax \texttt{Sheet!Cell} is the heart of linking; the 3D form \texttt{Sheet1:Sheet3!Cell} is powerful when all sheets share the same structure.
\end{exoresolu}

\courssub{What-If Analysis}

\begin{experience}[The ``reverse'' problem]
A student knows her CA mark (14/20) and the formula \texttt{Final = 0.4*CA + 0.6*Exam}. She wants a final mark of 12. What exam mark does she need? She could guess-and-check, or solve the equation. \cle{Goal Seek} does the algebra for her. This is \cle{what-if analysis}: changing inputs to hit a target output.
\end{experience}

\begin{methode}[Performing a what-if analysis]
\cle{What-if analysis} consists of changing input values to observe the effect on results computed by formulas.
\begin{enumerate}
\item \textbf{Build the model with formulas} (never typed results), so that outputs depend on inputs. Identify: \emph{input cells} (variables you can change), \emph{formula cells} (calculated results).
\item \textbf{Manual what-if}: change an input cell and observe the recalculated outputs. Good for quick exploration.
\item \textbf{Goal Seek} (\emph{Data $\rightarrow$ What-If Analysis $\rightarrow$ Goal Seek}): finds the input value needed to reach a target result. Fill in:
\begin{itemize}
\item \emph{Set cell}: the formula cell (the result you want to hit).
\item \emph{To value}: the target number.
\item \emph{By changing cell}: the input cell you can vary.
\end{itemize}
\item \textbf{Data Table} (one or two variables): shows the result of a formula for a whole series of input values. \emph{Data $\rightarrow$ What-If Analysis $\rightarrow$ Data Table}.
\item \textbf{Scenario Manager}: save and compare several sets of input values (e.g.\ Optimistic / Normal / Pessimistic budgets). \emph{Data $\rightarrow$ What-If Analysis $\rightarrow$ Scenario Manager}.
\end{enumerate}
\textbf{Reflex:} Goal Seek requires that the ``Set cell'' contains a \emph{formula} depending (directly or indirectly) on the ``By changing cell''.
\end{methode}

\begin{aretenir}[What-if tools comparison]
\begin{center}
\begin{tabularx}{\linewidth}{|l|X|X|}\hline
\textbf{Tool} & \textbf{Use when...} & \textbf{Output} \\ \hline
Manual change & Quick test, one or two values & Immediate recalculation \\ \hline
Goal Seek & One target, one unknown input & Single solution value \\ \hline
Data Table (1 var) & See results for many values of one input & Column/row of results \\ \hline
Data Table (2 var) & See results for combinations of two inputs & Matrix of results \\ \hline
Scenario Manager & Compare a few named scenarios (Best/Worst) & Switchable named sets \\ \hline
\end{tabularx}
\end{center}
\end{aretenir}

\begin{attention}[Goal Seek traps]
\begin{itemize}
\item ``Set cell'' must be a \textbf{formula}, not a constant.
\item ``By changing cell'' must be an \textbf{input} (no formula), and must affect the set cell.
\item Goal Seek finds \textbf{one} solution; non-linear models may have multiple solutions.
\item It uses numerical iteration; result may be approximate (e.g.\ 10.6666667). Round appropriately for context.
\end{itemize}
\end{attention}

\begin{exoresolu}[Goal Seek: mark needed at the GCE]
A candidate's final mark is computed as \texttt{Final = 0.4*CA + 0.6*Exam}, where CA (continuous assessment, in B1) is already known: 14/20, and the Exam mark is in B2. The final mark is in B3. (a) Write the formula in B3. (b) The candidate wants a final mark of at least 12/20. Use Goal Seek to find the minimum exam mark. (c) Compute the result by hand to verify.
\tcblower
\textbf{Solution.}
\begin{enumerate}
\item[(a)] In B3: \texttt{=0.4*B1+0.6*B2}.
\item[(b)] \emph{Data $\rightarrow$ What-If Analysis $\rightarrow$ Goal Seek}: Set cell = \texttt{B3}; To value = \texttt{12}; By changing cell = \texttt{B2}. Click \texttt{OK}. Goal Seek reports the value B2 $\approx 10.6667$. The candidate therefore needs at least about $10.67/20$ at the exam (round up to 11 if marks are integers).
\item[(c)] Verification by algebra:
\begin{align*}
0.4\times 14 + 0.6\times x &= 12\\
5.6 + 0.6x &= 12\\
0.6x &= 6.4\\
x &= \frac{6.4}{0.6} = \frac{64}{6} = \frac{32}{3} \approx 10.67
\end{align*}
The spreadsheet result is confirmed.
\end{enumerate}
\textbf{Key point:} Goal Seek automates exactly this ``reverse'' calculation; it only works because B3 is a formula linked to B2.
\end{exoresolu}

\courssub{Integration Activity: from Data to a Complete Report}

\begin{experience}[The real GCE Paper 3 task]
The practical exam gives you: a raw data file (CSV or Excel), a word processor template, and a list of requirements: ``Compute X, create chart Y, write a report with title, TOC, conclusions, page numbers, export to PDF.'' You have 2h30. The only way to finish is a rehearsed workflow: \textbf{Spreadsheet first (clean data, formulas, chart)}, then \textbf{Word (structure, paste linked objects, finalise)}.
\end{experience}

\begin{methode}[Integrating spreadsheet and word processor --- full workflow]
Typical GCE integration task: analyse data in a spreadsheet, then present conclusions in a word-processed document.
\begin{enumerate}
\item \textbf{Spreadsheet part}:
\begin{enumerate}
\item Import/open raw data; convert to Table (\texttt{Ctrl+T}); clean (remove duplicates, fix text-in-numbers).
\item Compute required columns with formulas (use \$ for constants).
\item Apply logical functions (\texttt{IF}, \texttt{COUNTIF}, \texttt{SUMIF}, \texttt{VLOOKUP}) as asked.
\item Create the required chart: select correct range $\rightarrow$ correct chart type $\rightarrow$ add titles, labels.
\item Verify: test formulas with simple values; check copy behaviour (relative/absolute).
\end{enumerate}
\item \textbf{Transfer to word processor}:
\begin{enumerate}
\item In spreadsheet: select table $\rightarrow$ \texttt{Ctrl+C}.
\item In word processor: \emph{Home $\rightarrow$ Paste $\rightarrow$ Paste Special $\rightarrow$ Paste Link $\rightarrow$ Microsoft Excel Worksheet Object} (or \texttt{LibreOffice Spreadsheet}). This keeps the document synchronised.
\item For chart: click chart $\rightarrow$ \texttt{Ctrl+C} $\rightarrow$ in Word: \emph{Paste Special $\rightarrow$ Paste Link $\rightarrow$ Chart Object}.
\end{enumerate}
\item \textbf{Document part}:
\begin{enumerate}
\item Apply \emph{Title}, \emph{Heading 1}, \emph{Heading 2}, \emph{Normal} styles.
\item Write introduction, methodology, findings (referencing the linked table/chart), conclusion.
\item Insert automatic Table of Contents (\emph{References $\rightarrow$ Table of Contents}).
\item Insert page numbers, headers/footers (name, candidate number, date).
\item \texttt{Ctrl+A} $\rightarrow$ \texttt{F9} to update all fields.
\item Save \texttt{.docx}, then \emph{Export to PDF}.
\end{enumerate}
\end{enumerate}
\textbf{Reflex:} a \emph{linked} paste keeps the document synchronised with the spreadsheet; a simple paste freezes a copy. In the exam, \textbf{always use Paste Link} unless explicitly told ``paste as picture''.
\end{methode}

\begin{aretenir}[Integration checklist --- tick before submitting]
\begin{itemize}
\item[$\square$] Spreadsheet: all formulas use cell references (no typed numbers), \$ used correctly.
\item[$\square$] Spreadsheet: chart type matches the question (evolution $\rightarrow$ line, comparison $\rightarrow$ column, proportion $\rightarrow$ pie).
\item[$\square$] Word: heading styles used $\rightarrow$ TOC generates correctly.
\item[$\square$] Word: table and chart pasted as \textbf{links} (right-click $\rightarrow$ Linked Worksheet Object $\rightarrow$ Update Link works).
\item[$\square$] Word: page numbers, footer with name/candidate number, title page.
\item[$\square$] Both files saved (\texttt{.xlsx}, \texttt{.docx}) and PDF exported.
\end{itemize}
\end{aretenir}

\begin{exoresolu}[Complete integration task: mobile money agency]
An MTN Mobile Money agency in Douala records, in a worksheet, for each of 6 agents: the number of daily transactions (column B) and the total amount in FCFA (column C). You must: (1) compute in D2 the average amount per transaction for the first agent; (2) display in E2 ``Bonus'' if the agent made more than 50 transactions, ``No bonus'' otherwise; (3) count the agents receiving a bonus; (4) produce a column chart comparing the agents' transaction numbers; (5) insert the table and the chart into a one-page word-processed report with a title, a short conclusion, and your name in the footer.
\tcblower
\textbf{Solution.}
\begin{enumerate}
\item In D2: \texttt{=C2/B2} (amount divided by number of transactions), formatted as Currency (FCFA, 0 decimals) via \emph{Format Cells $\rightarrow$ Currency}, then copied down to D7.
\item In E2: \texttt{=IF(B2>50;"Bonus";"No bonus")}, copied down to E7.
\item In a cell below, e.g.\ E9: \texttt{=COUNTIF(E2:E7;"Bonus")}.
\item Select the agents' names (A1:A7) and the transaction column (B1:B7) $\rightarrow$ \emph{Insert $\rightarrow$ Chart $\rightarrow$ Clustered Column}. Add Chart Title: ``Daily Transactions per Agent'', Horizontal Axis Title: ``Agent'', Vertical Axis Title: ``Number of Transactions''. Remove legend (single series).
\item In the word processor:
\begin{enumerate}
\item Type the title ``MTN Mobile Money Douala --- Daily Performance Report'' with \emph{Title} style.
\item Write one short paragraph of conclusion, e.g.: ``Agent 3 (Ngono) recorded the highest transactions (62) and qualifies for bonus. Three out of six agents exceeded the 50-transaction threshold. The agency should investigate Agent 2's low volume (28).''
\item Copy the table (A1:E7) from Excel $\rightarrow$ in Word: \emph{Paste Special $\rightarrow$ Paste Link $\rightarrow$ Microsoft Excel Worksheet Object}. Resize to fit page width.
\item Copy the chart from Excel $\rightarrow$ in Word: \emph{Paste Special $\rightarrow$ Paste Link $\rightarrow$ Chart Object}. Position below the table.
\item Insert footer: \emph{Insert $\rightarrow$ Footer $\rightarrow$ Blank} $\rightarrow$ type your name and candidate number.
\item Save as \texttt{MM\_Douala\_Report.docx}, update fields (\texttt{F9}), export to PDF \texttt{MM\_Douala\_Report.pdf}.
\end{enumerate}
\end{enumerate}
\textbf{Key point:} this task chains all the skills of the chapter --- formulae, logical functions, conditional counting, charts, linking and document production --- exactly as in the integration activities of the syllabus (Lesson 64).
\end{exoresolu}

\begin{exercice}[Practice: End-to-end GCE simulation]
You are given a file \texttt{ExamData.xlsx} with two sheets: \texttt{RawMarks} (columns: RegNo, Name, Class, Paper1, Paper2, Paper3) and \texttt{Coeffs} (Paper1=2, Paper2=3, Paper3=5 in B2:B4). In a new sheet \texttt{Results}:
\begin{enumerate}
\item Link the student data from \texttt{RawMarks} (use \texttt{=RawMarks!A2} etc.).
\item Compute \texttt{WeightedTotal} = Paper1*2 + Paper2*3 + Paper3*5 using absolute references to \texttt{Coeffs!B2:B4}.
\item Compute \texttt{Average} = WeightedTotal / 10 (sum of coefficients).
\item In \texttt{Decision} column: \texttt{=IF(Average>=12;"Pass";IF(Average>=10;"Probation";"Fail"))}.
\item Count Pass, Probation, Fail with \texttt{COUNTIF}.
\item Create a column chart of Average per student (sorted descending).
\item Copy the Results table and chart into a Word report with title, TOC, conclusion, page numbers, footer with your name. Export to PDF.
\end{enumerate}
\textit{Time yourself: target 45 minutes for spreadsheet, 30 minutes for report.}
\end{exercice}

\begin{corrige}[Exercise: End-to-end GCE simulation]
\begin{enumerate}
\item In \texttt{Results!A2}: \texttt{=RawMarks!A2}, copy across to D2 (Name, Class), then down.
\item In \texttt{Results!F2} (WeightedTotal): \texttt{=C2*\$Coeffs.\$B\$2+D2*\$Coeffs.\$B\$3+E2*\$Coeffs.\$B\$4} (assuming Paper1 in C, Paper2 in D, Paper3 in E). Copy down.
\item In \texttt{Results!G2} (Average): \texttt{=F2/SUM(\$Coeffs.\$B\$2:\$Coeffs.\$B\$4)}. Copy down.
\item In \texttt{Results!H2} (Decision): \texttt{=IF(G2>=12;"Pass";IF(G2>=10;"Probation";"Fail"))}. Copy down.
\item Counts: \texttt{=COUNTIF(H2:H100;"Pass")}, \texttt{=COUNTIF(H2:H100;"Probation")}, \texttt{=COUNTIF(H2:H100;"Fail")}.
\item Chart: sort \texttt{Results} by Average descending (Data $\rightarrow$ Sort), select Name and Average columns $\rightarrow$ Insert $\rightarrow$ Column Chart.
\item Word: follow Method 7 workflow. Verify linked objects update when you change a mark in \texttt{RawMarks}.
\end{enumerate}
\end{corrige}

\newpage

\coursec{Exercises}

\courssub{I check my knowledge}
\begin{exercice}[MCQ on cell referencing]
Which of the following formulas, when copied from cell E2 to F4, will keep the reference to cell C1 fixed?
\begin{enumerate}
\item =B2*C1+D2
\item =B2*\$C\$1+\$D2
\item =\$B\$2*C1+D\$2
\item =B2*C\$1+\$D\$2
\end{enumerate}
\end{exercice}

\begin{exercice}[True/False with justification]
State whether each statement is true or false and justify your answer.
\begin{enumerate}
\item In a spreadsheet, a relative reference changes when the formula is copied to another cell.
\item The function SUMIF can only be used with numeric criteria.
\item A chart of type ``Pie'' is the most appropriate to show the evolution of monthly income over a year.
\item The error \#NAME? appears when a formula tries to divide by zero.
\end{enumerate}
\end{exercice}

\begin{exercice}[Definitions]
Define the following terms as used in spreadsheet applications:
\begin{enumerate}
\item Absolute cell reference
\item Mixed cell reference
\item What‑If Analysis
\item Mail merge
\end{enumerate}
\end{exercice}

\begin{exercice}[Direct calculations]
A student obtained the following marks in three sequences: 8 (weight 1), 9 (weight 1), 11 (weight 2). Compute the weighted average.
\end{exercice}

\courssub{I practise}
\begin{exercice}[SUMIF formula for fee sheet]
A worksheet contains 50 rows with columns: A (Name), B (Class), C (Amount Paid). Write the formula that calculates the total amount paid by students in ``Form Five'' only.
\end{exercice}

\begin{exercice}[Relative vs absolute referencing]
Explain the difference between relative and absolute referencing. Give one practical situation where each type is required.
\end{exercice}

\begin{exercice}[Goal Seek scenario]
A student has averages of 8, 9 and 11 in three sequences weighted 1, 1 and 2 respectively. Using the logic of Goal Seek, determine the mark needed in a fourth sequence (weight 2) to achieve an overall weighted average of 10. Show the equation you would set up.
\end{exercice}

\begin{exercice}[Chart selection]
For each of the following data sets, choose the most appropriate chart type and justify your choice:
\begin{enumerate}
\item Monthly school income over a 12‑month period.
\item Percentage share of students per class (Form 1 to Upper Sixth).
\item Comparison of boys' and girls' average marks in five subjects.
\end{enumerate}
\end{exercice}

\courssub{I prepare for the GCE}
\begin{exercice}[Mail merge procedure (10 marks)]
List, in the correct order, the steps required to produce 200 personalised fee‑reminder letters using mail merge from an Excel data source. Name the three main components involved in the mail merge process.
\end{exercice}

\begin{exercice}[Error diagnosis (12 marks)]
A worksheet displays the error \#DIV/0! in the ``Average'' column and \#NAME? in another cell.
\begin{enumerate}
\item Diagnose the probable cause of each error.
\item Propose a correction for each error.
\end{enumerate}
\end{exercice}

\begin{exercice}[Linked worksheets and renaming (13 marks)]
A school keeps term results on three separate sheets named ``Term1'', ``Term2'' and ``Term3''. On a ``Summary'' sheet you want to add a student's three term totals.
\begin{enumerate}
\item Write the formula on the ``Summary'' sheet that adds the three term totals for a student whose total is in cell B10 on each term sheet.
\item Explain what happens to the link if the ``Term2'' sheet is renamed to ``SecondTerm''.
\item Suggest a method to avoid broken links when sheets are renamed.
\end{enumerate}
\end{exercice}

\newpage

\coursec{Solutions to the exercises}

\courssub{Solutions -- I check my knowledge}

\begin{corrige}[Exercise 1 -- MCQ on cell referencing]
\textbf{Correct answer: (b) =B2*\$C\$1+\$D2}

\textbf{Justification.} The requirement is that the reference to \textbf{C1} must remain \emph{fixed} when the formula is copied from E2 to F4. A reference is fixed only if \emph{both} its column letter and its row number are preceded by the dollar sign \$, i.e.\ \textbf{\$C\$1} (\cle{absolute reference}).

\begin{itemize}
\item Option (a): C1 is fully relative; copied from E2 to F4 (one column right, two rows down) it becomes D3. \textbf{Wrong.}
\item Option (b): C1 is written \$C\$1, so it stays C1 wherever the formula is copied. \textbf{Correct.}
\item Option (c): C1 is fully relative (it becomes D3); only D\$2 is partly fixed. \textbf{Wrong.}
\item Option (d): C\$1 fixes only the \emph{row} (1), but the column C changes to D when copied one column to the right, giving D\$1. \textbf{Wrong.}
\end{itemize}

Note that in option (b) the term \$D2 is a \cle{mixed reference}: the column D is fixed, the row is free. This does not affect the question, which concerns only C1.
\end{corrige}

\begin{corrige}[Exercise 2 -- True/False with justification]
\begin{enumerate}
\item \textbf{TRUE.} By definition, a \cle{relative reference} (e.g.\ B2) is interpreted as a position \emph{relative} to the cell containing the formula. When the formula is copied, the reference is automatically adjusted by the same displacement: copying =B2 from E2 down to E3 gives =B3.
\item \textbf{FALSE.} \cle{SUMIF} also accepts \emph{text} criteria and criteria with comparison operators, e.g.\ =SUMIF(B2:B50;``Form Five'';C2:C50) or =SUMIF(C2:C50;``>50000''). The criterion is not restricted to numbers.
\item \textbf{FALSE.} A \cle{pie chart} shows \emph{proportions of a whole at one moment}; it cannot display evolution over time. To show the evolution of monthly income over a year, a \cle{line chart} (or a column chart) is appropriate, with months on the horizontal axis.
\item \textbf{FALSE.} \#NAME? appears when the spreadsheet does not \emph{recognise a name} in the formula (misspelt function such as =SOMME(A1:A5) in an English version, or an undefined named range). Division by zero produces the error \textbf{\#DIV/0!}.
\end{enumerate}
\end{corrige}

\begin{corrige}[Exercise 3 -- Definitions]
\begin{enumerate}
\item \textbf{Absolute cell reference:} a cell reference in which both the column and the row are locked with the dollar sign (e.g.\ \$C\$1). It always designates the \emph{same physical cell}, whatever the destination of a copy of the formula.
\item \textbf{Mixed cell reference:} a reference in which only one of the two coordinates is locked: either the column (\$C1) or the row (C\$1). On copying, the locked part stays fixed while the free part is adjusted.
\item \textbf{What-If Analysis:} the set of spreadsheet techniques used to study how a change in input values affects the results of formulas. It includes tools such as \emph{Goal Seek} (find the input needed to reach a target result), \emph{Scenario Manager} and \cle{Data Tables}.
\item \textbf{Mail merge:} a word-processing feature that automatically combines a \emph{main document} (containing the fixed text and merge fields) with a \emph{data source} (a table of recipients, e.g.\ an Excel sheet) to produce many personalised documents (letters, certificates, fee reminders), one per record.
\end{enumerate}
\end{corrige}

\begin{corrige}[Exercise 4 -- Direct calculations]
The weighted average is the sum of (mark $\times$ weight) divided by the sum of the weights:
\[
\bar{x}=\frac{\sum x_i\,n_i}{\sum n_i}=\frac{8\times1+9\times1+11\times2}{1+1+2}=\frac{8+9+22}{4}=\frac{39}{4}=9.75.
\]
\textbf{The weighted average is 9.75 out of 20.}

In a spreadsheet, with marks in B2:B4 and weights in C2:C4, the same result is obtained by:
\[ =\text{SUMPRODUCT(B2:B4;C2:C4)/SUM(C2:C4)} \]
\textbf{Note:} In English-language versions of Excel the list separator is a comma (,); in French versions or LibreOffice Calc it is a semicolon (;). Use the separator corresponding to your regional settings.
\end{corrige}

\courssub{Solutions -- I practise}

\begin{corrige}[Exercise 5 -- SUMIF formula for fee sheet]
The \cle{SUMIF} function has the syntax \texttt{=SUMIF(range; criteria; sum\_range)}. Here the class is tested in column B and the amounts to add are in column C:
\[ =\text{SUMIF(B2:B51;``Form Five'';C2:C51)} \]
(50 data rows, from row 2 to row 51, assuming row 1 contains the headings.)

\textbf{Remark:} the text criterion must be written exactly as it appears in the sheet (``Form Five''), between double quotes. If the class name were stored in a cell, say E1, one would write \texttt{=SUMIF(B2:B51;E1;C2:C51)}, which is more flexible and avoids hard-coding.
\end{corrige}

\begin{corrige}[Exercise 6 -- Relative vs absolute referencing]
\textbf{Relative referencing} (e.g.\ B2): the reference designates a cell by its \emph{position relative} to the formula. When the formula is copied, the reference moves by the same displacement (copy one row down $\Rightarrow$ B2 becomes B3).

\textbf{Absolute referencing} (e.g.\ \$B\$2): the reference designates a \emph{fixed} cell; it never changes on copying.

\textbf{Practical situations:}
\begin{itemize}
\item \emph{Relative:} computing each student's term total with \texttt{=SUM(B2:D2)} in E2, then copying the formula down for the whole class -- each row must sum \emph{its own} marks.
\item \emph{Absolute:} converting amounts in CFA francs using an exchange rate stored in a single cell H1: the formula \texttt{=B2*\$H\$1} copied down the column must always read the \emph{same} rate cell. Similarly, applying a single school-fees amount or a VAT rate stored once to many rows.
\end{itemize}
\end{corrige}

\begin{corrige}[Exercise 7 -- Goal Seek scenario]
Let $x$ be the mark needed in the fourth sequence (weight 2). The overall weighted average must equal 10:
\begin{align*}
\frac{8\times1+9\times1+11\times2+x\times2}{1+1+2+2}&=10 \\
\frac{39+2x}{6}&=10 \\
39+2x&=60 \\
2x&=21 \\
x&=10.5.
\end{align*}
\textbf{The student needs 10.5 out of 20 in the fourth sequence.}

\textbf{With Goal Seek in the spreadsheet:} put the marks and weights in cells, compute the average in a cell (say F5) with \texttt{=SUMPRODUCT(...)/SUM(...)}, leaving the fourth mark cell empty. Then open \emph{Data $\rightarrow$ What-If Analysis $\rightarrow$ Goal Seek} and set: \textbf{Set cell} = F5, \textbf{To value} = 10, \textbf{By changing cell} = the fourth mark cell. Goal Seek returns 10.5 automatically.
\end{corrige}

\begin{corrige}[Exercise 8 -- Chart selection]
\begin{enumerate}
\item \textbf{Monthly school income over 12 months: line chart} (a column chart is also acceptable). Justification: a line chart is the standard tool for showing the \emph{evolution/trend} of a quantity over time; the months go on the horizontal axis.
\item \textbf{Percentage share of students per class: pie chart} (or 100\% stacked bar). Justification: the data are \emph{parts of a whole} (all shares sum to 100\%), which is exactly what a pie chart displays.
\item \textbf{Boys' vs girls' average marks in five subjects: clustered (grouped) column/bar chart} with two series. Justification: it allows a direct \emph{comparison of two categories} across several items (subjects), with one column per sex side by side for each subject.
\end{enumerate}
\end{corrige}

\courssub{Solutions -- I prepare for the GCE}

\begin{corrige}[Exercise 9 -- Mail merge procedure (10 marks)]
\textbf{The three main components of a mail merge} (3 marks):
\begin{enumerate}
\item The \textbf{main document} -- the letter template containing the fixed text and the merge fields;
\item The \textbf{data source} -- the Excel sheet (or CSV, Access table) with one row per student (Name, Class, Amount owed, \dots);
\item The \textbf{merged document(s)} -- the 200 personalised letters produced by the combination.
\end{enumerate}

\textbf{Steps in the correct order} (7 marks -- 1 mark per correctly ordered step, any 7 of the following):
\begin{enumerate}
\item Prepare and save the Excel data source: one header row with field names (e.g.\ \texttt{Name}, \texttt{Class}, \texttt{Amount}), one record per student, no blank rows.
\item Open the word processor (MS Word or LibreOffice Writer) and create (or open) the main document; type the fixed text of the fee-reminder letter.
\item Start the mail merge: \emph{Mailings $\rightarrow$ Start Mail Merge $\rightarrow$ Letters} (Word) or \emph{Tools $\rightarrow$ Mail Merge Wizard} (LibreOffice).
\item Select the recipients: \emph{Select Recipients $\rightarrow$ Use an Existing List}, then browse to the Excel file and choose the correct worksheet.
\item Insert the \textbf{merge fields} (e.g.\ \texttt{<<Name>>}, \texttt{<<Class>>}, \texttt{<<Amount>>}) at the appropriate positions in the letter.
\item (Optional) Set rules or filter/sort the recipient list if only some students should receive the letter (e.g.\ only those with Amount > 0).
\item \textbf{Preview the results} to check that the fields are correctly replaced for several records.
\item Complete the merge: \emph{Finish \& Merge $\rightarrow$ Edit Individual Documents} (produces all 200 letters in a new file) or \emph{Print Documents}.
\item Save and/or print the 200 personalised letters.
\end{enumerate}

\textbf{Examiner expectations:} the candidate must show that the data source is prepared \emph{before} the connection, that merge fields are inserted \emph{before} the merge, and must name the three components precisely. Vague answers such as ``click mail merge'' without the logical order lose marks.
\end{corrige}

\begin{corrige}[Exercise 10 -- Error diagnosis (12 marks)]
\textbf{(a) Diagnosis} (6 marks -- 3 per error):

\textbf{\#DIV/0! in the ``Average'' column:} the formula attempts a \emph{division by zero or by an empty cell}. Typical cause: \texttt{=SUM(B2:D2)/E2} where the divisor E2 (number of subjects or total coefficient) is empty or equal to 0 -- for example a newly added student whose coefficient has not yet been entered, or a copied formula whose divisor reference points to a blank cell.

\textbf{\#NAME?:} the spreadsheet \emph{does not recognise a text} in the formula. Typical causes: a misspelt function name (\texttt{=AVEARGE(B2:D2)} or the French \texttt{=MOYENNE(...)} in an English version), a named range that does not exist or was deleted, or text entered in a formula without double quotes.

\textbf{(b) Corrections} (6 marks -- 3 per error):

\textbf{For \#DIV/0!:}
\begin{itemize}
\item Enter the missing value in the divisor cell, or correct the divisor reference; \emph{and/or}
\item Make the formula robust with a test: \texttt{=IF(E2=0;``'';SUM(B2:D2)/E2)} or \texttt{=IFERROR(SUM(B2:D2)/E2;``--'')}, so that the sheet stays clean while data are incomplete.
\end{itemize}

\textbf{For \#NAME?:}
\begin{itemize}
\item Check and correct the spelling of the function (use the function wizard or autocomplete);
\item Recreate or correct the named range (\emph{Formulas $\rightarrow$ Name Manager});
\item Put any literal text used in the formula between double quotes.
\end{itemize}

\textbf{Examiner expectations:} one mark for identifying the error type, one for a plausible concrete cause, one for a valid correction, per error. Simply restating ``division by zero'' for \#DIV/0! without proposing a correction earns only partial credit.
\end{corrige}

\begin{corrige}[Exercise 11 -- Linked worksheets and renaming (13 marks)]
\textbf{(a) Formula on the Summary sheet} (4 marks):

A reference to another sheet uses the syntax \emph{SheetName!Cell}. The formula that adds the three term totals is:
\[ =\text{Term1!B10}+\text{Term2!B10}+\text{Term3!B10} \]
Equivalently, using a 3-D reference across consecutive sheets: \texttt{=SUM(Term1:Term3!B10)}. (2 marks for correct sheet!cell syntax, 2 marks for adding the three correct cells.)

\textbf{(b) Effect of renaming ``Term2'' to ``SecondTerm''} (4 marks):

In modern spreadsheets (Excel, LibreOffice Calc), renaming a sheet does \emph{not} break the link: the application \textbf{automatically updates} all formulas that reference the sheet, so the formula becomes \texttt{=Term1!B10+SecondTerm!B10+Term3!B10} and the result is unchanged. However, if the sheet is \emph{deleted} (or the reference is typed as plain text in an external file path), the link breaks and the formula returns the error \textbf{\#REF!}. (2 marks for stating the automatic update, 2 marks for mentioning the \#REF! risk in the deletion case.)

\textbf{(c) Method to avoid broken links} (5 marks -- any well-developed method, e.g.):
\begin{itemize}
\item Use \textbf{named ranges} (e.g.\ define the name \emph{TotalTerm2} for \texttt{Term2!\$B\$10}) and write formulas with the names: \texttt{=TotalTerm1+TotalTerm2+TotalTerm3}. Names refer to cells, not to sheet labels typed in formulas, so maintenance is easier and formulas are more readable.
\item Finalise sheet names \emph{before} building the summary formulas, and never delete a source sheet -- hide it instead if necessary.
\item Keep all linked sheets in the \emph{same workbook} rather than linking to external files, and avoid moving/renaming the file or its folder.
\end{itemize}

\textbf{Examiner expectations:} precise syntax with the exclamation mark in (a); in (b) the candidate must distinguish \emph{renaming} (link updated automatically) from \emph{deleting} (\#REF!); in (c) any coherent, technically valid preventive method with a brief explanation earns full marks.
\end{corrige}

\newpage

\coursec{Summary sheet}

\begin{popfigure}
\begin{tikzpicture}[grow cyclic, mindmap, every node/.style={concept, text=white, font=\small},
root concept/.append style={concept color=popdark, font=\small\bfseries},
level 1/.append style={level distance=3.4cm, sibling angle=72},
level 1 concept/.append style={font=\scriptsize\bfseries},
level 2/.append style={level distance=2.2cm, sibling angle=45},
level 2 concept/.append style={font=\tiny, concept color=popmuted}]
\node[root concept]{Word Processing \& Spreadsheets}
 child[concept color=popblue]{ node {Producing Documents}
   child { node {Styles, headers, footers} }
   child { node {Mail merge, TOC} }
 }
 child[concept color=popgreen]{ node {Formatting Data}
   child { node {Sort, filter, validate} }
   child { node {Conditional formatting} }
 }
 child[concept color=poporange]{ node {References \& Formulae}
   child { node {Relative A1} }
   child { node {Absolute \$A\$1} }
   child { node {SUM, AVERAGE, IF} }
 }
 child[concept color=poppurple]{ node {Functions \& Charts}
   child { node {VLOOKUP, COUNTIF} }
   child { node {Column, pie, line charts} }
   child { node {Linking sheets} }
 }
 child[concept color=poppink]{ node {What-If Analysis}
   child { node {Goal Seek} }
   child { node {Scenarios, data tables} }
 };
\end{tikzpicture}
\legende{Mind map of the chapter: from document production to what-if analysis.}
\end{popfigure}

\begin{bilan}
\textbf{Key definitions}
\begin{itemize}
\item A \cle{word processor} (MS Word, LibreOffice Writer) is software used to create, edit, format, save and print text documents.
\item \cle{Formatting} means changing the appearance of a document: character formatting (font, size, bold, colour), paragraph formatting (alignment, line spacing, indentation) and page formatting (margins, orientation, headers and footers).
\item A \cle{style} is a named set of formatting attributes applied in one click; modifying the style updates the whole document automatically.
\item \cle{Mail merge} combines a main document (letter) with a data source (list of names) to produce personalised documents.
\item A \cle{spreadsheet} (MS Excel, LibreOffice Calc) is a grid of \cle{cells} organised in rows (numbered) and columns (lettered); the intersection $B3$ is a cell, a rectangular block like $A1:D10$ is a \cle{range}.
\item A \cle{formula} always begins with \texttt{=} and combines values, cell references and functions; a \cle{function} is a predefined formula such as \texttt{SUM} or \texttt{AVERAGE}.
\item \cle{What-if analysis} is the technique of changing input values to observe the effect on results computed by formulae.
\end{itemize}

\textbf{Formulae and functions to know}
\begin{itemize}
\item Basic: \texttt{=SUM(A1:A10)}, \texttt{=AVERAGE(B2:B20)}, \texttt{=MAX(...)}, \texttt{=MIN(...)}, \texttt{=COUNT(...)} (numbers only), \texttt{=COUNTA(...)} (non-empty cells).
\item Logic: \texttt{=IF(C2>=10,"Admitted","Failed")} --- syntax \texttt{IF(test; value\_if\_true; value\_if\_false)}.
\item Advanced: \texttt{=COUNTIF(A2:A50,"M")}, \texttt{=SUMIF(B2:B50,">10000")}, \texttt{=VLOOKUP(code; table; column; FALSE)}.
\item Referencing: \textbf{relative} \texttt{A1} changes when copied; \textbf{absolute} \texttt{\$A\$1} never changes; \textbf{mixed} \texttt{\$A1} or \texttt{A\$1} locks only the column or the row.
\item Linking worksheets: \texttt{=Sheet2!B5} reads a cell from another sheet; \texttt{='[Book2.xlsx]Sheet1'!A1} reads from another workbook.
\item Percentage of a total: $\text{share} = \dfrac{\text{part}}{\text{total}} \times 100$.
\end{itemize}

\textbf{Methods}
\begin{itemize}
\item Producing a document: type the text $\rightarrow$ save early $\rightarrow$ apply character and paragraph formatting $\rightarrow$ insert page numbers, header/footer $\rightarrow$ use heading styles $\rightarrow$ generate the table of contents $\rightarrow$ proofread and print preview.
\item Building a marksheet: enter labels and data $\rightarrow$ write one formula with correct references $\rightarrow$ copy it down with the fill handle $\rightarrow$ format (borders, number format, conditional formatting) $\rightarrow$ sort or filter if required.
\item Goal Seek: \emph{Data $\rightarrow$ What-If Analysis $\rightarrow$ Goal Seek}: set the result cell, to the target value, by changing the input cell.
\item Charts: select the data (with labels) $\rightarrow$ Insert $\rightarrow$ choose chart type (column for comparisons, pie for shares, line for trends) $\rightarrow$ add title and axis labels.
\end{itemize}

\textbf{Common mistakes}
\begin{itemize}
\item Forgetting the \texttt{=} sign: the entry is stored as text, not computed.
\item Copying a formula with a relative reference where an absolute one was needed (e.g.\ a single VAT cell must be \texttt{\$B\$1}).
\item Confusing \texttt{COUNT} (counts numbers) with \texttt{COUNTA} (counts all non-empty cells).
\item Using a pie chart for time series, or a chart without title and labels.
\item Typing results by hand instead of using formulae: the sheet no longer updates when data change.
\item In mail merge, forgetting to link the data source before inserting merge fields.
\end{itemize}

\textbf{GCE tips}
\begin{itemize}
\item Read the question twice: the exam often asks for the \emph{exact formula} to write in a named cell --- quote it fully, e.g.\ \texttt{=IF(D2>=50,"Pass","Fail")}.
\item State the chart type \emph{and justify} the choice in one sentence.
\item For practical papers, save your work regularly with the exact file name requested.
\item Know the vocabulary: cell, range, worksheet, workbook, fill handle, sorting key, filtering criterion.
\item When asked to distinguish, contrast in two clear columns (e.g.\ relative vs absolute referencing).
\end{itemize}
\end{bilan}

\findecours

\end{document}
